\documentclass{article}

\usepackage[utf8]{inputenc}
\usepackage[T1]{fontenc}
\usepackage{csquotes}
\usepackage[english]{babel}
\usepackage[a4paper, margin=1in]{geometry}
\usepackage[section]{placeins}
\usepackage[backend=biber,maxbibnames=9,sorting=ynt,style=alphabetic]{biblatex}
\usepackage[hidelinks]{hyperref}
\usepackage{appendix}
\usepackage{dutchcal}
\DeclareMathAlphabet{\pazocal}{OMS}{cmsy}{m}{n}
\usepackage{amsmath}
\usepackage[capitalize,noabbrev]{cleveref}
\usepackage{amssymb}
\usepackage{amsthm}
\usepackage{stmaryrd}
\usepackage{mathpartir}
\usepackage{mathtools}
\usepackage{bbm}
\usepackage{xspace}
\usepackage{xstring}
\usepackage{suffix}
\usepackage{enumerate}
\usepackage[theorems,breakable,skins]{tcolorbox}
\usepackage{xcolor}
\usepackage{tikz}
\usetikzlibrary{arrows.meta,calc,positioning,shapes.multipart,cd,decorations.pathreplacing,decorations.pathmorphing}
\usepackage{tikz-cd}
\usepackage{mdframed}
\usepackage{aliascnt} %% for fixing cleveref with theorems
% \usepackage{parskip}  % adds a line between paragraphs

\addbibresource{bibliography.bib}

% Make paragraph bold
\makeatletter
\def\@parfont{\bfseries}
\makeatother

% Shorter section references
\crefformat{section}{\S#2#1#3}

% Removes small-caps font from label in inferrule
\renewcommand{\LabTirName}[1]{{#1}}

% --------
% 1. Misc.
% ========

% Colors
\newcommand{\blue}[1]{\textcolor{blue}{#1}}
\newcommand{\gray}[1]{\textcolor{gray}{#1}}
\newcommand{\red}[1]{\textcolor{red}{#1}}

% Notes
\newcommand{\orpheas}[1]{\nl\red{\textit{Orpheas: }#1}}

% Writing
\def\eg{e.g.,\xspace}
\def\ie{i.e.,\xspace}

% Spacing
\def\spac{\,}
\def\nl{\mbox{}\\}

% Other
\def\dash{-}

\def\tinym{\scriptscriptstyle}
\def\smallm{\scriptstyle}

% minimum value: 7
\newcommand{\ndsize}{7}

\newcommand{\shpscale}{0.04*\ndsize}
\def\gap{0.2*\ndsize}
\def\vshift{\ndsize}
\def\hshift{0.8*\ndsize}

% Tikz Styles for heaplets
\tikzstyle{ptr}=[draw, thin, {Circle[scale=\shpscale]}-{Triangle[scale=\shpscale]}]
\tikzstyle{ccell}=[draw, 
    minimum size=\ndsize pt, outer sep=0pt, inner sep=0pt,
    label={[shift={(-0.8*\ndsize pt,-0.4*\ndsize pt)},color=black,scale=\shpscale]#1}]
\tikzstyle{vcell}=[draw, dashed, 
    minimum size=\ndsize pt, outer sep=0pt, inner sep=0pt,
    label={[shift={(-0.8*\ndsize pt,-0.4*\ndsize pt)},color=black,scale=\shpscale]#1}]

\tikzstyle{brace}=[draw,decorate,decoration={brace,amplitude=5pt,raise=5ex}]
\tikzstyle{labbrace}=[midway,xshift=3.5em,yshift=1pt]

\newcommand{\pbsquare}[1]{
  \IfSubStr{#1}{l}
    {\IfSubStr{#1}{u}{\arrow[#1, phantom, pos=0.15, "{\ulcorner}"]}{\arrow[#1, phantom, pos=0.15, "{\llcorner}"]}}
    {\IfSubStr{#1}{u}{\arrow[#1, phantom, pos=0.15, "{\urcorner}"]}{\arrow[#1, phantom, pos=0.15, "{\lrcorner}"]}}
}
\newcommand{\posquare}[1]{
  \IfSubStr{#1}{u}
    {\arrow[#1, phantom, pos=0.15, "{\ulcorner}"]}
    {\arrow[#1, phantom, pos=0.15, "{\llcorner}"]}
}
\newcommand{\lisquare}[1]{\arrow[#1, phantom, pos=0.15, "{\scriptstyle\vee}"]} 
\tikzset{rmonotip/.tip={Glyph[glyph math command=rto]}}
\tikzset{alloctip/.tip={Glyph[glyph math command=ato]}}
\tikzset{malloctip/.tip={Glyph[glyph math command=mato]}}

\newenvironment{arrow-diagram}[1][normal]{
\begin{tikzcd}[
    cramped,
    nodes={every node/.style={inner sep=0pt, outer sep=0pt}},
    arrows={start anchor=real center, end anchor=real center, shorten >= 1pt, shorten <= 1pt},
    sep={#1}]
}{\end{tikzcd}}

% Tikz arrow movements
\def\north{(0,\vshift pt)}
\def\nnorth{(0,\vshift-\gap pt)}
\def\nnnorth{(0,\vshift-2*\gap pt)}
\def\west{(-\hshift pt,0)}
\def\wwest{(-0.4*\hshift-\hshift pt,0)}
\def\east{(\hshift pt,0)}
\def\eeast{(0.4*\hshift+\hshift pt,0)}
\def\south{(0,-\vshift pt)}
\def\ssouth{(0,-\vshift+\gap pt)}
\def\sssouth{(0,-\vshift+2*\gap pt)}

% ---------------
% 2. Mathematics.
% ===============

\renewcommand\qedsymbol{\hbox{\rlap{$\sqcap$}$\sqcup$}}

% Misc
\def\eqdef{:=}
\def\defeq{=:}
\def\inddef{::=}
\newcommand{\overeq}[1]{\stackrel{#1}{=}}
\def\iso{\cong}
\def\mod{\mathsf{mod}}
\newcommand{\img}[1]{\mathsf{Img}(#1)}
\newcommand{\dom}[1]{\mathsf{Dom}(#1)}
\newcommand{\codom}[1]{\mathsf{Codom}(#1)}
\WithSuffix\newcommand\inv^[1]{#1^{\text{-}1}}
\renewcommand{\vec}{\overline}
\def\from{\leftarrow}
\def\xfrom{\xleftarrow}
\def\xto{\xrightarrow}
\def\injto{\rightarrowtail}
\def\inclto{\hookrightarrow}
\def\xinclto{\xhookrightarrow}
\def\surjto{\twoheadrightarrow}
\newcommand{\xsurjto}[1]{\xtwoheadrightarrow{#1}}
\newcommand{\xrto}[1]{\xrightarrow{#1}_{\mathrlap{\hspace{-1.5pt}{\scriptscriptstyle\mathsf{r}}}}}
\def\rto{\xrto\relax}
\newcommand{\xato}[1]{\xrightarrow{#1}_{\mathrlap{\hspace{-2.7pt}{\scriptscriptstyle\mathrm\nu}}}\,}
\def\ato{\xato\relax}
\def\fnu{\rotatebox[origin=c]{180}{$\scriptscriptstyle\mathrm\nu$}}
\newcommand{\xmato}[1]{\xrightarrow{#1}^{\mathrlap{\hspace{-2.7pt}{\fnu}}}\,}
\def\mato{\xmato\relax}
\def\To{\Rightarrow}
\def\From{\Leftarrow}
\def\dashto{\dashrightarrow}
\makeatletter
\newcommand{\xdashrightarrow}[2][]{\ext@arrow 0359\rightarrowfill@@{#1}{#2}}
\newcommand{\xdashleftarrow}[2][]{\ext@arrow 3095\leftarrowfill@@{#1}{#2}}
\newcommand{\xdashleftrightarrow}[2][]{\ext@arrow 3359\leftrightarrowfill@@{#1}{#2}}
\def\rightarrowfill@@{\arrowfill@@\relax\relbar\rightarrow}
\def\leftarrowfill@@{\arrowfill@@\leftarrow\relbar\relax}
\def\leftrightarrowfill@@{\arrowfill@@\leftarrow\relbar\rightarrow}
\def\arrowfill@@#1#2#3#4{%
  $\m@th\thickmuskip0mu\medmuskip\thickmuskip\thinmuskip\thickmuskip
   \relax#4#1
   \xleaders\hbox{$#4#2$}\hfill
   #3$%
}
\makeatother
\def\xdashto{\xdashrightarrow}
\newcommand{\xsquigrightarrow}[1]{\overset{#1}{\rightsquigarrow}}
\def\xsquigto{\xsquigrightarrow}
\def\xiff{\xLeftrightarrow}

\def\mwhere{\mathit{where}\;}
\def\mif{\mathit{if}\;}
\def\suchthat{\textit{s.t.}}

% Term list.
\newcommand{\tlist}[2]{{#1}_1,...,{#1}_{#2}}
\newcommand{\tlistr}[2]{(#1)_{#2}}
\newcommand{\tsetr}[2]{\{#1\}_{#2}}

% Common metavariables.
% -- Vectors of binders.
\def\vxs{\vec\x}
\def\vys{\vec\y}
% -- Variables.
\def\va{\mathit{a}}
\def\vb{\mathit{b}}
\def\vc{\mathit{c}}
\def\vd{\mathit{d}}
\def\ve{\mathit{e}}
\def\vf{\mathit{f}}
\def\vg{\mathit{g}}
\def\vh{\mathit{h}}
\def\vi{\mathit{i}}
\def\vj{\mathit{j}}
\def\vk{\mathit{k}}
\def\vl{\mathit{l}}
\def\vm{\mathit{m}}
\def\vn{\mathit{n}}
\def\vo{\mathit{o}}
\def\vp{\mathit{p}}
\def\vq{\mathit{q}}
\def\vr{\mathit{r}}
\def\vs{\mathit{s}}
\def\vt{\mathit{t}}
\def\vu{\mathit{u}}
\def\vv{\mathit{v}}
\def\vw{\mathit{w}}
\def\vx{\mathit{x}}
\def\vy{\mathit{y}}
\def\vz{\mathit{z}}
\def\vA{\mathit{A}}
\def\vB{\mathit{B}}
\def\vC{\mathit{C}}
\def\vD{\mathit{D}}
\def\vE{\mathit{E}}
\def\vF{\mathit{F}}
\def\vG{\mathit{G}}
\def\vH{\mathit{H}}
\def\vI{\mathit{I}}
\def\vJ{\mathit{J}}
\def\vK{\mathit{K}}
\def\vL{\mathit{L}}
\def\vM{\mathit{M}}
\def\vN{\mathit{N}}
\def\vO{\mathit{O}}
\def\vP{\mathit{P}}
\def\vQ{\mathit{Q}}
\def\vR{\mathit{R}}
\def\vS{\mathit{S}}
\def\vT{\mathit{T}}
\def\vU{\mathit{U}}
\def\vV{\mathit{V}}
\def\vW{\mathit{W}}
\def\vX{\mathit{X}}
\def\vY{\mathit{Y}}
\def\vZ{\mathit{Z}}
\def\var{\mathit{x}}
\def\Var{\mathit{X}}
\def\val{\mathit{V}}
\def\wal{\mathit{W}}
\def\vop{\mathit{op}}
\def\compM{\mathit{M}}
\def\compN{\mathit{N}}
\def\pvarA{\mathit{a}}
\def\pvarB{\mathit{b}}
\def\pvarC{\mathit{c}}
\def\pvarD{\mathit{d}}
% Greek meta-variables (μεταβλητές)
\def\me{\varepsilon}
\def\t{\tilde}
\def\lt{\widetilde}
\def\ms{\sigma}
\def\mt{\tau}
\def\mm{\mu}
\def\mz{\zeta}

% Nat.
\def\nat{\mathbb{N}}
\def\Succ{\mathit{S}\,}

% Maps.
\def\mvarM{\eta}
\def\mapempty{\varnothing}
\def\mapbinding{\mapsto}
\newcommand{\maplookup}[2]{#1(#2)}
\newcommand{\mapupdate}[2]{[{#1} \mapbinding {#2}]}

% Lists.
\def\listnil{[]}
\def\listcons{;}
\newcommand{\listsing}[1]{[#1]}
\newcommand{\listtwo}[2]{[#1; #2]}
\def\listapp{+\hspace{-4pt}+\spac}
\newcommand{\listlength}[1]{|#1|}

% Sets
\newcommand{\singset}[1]{\{ #1 \}}
\newcommand{\powerset}[1]{\mathcal{P}(#1)}
\def\setDiff{\setminus}
\def\setInt{\cap}
\def\setUnion{\cup}
\def\setDUnion{\uplus}
\def\fin{\mathsf{fin}}
\def\subseteqfin{\subseteq_\fin}
\newcommand{\card}[1]{|#1|}
\newcommand{\downset}[1]{[#1]}
\newcommand{\klstar}[1]{{#1}^*}
\newcommand\quotient[2]{{^{#1}}\hspace{-3pt}\Big/{_{#2}}}
\def\quot{q}
\newcommand\eqclass[1]{q_{#1}}
\newcommand\meduplus[1]{\scalebox{1}{$\displaystyle\biguplus_{#1}$}}
\newcommand\disunion[1]{\displaystyle\meduplus{#1}}
\newcommand{\const}[1]{\mathsf{K}(#1)}

% -------------------------------------
% 3. Programming Language (CBPV-Style).
% =====================================

\newcommand{\exprkw}[1]{{\rm {\bf #1}}}
\newcommand{\weaken}[2]{\prescript{#1}{}{\hspace{-1pt}#2}}

% Identifiers.
\newcommand{\ident}[1]{\ensuremath{\mathit{#1}}\xspace}
% Literals.
\def\True{\exprkw{T}}
\def\False{\exprkw{F}}
% Abstraction.
\newcommand{\Lam}[1]{\lambda#1.\spac}
\newcommand{\LamT}[2]{\lambda(#1 \,: #2).\spac}
\newcommand{\LamTwo}[2]{\lambda#1\spac#2.\spac}
\newcommand{\LamTwoT}[4]{\lambda(#1 \spac: #2)\spac(#3 \spac: #4).\spac}
\newcommand{\LamThree}[3]{\lambda#1\spac#2\spac#3.\spac}
\newcommand{\LamFour}[4]{\lambda#1\spac#2\spac#3\spac#4.\spac}
% Application.
\def\App{\spac}
\newcommand{\BApp}[1]{(#1)}
% Return
\def\Ret{\exprkw{return}\spac}
% Do
\def\revto{\leftarrow}
\newcommand{\Do}[2]{\exprkw{do}\spac\mathit{#1}\revto{#2}\spac\exprkw{in}\spac}
\newcommand{\DoV}[3]{
\begin{array}[t]{@{}l}
    \exprkw{do}\spac\mathit{#1}\revto{#2}                                            \\
    \exprkw{in}\spac{#3}
\end{array}}
% Thunks
\def\Thunk{\exprkw{thunk}\spac}
\def\Force{\exprkw{force}\spac}
% Operation Call
\newcommand{\OCall}[3]{{#1}\BApp{{#2};{#3}}}
% If-Then
\newcommand{\Ite}[2]{\exprkw{if}\spac{#1}\spac\exprkw{then}\spac{#2}\spac\exprkw{else}\spac}
% Pairs
\def\Unit{\langle\rangle}
\newcommand{\VPair}[2]{( #1 , #2 )}
\newcommand{\Pair}[2]{\langle #1 , #2 \rangle}
\newcommand{\Triple}[3]{\langle #1 , #2, #3 \rangle}
\def\Fst{\exprkw{fst}\spac}
\def\Snd{\exprkw{snd}\spac}
\def\Thd{\exprkw{thd}\spac}
\newcommand{\Split}[3]{\exprkw{split}\spac{#1}\spac\exprkw{as}\spac\Pair{#2}{#3}.\spac}
\newcommand{\SplitV}[4]{
\begin{array}[t]{@{}l}
\exprkw{split}\spac{#1}\spac\{ \\
    \qquad \Pair{#2}{#3}. \spac #4 \\
\}
\end{array}}
% Variants
\def\InjL{\exprkw{inl}\spac}
\def\InjR{\exprkw{inr}\spac}
\newcommand{\Case}[5]{\exprkw{case}\spac{#1}\spac\{ #2.\spac #3 ; #4. \space #5 \}}
\newcommand{\CaseV}[5]{
\begin{array}[t]{@{}l}
    \exprkw{case}\spac{#1}\spac \\
    \quad#2. \spac #3 \\
    \quad#4. \spac #5
\end{array}}
\newcommand{\Absurd}[1]{\exprkw{case}\spac{#1}\{\}}

% ----------------------
% 4. Types (CBPV style).
% ======================

% Keywords.
\newcommand{\typekw}[1]{\textsf{#1}}

% Type variables
\def\tvarA{\alpha}
\def\tvarB{\beta}
\def\tvarC{\gamma}
\def\tvarsA{\vec\alpha}

% Value type metavariables
\def\vtmvarA{A}
\def\vtmvarB{B}
\def\vtmvarC{C}
\def\vtmvarD{D}

% Computation type metavariables
\newcommand{\ctype}[1]{\underline{\textnormal{#1}}}
\def\ctmvarA{\ctype\vtmvarA}
\def\ctmvarB{\ctype\vtmvarB}
\def\ctmvarC{\ctype\vtmvarC}
\def\ctmvarD{\ctype\vtmvarD}

% Value types
\def\tyZero{\typekw{0}}
\def\tyUnit{\typekw{1}}
\def\tyBool{\typekw{B}}
\def\tyNat{\typekw{N}}
\newcommand{\tyRef}[1]{\typekw{ref}_{#1}}
\def\tyU{\typekw{U}\spac}
\def\tySum{\spac\typekw{$+$}\spac}
\def\tyProd{\spac\typekw{$\times$}\spac}

% Computation types
\def\tyF{\typekw{F}}
\def\tyFun{\rightarrow}
\def\tyCProd{\spac\typekw{$\&$}\spac}

% --------------------
% 5. Typing Judgments.
% ====================

\def\parctx{\Delta}
\def\parctxB{\parctx'}
\def\valctx{\Gamma}
\def\valctxB{\valctx'}
\def\ctxPlus{,}
\def\ctxNil{\cdot}

\newcommand{\FVars}[1]{\mathsf{FreeVar}\spac#1}

\newcommand{\ctyped}[4]{#1 \mid #2 \vdash_{\mathsf{c}} #3 : #4}
\newcommand{\vtyped}[4]{#1 \mid #2 \vdash_{\mathsf{v}} #3 : #4}

\newcommand{\atjudg}[3]{#1 \vdash #2 : #3}
\newcommand{\ateqjudg}[3]{#1 \vdash #2 \equiv #3}
\newcommand{\patjudg}[3]{#1 \mid #2 \vdash #3}
\newcommand{\pateqjudg}[4]{#1 \mid #2 \vdash & #3 \equiv #4}

\newcommand{\RULEDEF}[2]{\hypertarget{#1}{{(#2)}}\xspace}
\newcommand{\RULE}[2]{\hyperlink{#1}{(#2)}\xspace}

\def\nameTVar{T-Var}
\def\nameTThunk{T-Thunk}
\def\nameTInjL{T-Injl}
\def\nameTInjR{T-Injr}
\def\nameTUnit{T-Unit}
\def\nameTApp{T-App}
\def\nameTLam{T-Lam}
\def\nameTForce{T-Force}
\def\nameTReturn{T-Ret}
\def\nameTDo{T-Do}
\def\nameTPair{T-Pair}
\def\nameTFst{T-Fst}
\def\nameTSnd{T-Snd}
\def\nameTSplit{T-Split}
\def\nameTAbsurd{T-Absurd}
\def\nameTCase{T-Case}

% -----------------------------------
% 6. Parametrised Algebraic Theories.
% ===================================

\def\vMod{\mathfrak{M}}
\def\vParam{\mathfrak{p}}

% Local State specific
\def\eqOp{\ensuremath{?_=}}
\newcommand{\Eq}[2]{\,?_{#1 = #2}\,}
\def\eqOpW{\mathit{eq}}
\newcommand{\EqW}[4]{\eqOpW(#1, #2, #3, #4)}
\def\newOp{\ensuremath{\nu}}
\newcommand{\New}[2]{\newOp_{#1}{#2}.\spac}
\def\readOp{\mathit{r}}
\newcommand{\Read}[3]{\readOp(#1, #2, #3)}
\def\writeOp{\mathit{w}}
\newcommand{\Write}[3]{\writeOp_{#1}(#2, #3)}

\newcommand{\MatchSing}[2]{\exprkw{match}\spac #1 \To\spac #2}
\newcommand{\MatchDual}[4]{\exprkw{match}\spac #1 \To\spac #2 ; #3 \To\spac #4}
\newcommand{\MatchDualV}[4]{
\begin{array}[t]{@{}l}
    \exprkw{match} \\
    \quad #1 \To \spac #2; \\
    \quad #3 \To \spac #4
\end{array}}
\newcommand{\MatchVec}[2]{\exprkw{match}\spac \{ #2 \}_{#1}}

% -------------------
% 7. Category Theory.
% ===================

\def\funmvarT{T}
\def\funmvarF{F}
\def\funmvarG{G}
\def\mcatI{\pazocal{I}}
\def\mcatC{\pazocal{C}}
\def\mcatD{\pazocal{D}}
\def\mcatJ{\pazocal{J}}

\def\comp{\circ}

\newcommand{\objC}[1]{\mathsf{Obj}({#1})}
\def\setC{\mathsf{Set}}
\def\catC{\mathsf{Cat}}
\newcommand{\funC}[2]{[#1, #2]}
\newcommand{\pshC}[1]{\mathsf{Psh}({#1})}
\newcommand{\famC}[1]{\mathsf{Fam}({#1})}
\newcommand{\fsreprC}[1]{\pshC{#1}_\sqcup}
\newcommand{\opC}[1]{{#1}^{\mathsf{op}}}
\newcommand{\algC}[1]{{#1}\mathsf{-Alg}}
\newcommand{\commaC}[2]{{#1}\Downarrow{#2}}
\newcommand{\sliceC}[2]{{#1}/{#2}}
\newcommand{\arrC}[1]{\mathsf{Arr}({#1})}
\newcommand{\elemC}{\oint\,}
\def\equivC{\simeq}

\def\finC{\mathsf{Fin}}
\def\injC{\mathsf{Inj}}
\def\worldC{\mathbb{W}}
\def\inst{\mathbb{I}}
\def\instC{\worldC\inst}
\def\tempC{\mathbb{T}}
\def\tinstC{{\tempC\inst}}
\newcommand\allocC[1]{\mathsf{alloc}(#1)}

\newcommand{\homF}[3]{{#1}({#2}, {#3})}
\newcommand{\diagF}[1]{\Delta_{#1}}
\newcommand{\Ran}[2]{\mathrm{Ran}_{#1}\,{#2}}
\newcommand{\Colimit}[1]{\mathit{colimit}\,{#1}}
\newcommand{\Prod}[1]{\prod_{#1}\spac}
\newcommand{\Coprod}[1]{\coprod_{#1}\spac}
\newcommand{\Exp}[2]{{#1}^{#2}}
\newcommand{\End}[1]{\displaystyle{\int_{#1}}\spac}
\newcommand{\Coend}[1]{\displaystyle{\int^{\ensuremath{#1}}}\spac}
\newcommand{\pltimes}[1]{\otimes_{#1}}
\newcommand{\mplus}{\oplus}
\newcommand{\liplus}[1]{\oplus_{#1}}
\newcommand{\uwedge}[1]{\owedge_{#1}}
\newcommand{\eqz}[2]{\mathsf{eqz}(#1, #2)}
\newcommand{\coeqz}[2]{\mathsf{coeqz}(#1, #2)}

\def\yoneda{\vy}
\newcommand\indF[1]{|{#1}|}
\def\monFLS{\mathcal{T}}
\def\monFLSFin{\monFLS^\fin}
\def\heapF{\mathbbm{H}}
\def\dheapF{\underline\heapF}
\def\theapF{{\tempC\heapF}}
\newcommand\resF[1]{\mathsf{R}^{#1}}
\def\hideMon{{\mathbb{P}}}
\def\thideMon{{\tempC\hideMon}}
\def\hmonIso{\,{\overset{\mbox{\tiny $\hideMon$}}{\iso}}\,}
\newcommand\stencilF[1]{\mathcal{S}_{#1}}
\newcommand\stencilRebase[1]{\mathcal{S}^{#1}}
\newcommand{\tempCompl}[1]{{#1}^*}
% \newcommand{\inclITI}[1]{\IfStrEq{#1}{\dash}{1_{(\dash)}}{1_{#1}}}
\newcommand{\inclITI}[1]{1_{#1}}
\newcommand{\extF}[2]{(#1 \oplus #2)}
\newcommand{\Plus}[1]{{#1}^+}
\newcommand{\UPlus}[1]{(#1)^{\hspace{-1pt}+}\hspace{-1pt}}
\newcommand{\Minus}[1]{{#1}^-}
\newcommand{\UMinus}[1]{(#1)^{\hspace{-1pt}-}\hspace{-1pt}}
\def\instU{\vu}
\def\tinstU{\vu}
\newcommand{\concL}[1]{\lfloor#1\rfloor}
\newcommand{\varL}[1]{\lceil#1\rceil}
\newcommand{\anyL}[1]{[#1]}
\def\isIn{\;\mathsf{is}\,\mathsf{in}\;}
\newcommand{\uncon}[1]{\mathsf{uc}(#1)}
\newcommand{\con}[1]{\mathsf{c}(#1)}
\newcommand{\dpres}{\mathsf{dpres}}
\newcommand{\ureach}{\mathsf{ur}}
\newcommand{\reach}{\mathsf{r}}
\newcommand{\GU}[2]{\mathsf{GU}(#1, #2)}
\newcommand\realise[3]{#1,#2 \Vdash #3}

\def\unit{\eta}
\def\counit{\varepsilon}
\def\bind{\ensuremath{\gg\!\!=}}
\def\mul{\mu}
\WithSuffix\newcommand\ext^[1]{#1^*}
\WithSuffix\newcommand\res_[1]{{#1}_*}
\WithSuffix\newcommand\res|[1]{|^{#1}}
\WithSuffix\newcommand\uext^[1]{#1\rangle}
\WithSuffix\newcommand\proj_[1]{\pi_{#1}}
\newcommand{\pair}[2]{\langle#1, #2\rangle}
\newcommand{\ipair}[2]{\langle \hdots, #1 \mapsto #2, \hdots \rangle}
\newcommand{\triple}[3]{\langle#1, #2, #3\rangle}
\newcommand{\copair}[2]{[#1, #2]}
\newcommand{\cotriple}[3]{[#1, #2, #3]}
\def\projOne{\proj_{1}}
\def\projTwo{\proj_{2}}
\def\projThree{\proj_{3}}
\def\projFour{\proj_{4}}
\def\projFive{\proj_{5}}
\WithSuffix\newcommand\inj_[1]{\mathit{inj}_{#1}}
\WithSuffix\newcommand\inj^[1]{\mathit{inj}^{#1}}
\def\inl{\mathit{inl}}
\WithSuffix\newcommand\inl_[1]{\mathit{inl}^{\tinym#1}}
\WithSuffix\newcommand\incl_[1]{\mathit{i}^{\tinym#1}}
\def\inr{\mathit{inr}}
\newcommand{\curry}[1]{\lambda_{#1}}
\def\incl{\mathit{i}}
\def\eval{\mathsf{eval}}

\def\cov{\mathsf{cov}}
\def\rebase{\mathsf{rebase}}
\def\embed{\mathsf{embed}}
\def\ev{\mathsf{ev}}
\def\hideSt{\mathsf{hide}}
\newcommand{\compl}[1]{{#1}^\mathsf{c}}
\newcommand{\extNat}[2]{\mathsf{ext}_{#1}^{#2}}
\newcommand{\rebaseF}[1]{\elemC{\stencilRebase{#1}}}
\newcommand{\reprFinF}[2]{\homF\finC{#1}{#2}}
\newcommand{\hideHMon}[2]{\mathsf{hide}^{#1}_{#2}}
\newcommand{\hideRes}[2]{\mathsf{hide}^{#1}_{#2}}

\def\eto{\surjto}
\newcommand{\xtwoheadrightarrow}[2][]{%
  \xrightarrow[#1]{#2}\mathrel{\mkern-14mu}\rightarrow
}
\newcommand{\xeto}[1]{\xtwoheadrightarrow{#1}}
\def\rmto{\to_{\hspace{-1pt}\mathsf{r}}}
\newcommand{\xrmto}[1]{\xrightarrow{#1}_{\hspace{-1pt}\mathsf{r}}}
\newcommand{\aiso}[1]{\iso_\nu^{#1}}

% -------------
% 7. Semantics.
% =============

\def\typeLoc{\mathbb{L}}
\def\typeVal{\mathsf{Val}}
\def\typeSort{\mathsf{Sort}}
\def\typeFG{\mathsf{FGType}}
\def\typeWorld{\mathsf{World}}
\def\typeEctx{\mathsf{ECtx}}
\def\typeRes{\mathsf{Res}}

\def\ectx{\mathit{E}}
\def\hole{[\,]}

\newcommand{\pointers}[2]{\mathsf{locs}^{#1}(#2)}

\newcommand{\Subst}[2]{[#1 / #2]}
\newcommand{\SubstTwo}[4]{[#1 / #2, #3 / #4]}

\def\stepsto{\rightarrow}
\newcommand{\lstepsto}[4]{\xrightarrow{ #1(#2)_{#3} \rightsquigarrow {#4}}}
\newcommand{\lterminates}[2]{\downarrow_{#1(#2)}}
\def\Stepsto{\Downarrow}

\newcommand{\sem}[1]{[\![#1]\!]}
\newcommand{\semM}[2]{\sem{#1}_{#2}}

\newcommand{\inferRuleSem}[2]{\left\llbracket\begin{array}{c}\inferrule[]{#1}{#2}\end{array}\right\rrbracket}

% ------------------------
% 8. Logic
% ========================

\newcommand{\All}[1]{\forall #1.\spac}
\newcommand{\AllTwo}[2]{\forall #1,#2.\spac}
\newcommand{\AllThree}[3]{\forall #1,#2,#3.\spac}
\newcommand{\existsU}{\exists!}
\newcommand{\ExistsU}[1]{\exists! #1.\spac}
\newcommand{\Exists}[1]{\exists #1.\spac}
\newcommand{\ExistsTwo}[2]{\exists #1, #2.\spac}
\newcommand{\ExistsThree}[3]{\exists #1, #2, #3.\spac}
\def\Land{\wedge}
\def\Lor{\vee}


\newif\iftodo
% \todotrue
\todofalse

\theoremstyle{definition}
\newtheorem{theorem}{Theorem}

\theoremstyle{definition}
\newtheorem{definition}{Definition}[subsection]
\newaliascnt{lemma}{definition}
\newaliascnt{remark}{definition}
\newaliascnt{example}{definition}
\newaliascnt{corollary}{definition}

\theoremstyle{definition}
\newtheorem{lemma}[lemma]{Lemma}

\theoremstyle{remark}
\newtheorem{remark}[remark]{Remark}

\theoremstyle{definition}
\newtheorem{corollary}[corollary]{Corollary}

\theoremstyle{remark}
\newtheorem{example}[example]{Example}

\newtheoremstyle{proof}
    {3pt}       % Space above
    {3pt}       % Space below
    {}          % Body font
    {}          % Indent amount
    {\itshape}  % Theorem head font
    {.}         % Punctuation after theorem head
    {.5em}      % Space after theorem head
    {}          % Theorem head spec (can be left empty, meaning ‘normal’)
\theoremstyle{proof}

\newtheoremstyle{notation}
    {3pt}
    {3pt}
    {}
    {}
    {\itshape}
    {.}
    {.5em}
    {}
\newtheorem*{notation}{Notation}

\aliascntresetthe{lemma}
\aliascntresetthe{remark}
\aliascntresetthe{corollary}
\aliascntresetthe{example}

\tcolorboxenvironment{definition}{
    colframe=green!75!black,
    colback=green!5,
    boxrule=0.5pt,
    before skip=10pt,after skip=10pt,
    left=4pt, right=4pt
}
\tcolorboxenvironment{lemma}{
    colframe=red!75!black,
    colback=red!5,
    boxrule=0.5pt,
    interior hidden,
    before skip=10pt,after skip=10pt, 
    left=4pt, right=4pt
}
\tcolorboxenvironment{theorem}{
    colframe=red!75!black,
    colback=red!5,
    boxrule=0.5pt,
    interior hidden,
    before skip=10pt,after skip=10pt, 
    left=4pt, right=4pt
}
\tcolorboxenvironment{corollary}{
    colframe=red!75!black,
    colback=red!5,
    boxrule=0.5pt,
    interior hidden,
    before skip=10pt,after skip=10pt, 
    left=4pt, right=4pt
}
\tcolorboxenvironment{proof}{
    colframe=gray!10,
    colback=gray!10,
    interior hidden,
    breakable,
    before skip=10pt,after skip=10pt, left skip=0pt, right skip=0pt,
    left=4pt, right=4pt
}
\tcolorboxenvironment{example}{
    colframe=blue!75!black,
    colback=blue!5,
    boxrule=1pt,
    interior hidden,
    breakable,before skip=10pt,after skip=10pt, 
    left=4pt, right=4pt
}
\tcolorboxenvironment{remark}{
    colframe=yellow!75!black,
    colback=yellow!5,
    boxrule=0.5pt,
    interior hidden,
    breakable,before skip=10pt,after skip=10pt, 
    left=4pt, right=4pt
}
\tcolorboxenvironment{notation}{
    colframe=orange!75!black,
    colback=orange!5,
    boxrule=0.5pt,
    interior hidden,
    breakable,before skip=10pt,after skip=10pt, 
    left=4pt, right=4pt
}

\begin{document}

\title{Template Heaps for Full-Ground Local State}
\author{Orpheas van Rooij}

\maketitle


\section{A theory of template heaps}

In this section we introduce \emph{template heaps}, a notion that describes multiple heaps with a common segment.
We develop a small theory around template heaps in order to achieve our main objective: 
to give a finitary submonad of the Full-Ground Local State monad \cite{kammar2017monad}.

We begin by revisiting the categories $\worldC$ (worlds and injections) and $\instC$ (world and injections with instantiation data) where we mainly focus on their local independence structure~\cite{simpson2018category}.
We introduce the category of templates $\tempC$ that introduces a notion of location variables.
We extend the local independence structure to $\tempC$ in order to define a suitable generalisation of the hiding monad on template heaps.
We give the category of template instantiations $\tinstC$, which introduces instantiation data to $\tempC$.


\subsection[Category of Worlds and Injections]{The category of worlds and injections $\worldC$}

Given a set of sorts $\vC \in \typeSort$, the category $\worldC$ is defined as
\begin{itemize}
    \item Objects: worlds $\vw = \{ \vl_1 : \vC_1, \dots, \vl_\vn : \vC_\vn \}$.
    \item Morphisms: Injective, sort-preserving $\sigma : \vw \to \vw'$.
\end{itemize}

\begin{notation}
We use $\img\sigma$ to refer to its image and $\vw' \ominus \sigma$ for its complement, that is, $\vw' \setminus \img\sigma$, 
which induces an inclusion $\compl\sigma : \vw' \ominus \sigma \to \vw'$.

We denote morphisms from the empty world as $?_\vw : \emptyset \to \vw$ and inclusions as $\incl_\vi^{\vw_1,\vw_2} : \vw_\vi \to \vw_1 \mplus \vw_2$, 
where $\oplus : \worldC \times \worldC \to \worldC$ is the independent coproduct functor (see \cite{kammar2017monad} for definition).
We also denote by $\iso^{\sigma : \vw \to \vw'} : \vw' \iso \vw \oplus (\vw' \ominus \sigma)$ the canonical isomorphism.
\end{notation}

\begin{lemma}[Monomorphisms in $\worldC$]\label{lem:monos-W}
    Every morphism in $\worldC$ is mono.
\end{lemma}
\begin{proof}Immediate from their injectivity.\end{proof}

\begin{lemma}[Initial Object in $\worldC$]
    The category $\worldC$ has an initial object.
\end{lemma}
\begin{proof}Consider the empty world.\end{proof}

\begin{lemma}[Equalizers in $\worldC$]
    The category $\worldC$ has all equalizers.
\end{lemma}
\begin{proof}
    Define the equalizer in 
    \begin{tikzcd}{\eqz{\sigma_1}{\sigma_2}} \ar[r, hook, "\ve"] & {\vw_1} \ar[r, yshift=3pt, "\sigma_1"] \ar[r, yshift=-3pt, "\sigma_2"'] & {\vw_2}\end{tikzcd} 
    to be
    \[ \eqz{\sigma_1}{\sigma_2} \eqdef \{ \va : \vC \in \vw_1 \mid \sigma_1(\va) = \sigma_2(\va) \} \]
    with $\ve$ being the inclusion into $\vw_1$.

    Given a $\sigma_1 \comp \hat\sigma = \sigma_2 \comp \hat\sigma$, let $\vh : \hat\vw \to \eqz{\sigma_1}{\sigma_2}$ be given by $\hat\sigma$.  
    It is clear that it is well-defined, and uniqueness follows from $\ve$ being mono.
\end{proof}
\begin{example}
    This is an equalizer diagram 
    \begin{center}
    \begin{tikzcd}
        {\vb : \vC'} \ar[r, "\incl_2"] 
            & {\va : \vC \mplus \vb : \vC'} 
                \ar[r, yshift=3pt, "\incl_1 \oplus 1"] \ar[r, yshift=-3pt, "\incl_2 \oplus 1"'] 
            & {\va_1 : \vC \oplus \va_2 : \vC \oplus \vb : \vC'}
    \end{tikzcd}
    \end{center}
\end{example}
\begin{notation}
    We sometimes overload $\eqz{\sigma_1}{\sigma_2}$ to refer both to the object and the morphism when it is clear from context, \eg
    \[ \begin{tikzcd}{\eqz{\sigma_1}{\sigma_2}} \ar[rr, "\eqz{\sigma_1}{\sigma_2}"] && {\vw_1} \ar[r, yshift=3pt, "\sigma_1"] \ar[r, yshift=-3pt, "\sigma_2"'] & {\vw_2}\end{tikzcd}\]
\end{notation}

\begin{lemma}[Pullbacks in $\worldC$]
    The category $\worldC$ has all pullbacks.
\end{lemma}
\begin{proof}
    Define the pullback of injections $\vw_1 \xto{\sigma_1} \hat\vw \xfrom{\sigma_2} \vw_2$ to be
    \begin{align*} 
        \vw_1 \pltimes{\hat\vw} \vw_2 & \eqdef \img{\sigma_1} \setInt \img{\sigma_2} \\
        \res_{\sigma_1}\sigma_2 & \eqdef \inv^{\sigma_1} \\
        \res_{\sigma_2}\sigma_1 & \eqdef \inv^{\sigma_2}
    \end{align*}
    It is clear that the square 
        \begin{tikzcd}{\vw_1 \pltimes{\hat\vw} \vw_2} \ar[d, "\res_{\sigma_1}\sigma_2"'] \ar[r, pos=0.45, "\res_{\sigma_2}\sigma_1"] & {\vw_2} \ar[d, "\sigma_2"] \\
            {\vw_1} \ar[r, "\sigma_1"'] & {\hat\vw}
        \end{tikzcd}
        commutes.
    Further assume $\vw' \xto{\sigma'_1} \vw_1$, $\vw' \xto{\sigma'_2} \vw_2$.
    Define the induced morphism $\vh : \vw' \to \vw_1 \pltimes{\hat\vw} \vw_2$ as $\sigma_1 \comp \sigma'_1$. 
    It is well-defined by $\sigma_1 \comp \sigma'_1 = \sigma_2 \comp \sigma'_2$.
    Uniqueness follows from injectivity of $\inv^{\sigma_1}$, $\inv^{\sigma_2}$.
\end{proof}

\begin{example}
    This is a pullback diagram 
    \begin{tikzcd} 
        {\vb : \vC'}
            \ar[d, hook, "\res_\vi\vj"]  \ar[r, hook, "\res_\vj\vi"] \pbsquare{dr}
            & {\vb : \vC', \vc : \vC''}  
                \ar[d, hook, "\vj"] 
        \\
        {\va : \vC, \vb : \vC'} 
            \ar[r, hook, "\vi"] 
            & {\va : \vC, \vb : \vC', \vc : \vC''}
    \end{tikzcd}
\end{example}

\subsubsection{Local Independence Structure}

Pushouts in $\worldC$ do not exist in general as witnessed by $\worldC \inclto \commaC\finC\typeSort$ not reflecting pushouts.
The reason for this stems from the inherent injectivity of morphisms in $\worldC$
certain induced morphisms coming out of the pushout in $\commaC\finC\typeSort$ are not injective.
We instead consider a local independence structure on $\worldC$ which recovers the initia(L.I.C.) lity of the desired construct.
Instead of considering any square as a solution to the problem of closing two morphisms $\vw_1 \xfrom{\sigma_1} \vw \xto{\sigma_2} \vw_2$,
we consider independent squares that only allow locations from $\vw$ to be conflated.
The pushout in $\commaC\finC\typeSort$ becomes their local independent coproduct which takes the disjoint union of their coimages.
We recall here the local independent structure on $\worldC$~\cite{simpson2018category}.

\begin{definition}[Local Independent Squares in $\worldC$]
    \begin{tikzcd}[sep=small] 
                    & {\vw} \arrow[ld, "\sigma_1"'] \arrow[rd, "\sigma_2"] & \\
         {\vw_1} \arrow[rd, "{\sigma_1'}"'] &            & {\vw_2} \arrow[ld, "{\sigma_2'}"] \\
                                            &  {\hat\vw} &
    \end{tikzcd} 
    is a local independent square (L.I.S.) when it is a pullback
    \begin{center}
    \begin{tikzcd}[] 
        {\vw} \arrow[d, "\sigma_1"'] \arrow[r, "\sigma_2"] \pbsquare{dr} & {\vw_2} \arrow[d, "{\sigma_1'}"] \\
        {\vw_1} \arrow[r, "{\sigma_1'}"']                                & {\hat\vw} 
    \end{tikzcd} 
    \end{center}
\end{definition}

\begin{lemma}[Local Independence Structure on $\worldC$]\label{lem:local-ind-struct-W}
We can equip $\worldC$ with a local independent structure~\cite{simpson2018category} by defining the collection of local independent multicospans 
$ \mcatI = \{ \sigma_\vi \xto{\sigma_\vi'} \hat\sigma \}_{\vi \in \vI} $ as multicospans that satisfy (note that $\vI$ is finite): 
\[ \All{\vi, \vj \in \vI} \vi \neq \vj \implies 
    \begin{tikzcd}[sep=small] 
                    & {\vw} \arrow[ld, "\sigma_\vi"'] \arrow[dd, "{\hat\sigma}"] \arrow[rd, "\sigma_\vj"] & \\
        {\vw_\vi} \arrow[rd, "{\sigma_\vi'}"'] &            & {\vw_\vj} \arrow[ld, "{\sigma_\vj'}"] \\
                                          &  {\hat\vw} &
    \end{tikzcd}\text{ is a local independent square.} \]
\end{lemma}
\begin{proof}\nl
    We only need to show that the collection $\mcatI$ is closed under composition, all other axioms follow immediately. 
    Consider multicospan $\{ \sigma_\vi \xto{\sigma'_\vi} \hat\sigma \}_{\vi \in \vI} \in \mcatI$ 
    and a family of multicospans $\{ \{ \tau_{\vi\vj} \xto{\tau'_{\vi\vj}} \sigma_\vi \}_{\vj \in \vJ_\vi} \in \mcatI \}_{\vi \in \vI}$.
    We need $\{ \tau_{\vi\vj} \xto{\sigma'_\vi \comp \tau'_{\vi\vj}} \hat\sigma \}_{\vj \in \vJ_\vi, \vi \in \vI } \in \mcatI$, which follows from
    \begin{enumerate}
        \item For all $\vi \in \vI, \vk \neq \vl \in \vJ_\vi$, $\tau_{\vi\vk} \xto{\sigma'_\vi \comp \tau'_{\vi\vk}} \hat\sigma \xfrom{\sigma'_\vi \comp \tau'_{\vi\vl}} \tau_{\vi\vl}$ is a local independent square.
                \begin{arrow-diagram}
                    {} \ar[dd, "\tau_{\vi\vk}"'] \ar[rr, "\tau_{\vi\vl}"] \pbsquare{ddrr} && {} \ar[d, "\tau'_{\vi\vl}"] \\
                                                                                          && {} \ar[d, "\sigma'_\vi"] \\
                    {} \ar[r, "\tau'_{\vi\vk}"'] & {} \ar[r, "\sigma'_\vi"'] & {}
                \end{arrow-diagram}
                follows from
                \begin{arrow-diagram}
                    {} \ar[d, "\tau_{\vi\vk}"'] \ar[r, "\tau_{\vi\vl}"] \pbsquare{dr} & {} \ar[d, "\tau'_{\vi\vl}"] \\
                    {} \ar[r, "\tau'_{\vi\vk}"'] & {}
                \end{arrow-diagram}
                and injectivity of $\sigma'_\vi$ (pullbacks are preserved under post-composition with monos, \cref{lem:pullback-mono-post-comp}).
        \item For all $\vi \neq \vj \in \vI, \vk \in \vJ_\vi, \vl \in \vJ_\vj$, $\tau_{\vi\vk} \xto{\sigma'_\vi \comp \tau'_{\vi\vk}} \hat\sigma \xfrom{\sigma'_\vj \comp \tau'_{\vj\vl}} \tau_{\vj\vl}$ is a local independent square.
                We need
                \begin{arrow-diagram}
                    {} \ar[dd, "\tau_{\vi\vk}"'] \ar[rr, "\tau_{\vj\vl}"] \pbsquare{ddrr} && {} \ar[d, "\tau'_{\vj\vl}"] \\
                                                                                          && {} \ar[d, "\sigma'_\vj"] \\
                    {} \ar[r, "\tau'_{\vi\vk}"'] & {} \ar[r, "\sigma'_\vi"']              & {}              
                \end{arrow-diagram}
                which follows from pullback pasting
                \begin{arrow-diagram}
                    {} \ar[d, Rightarrow, "1"'] \ar[rr, "\tau_{\vj\vl}"] \pbsquare{dr} && {} \ar[d, "\tau'_{\vj\vl}"] \\
                    {} \ar[d, "\tau_{\vi\vk}"'] \ar[r, Rightarrow, "1"] \pbsquare{dr} & {} \ar[r, "\sigma_\vj"] \ar[d, "\sigma_\vi"] \pbsquare{dr}  & {} \ar[d, "\sigma'_\vj"] \\
                    {} \ar[r, "\tau'_{\vi\vk}"'] & {} \ar[r, "\sigma'_\vi"']          & {}
                \end{arrow-diagram}
                (see also \cref{lem:pullback-iso-mono})
    \end{enumerate}
\end{proof}

% \begin{remark}[Future Concern for \cref{sec:local-ind-squares-einitC}]\label{rem:failure-lis-einitC}
%     In the proof of \cref{lem:local-ind-struct-W}, case (2), we used the fact that coimages are preserved under composition
%     \[ \vg[\coimg{\vf}] \subseteq \coimg{\vg \comp \vf} \]
%     An analogous property when we attempt to define an independence structure in \cref{sec:local-ind-squares-einitC} is
%         Unreachable locations are preserved under composition in $\einitC$.
%     Unfortunately such a property does not hold, and thus the independence structure we give for $\einitC$ does not satisfy the composition axiom. 
% \end{remark}

\begin{definition}[Local Independent Coproducts from \cite{kammar2017monad}]\label{def:local-ind-coproducts-W}
Recall that a local independent coproduct of $\vw_1 \xfrom{\sigma_1} \vw \xto{\sigma_2} \vw_2$ from \textcite{kammar2017monad} is given by the following square
\begin{center}
    \begin{tikzcd}[sep=small] 
                    & {\vw} \arrow[ld, "\sigma_1"'] \arrow[rd, "\sigma_2"] & \\
        {\vw_1} \arrow[rd, "\ext^{\sigma_1}\sigma_2"'] &            & {\vw_2} \arrow[ld, "\ext^{\sigma_2}\sigma_1"] \\
                                                       &  {\vw_1 \liplus\vw \vw_2} &
    \end{tikzcd} 
\end{center}
defined by
\begin{align*}
    \vw_1 \liplus\vw \vw_2 & \eqdef \vw \oplus (\vw_1 \ominus \sigma_1) \oplus (\vw_2 \ominus \sigma_2) \\
    \ext^{\sigma_1}\sigma_2 & \eqdef 1 \oplus \incl_1 \comp \iso^{\sigma_1} \\
    \ext^{\sigma_2}\sigma_1 & \eqdef 1 \oplus \incl_2 \comp \iso^{\sigma_2}
\end{align*}
\end{definition}

\begin{lemma}[Local Independent Coproducts on $\worldC$]
    The local independent coproducts (L.I.C.) from \textcite{kammar2017monad} give the structure for the L.I.C. in \textcite{simpson2018category} (considering the independent structure above). 
\end{lemma}
\begin{proof}
    One needs to show: 
    \begin{itemize} 
        \item $\{ \ext^{\sigma_1}\sigma_2, \ext^{\sigma_2}\sigma_1 \}$ is $\mcatI$-neutral and independent (Definition 3.5).
        \item The unique morphism exists (Definition 3.5).
        \item Independent squares compose (Definition 5.2).
        \item The pasting-like law for L.I.C. (Definition 6.2).
    \end{itemize}
\end{proof}

\begin{remark}\label{rem:local-ind-coproducts-W-data}
In a L.I.C. square of span $\vw_1 \xfrom{\sigma_1} \vw \xto{\sigma_2} \vw_2$,
if $\sigma_1$ is equipped with instantiation data $\mvarM : \dheapF(\vw_1 \ominus \sigma_1, \vw_1)$,
then $\ext^{\sigma_2}\sigma_1$ can be equipped with instantiation data as well (See \cite{kammar2017monad} for definition of the heaplet bifunctor $\dheapF : \opC\worldC \times \worldC \to \setC$)
\[ \dheapF(\iso, \ext^{\sigma_1}\sigma_2)(\mvarM) : \dheapF((\vw_1 \liplus\vw \vw_2) \ominus \ext^{\sigma_2}\sigma_1, \vw_1 \liplus\vw \vw_2) \]
using the canonical iso $(\vw_1 \liplus\vw \vw_2) \ominus \ext^{\sigma_2}\sigma_1 \iso \vw_1 \ominus \sigma_1$.
\end{remark}


\subsection[Category of Instantiations]{The category of instantiations $\instC$}

The category $\instC$ is defined as
\begin{itemize}
    \item Objects: worlds $\vw = \{ \vl_1 : \vC_1, \dots, \vl_\vn : \vC_\vn \}$.
    \item Morphisms: Injective, sort-preserving $\rho : \vw \to \vw'$ and instantiation data $\mvarM : \dheapF(\vw' \ominus \rho, \vw')$.
\end{itemize}
We have a forgetful functor $\instU : \instC \to \worldC$ that forgets the instantiation data.

\begin{definition}[Inhabited Sorts]\label{def:inhabited-sorts}
    A sort $\vC \in \typeSort$ is said to be \emph{inhabited} if
    there exists a world $\vw$, renaming $\va : \vC \to \vw$, and $\mvarM : \heapF(\vw)$.
\end{definition}
Similarly we say that a set of sorts is inhabited if every sort in the set is inhabited.

\begin{lemma}\label{lem:lic-instC}
    The local independent structure on $\worldC$ can be directly extended to $\instC$ such that $\instU$ reflects local independent squares.
\end{lemma}

Recall that the hiding monad $\hideMon$ on $\instC$ is given by the colimit $\hideMon\vA\vw = \Colimit{(\sliceC\vw\instU \xto\vU \instC \xto\vA \setC)}$.
The following lemma describes the quotienting that happens by this colimit.
\begin{lemma}[Quotienting in the Hiding Monad]\label{lem:hmon-quotient}
    For $\sigma_1 : \vw \to \vw_1$, $\va_1 \in \vA(\vw_1)$, and $\sigma_2 : \vw \to \vw_2$, $\va_2 \in \vA(\vw_2)$, 
    we have that $\eqclass{\sigma_1}\va_1 = \eqclass{\sigma_2}\va_2 \in \hideMon\vA\vw$ iff there exist 
    \begin{tikzcd}
        & {\hat\vw}
    \\ 
    {\vw_1} 
        \ar[ur, "\tinstU\me_1'"] 
            && {\vw_2} 
                \ar[ul, "\tinstU\me_2'"']
    \\ 
        & {\vw} 
            \ar[ul, "\sigma_1"] \ar[ur, "\sigma_2"']
    \end{tikzcd}
    s.t. $\vA(\me_1')(\va_1) = \vA(\me_2')(\va_2)$.
\end{lemma}
\begin{proof}
    We first show that the relation
    \[ \inj_{\sigma_1}(\va_1) \sim \inj_{\sigma_2}(\va_2) \iff 
       \ExistsTwo{\me_1'}{\me_2'} \instU\me_1' \comp \sigma_1 = \instU\me_2' \comp \sigma_2 \Land \vA(\me_1')(\va_1) = \vA(\me_2')(\va_2) \]
    is an equivalence.
    Reflexivity and symmetry is immediate.
    For transitivity, we have
    \begin{center}
        \begin{tikzcd}
                && {\vw_{12} \liplus{\vw_2} \vw_{23}}
        \\
            & {\vw_{12}} 
                \ar[ur, "\ext^{\me_{22}}\me_{23}"]
                    && {\vw_{23}}
                        \ar[lu, "\ext^{\me_{23}}\me_{22}"']
        \\
        {\vw_1}
            \ar[ur, "\tinstU\me_1"]
                && {\vw_2}
                        \ar[ul, "\tinstU\me_{22}"] \ar[ur, "\tinstU\me_{23}"']
                        && {\vw_3}
                            \ar[ul, "\tinstU\me_3"']
        \\\\
                && {\vw_1}
                    \ar[uull, "\sigma_1"] \ar[uu, "\sigma_2"] \ar[uurr, "\sigma_3"']
        \end{tikzcd}
    \end{center}
    Notice that the top square commutes by \cref{lem:lic-instC}.
    which implies that $\inj_{\sigma_1}(\va_1) \sim \inj_{\sigma_3}(\va_3)$.

    To prove that the generated relation by the colimit is exactly $\sim$, one shows direction $\to$ by inducting on the list of related elements,
    and for $\from$, notice that if $\inj_{\sigma_1}(\va_1) \sim \inj_{\sigma_3}(\va_3)$, then
    \[ \inj_{\sigma_1}\va_1 \leq \inj_{\instU\me_1 \comp \sigma_1}\vA(\me_1)(\va_1) \ge \inj_{\sigma_2}\va_2 \]

\end{proof}

\subsection[Category of Templates]{The category of templates $\tempC$}

The first step towards introducing template heaps requires us to distinguish between two different kinds of locations.
We have (i) concrete locations that are associated with some cell contents and 
(ii) location variables that stand for \emph{any arbitrary cell contents}.

To model these two kinds of locations we make use of the category of templates defined as the arrow category $\arrC\worldC$
\begin{itemize}
    \item Objects: Sort-preserving injective renamings between worlds $\Minus\vw \xto\sigma \Plus\vw$.
    \item Morphisms: Injective, sort-preserving $\concL\rho : \Minus\vw \to \Plus{\vw'}$ and $\anyL\rho : \Minus\vw \to \Plus{\vw'}$, such that: 
                    \begin{center}
                    \begin{tikzcd}{\Plus\vw} \ar[r, tail, "{\anyL\rho}"] & {\Plus{\vw'}} \\
                                  {\Minus\vw} \ar[u, tail, "\sigma"] \ar[r, tail, "{\concL\rho}"'] & {\Minus{\vw'}} \ar[u, tail, "\sigma'"']
                     \end{tikzcd} 
                    \end{center}
\end{itemize}
We have two forgetful functors $\concL\dash, \anyL\dash : \tempC \to \worldC$, projecting the $-$ and $+$ components respectively.

\begin{notation}
\noindent
\begin{itemize}
    \item We refer to the elements of $\concL\sigma$ as the concrete locations of $\sigma \in \objC{\tempC}$.
    \item We refer to the elements of $\Plus\vw \ominus \sigma$ as the location variables of $\sigma \in \objC{\tempC}$, notation $\varL\sigma$.
    \item We refer to elements of $\anyL\sigma$ as \emph{all} the locations of $\sigma \in \objC{\tempC}$.
    \item If it is clear from the context we overload the notations $\concL\sigma = \concL\sigma \xto\sigma \anyL\sigma$ and
        $\varL\sigma = \varL\sigma \xto{\compl\sigma} \anyL\sigma$ for the inclusions into $\anyL\sigma$.
    \item We have an induced morphism $\varL\rho : \varL\sigma \xto{\varL\sigma} \anyL\sigma \xto{\anyL\rho} \anyL{\sigma'}$.
    \item We refer to elements $\anyL{\sigma'} \ominus \anyL\rho$ as the new locations introduced by $\rho$, 
          and $\concL{\sigma'} \ominus \concL\rho$ as the concrete locations introduced by $\rho$.
    \item We refer to the variables $\varL\sigma$ that are instantiated by $\rho$ as $\concL{\sigma'} \pltimes{\anyL{\sigma'}} \varL\sigma$: 
        \begin{center}
            \begin{tikzcd}
                {\concL{\sigma'} \pltimes{\anyL{\sigma'}} \varL\sigma} \ar[r] \ar[d] \pbsquare{dr} & {\varL\sigma} \ar[d, "\varL\rho"] \\
                {\concL{\sigma'}} \ar[r, "\concL{\sigma'}"']                    & {\anyL{\sigma'}}
            \end{tikzcd}
        \end{center}
            which are exactly the $\varL\sigma$ that are mapped to concrete locations $\concL{\sigma'}$ by $\rho$.
    \item We refer to the variables $\varL{\sigma'}$ that are fresh w.r.t. $\rho$ as $\varL{\sigma'} \pltimes{\anyL{\sigma'}} (\anyL{\sigma'} \ominus \anyL\rho)$: 
        \begin{center}
            \begin{tikzcd}
                {\varL{\sigma'} \pltimes{\anyL{\sigma'}} (\anyL{\sigma'} \ominus \anyL\rho)}
                    \ar[r] \ar[d] \pbsquare{dr} 
                        & {\anyL{\sigma'} \ominus \anyL\rho} 
                            \ar[d, "\compl{\anyL\rho}"] 
                \\
                {\varL{\sigma'}}    
                    \ar[r, "\varL{\sigma'}"'] 
                        & {\anyL{\sigma'}} 
            \end{tikzcd}
        \end{center}
            which are exactly the $\varL{\sigma'}$ that are not mapped from any location $\anyL\sigma$ by $\rho$.
    \item Similarly, $\varL{\sigma'} \pltimes{\anyL{\sigma'}} \varL\sigma$ refers to the variables $\varL{\sigma'}$ that are mapped from variables in $\varL\sigma$ by $\rho$. 
        \begin{center}
            \begin{tikzcd}
                {\varL{\sigma'} \pltimes{\anyL{\sigma'}} \varL\sigma} \ar[r] \ar[d] \pbsquare{dr} &  {\varL\sigma} \ar[d, "\varL\rho"] \\
                {\varL{\sigma'}} \ar[r, "\varL{\sigma'}"'] & {\anyL{\sigma'}} 
            \end{tikzcd}
        \end{center}
    \item Notice that from $\concL{\sigma'}$ and $\varL{\sigma'}$ being disjoint, we have that concrete locations in $\concL\sigma$ cannot become uninstantiated
        \begin{center}
            \begin{tikzcd}
                {\emptyset} \ar[rr] \ar[d] \pbsquare{drr} && {\varL{\sigma'}} \ar[d, "\varL{\sigma'}"] \\
                {\concL\sigma} \ar[r, "\concL\rho"'] & {\concL{\sigma'}} \ar[r, "\concL{\sigma'}"'] & {\anyL{\sigma'}} 
            \end{tikzcd}
        \end{center}
    \end{itemize}
\end{notation}

\begin{example}\label{ex:tempC-overview}
    Consider the following morphism in $\tempC$
    \begin{center}
    \begin{tikzcd}
        {\va : \vC, \vb : \vC'} \ar[rr, hook, "\anyL\rho = \incl^{\va,\vb \inclto \va,\vb,\vc}" yshift=5pt]               && {\va : \vC, \vb : \vC', \vc : \vC''} \\
        {\va : \vC} \ar[u, hook, "\incl^{\va \inclto \va,\vb}"] \ar[rr, hook, "\concL\rho"'] && {\va : \vC, \vb : \vC'} \ar[u, hook, "\incl^{\va,\vb \inclto \va,\vb,\vc}"']
    \end{tikzcd}
    \end{center}
    \begin{itemize}
        \item We have variables $\varL{\incl^{\va \inclto \va,\vb}} = \{ \vb : \vC' \}$, and $\varL{\incl^{\va,\vb \inclto \va,\vb,\vc}} = \{ \vc : \vC''\}$.
        \item Variables instantiated by $\rho$: 
                \[\concL{\incl^{\va,\vb \inclto \va,\vb,\vc}} \pltimes{\{ \va, \vb, \vc \}} \varL{\incl^{\va \inclto \va,\vb}} = \{ \vb : \vC' \} \]
        \item Variables $\varL{\incl^{\va,\vb \inclto \va,\vb,\vc}}$ that are fresh w.r.t. to $\rho$: 
                \[ \varL{\incl^{\va,\vb \inclto \va,\vb,\vc}} \pltimes{\{ \va,\vb,\vc\}} (\anyL{\incl^{\va,\vb \inclto \va,\vb,\vc}} \ominus \anyL\rho) = \{ \vc : \vC''\} \]
        \item Variables $\varL{\incl^{\va,\vb \inclto \va,\vb,\vc}}$ mapped from variables in $\varL{\incl^{\va \inclto \va,\vb}}$ by $\rho$: 
                \[ \varL{\incl^{\va,\vb \inclto \va,\vb,\vc}} \pltimes{\{ \va,\vb,\vc\}} \varL{\incl^{\va \inclto \va,\vb}} = \emptyset \]
    \end{itemize}
\end{example}

\begin{definition}[Variable Extension Functor]
We define the variable extension functor as a means to add $\vw$-many variables.
\[\begin{array}{r@{\,}c@{\,}l}
    \extF\dash\vw & : & \tempC \to \tempC \\
    \extF{(\concL\sigma \xto{\concL\sigma} \anyL\sigma)}{\vw} & = & \begin{tikzcd}{\anyL\sigma \mplus \vw} \\ 
                                                                                  {\concL\sigma} \ar[u, "\incl_1 \comp \concL\sigma"']
                                                                    \end{tikzcd} \\
    \extF\rho\vw & = & \pair{\concL\me}{\anyL\me \mplus 1_\vw}
\end{array}\]
Further, for every functor $\vG : \tempC \to \setC$ and $\vw$, we have a $\extNat\vG\vw : \vG \To \vG\extF\dash{\vw'}$.
\end{definition}
\begin{example}
    The action of $\extF\dash{\{ \vd : \vC''' \}}$ on morphism $\rho$ from \cref{ex:tempC-overview} is $\extF\rho{\{ \vd : \vC'''\}}$ defined as
    \begin{center}
    \begin{tikzcd}
        {\{ \va : \vC, \vb : \vC' \} \oplus \{ \vd : \vC''' \} } 
            \ar[r, "\anyL\rho \oplus 1_\vd" yshift=2pt] 
            & {\{ \va : \vC, \vb : \vC', \vc : \vC'' \} \oplus \{ \vd : \vC'''\}} 
        \\
        {\va : \vC} 
            \ar[u, hook, "\incl_1 \comp \incl^{\va \inclto \va,\vb}"] \ar[r, hook, "\concL\rho"'] 
            & {\va : \vC, \vb : \vC'} \ar[u, hook, "\incl_1 \comp \incl^{\va,\vb \inclto \va,\vb,\vc}"']
    \end{tikzcd}
    \end{center}
\end{example}

\begin{lemma}[Monomorphisms in $\tempC$]\label{lem:monos-T}
    Every morphism in $\tempC$ is mono.
\end{lemma}
\begin{proof}Follows directly from injectivity of underlying morphisms.\end{proof}

\begin{lemma}[Initial Object in $\tempC$]
    The category $\tempC$ has an initial object.
\end{lemma}
\begin{proof}Consider template $?_\emptyset : \emptyset \to \emptyset$.\end{proof}

\begin{lemma}[Equalizers in $\tempC$]
    The category $\tempC$ has all equalizers, and $\concL\dash$, $\anyL\dash$ preserve equalizers.
\end{lemma}
\begin{proof}
    The equalizer
    \begin{tikzcd}{\eqz{\rho_1}{\rho_2}} \ar[r] & {\sigma_1} \ar[r, yshift=3pt, "\rho_1"] \ar[r, yshift=-3pt, "\rho_2"'] & {\sigma_2}\end{tikzcd} 
    is given componentwise
    \begin{center}
    \begin{tikzcd}
        {\concL{\hat\sigma}} 
            \ar[r, dashed, "\concL\vh"] \ar[rrr, yshift=4pt, bend left=30, "\concL{\hat\rho}"] \ar[d, "\concL{\hat\sigma}"]
            & {\eqz{\concL{\rho_1}}{\concL{\rho_2}}} 
                \ar[rr, "\eqz{\concL{\rho_1}}{\concL{\rho_2}}"] \ar[d, dashed] 
                    && {\concL{\sigma_1}} \ar[d, "\concL{\sigma_1}"]
        \\
        {\anyL{\hat\sigma}}
            \ar[r, dashed, "\anyL\vh"] \ar[rrr, bend right=30, "\anyL{\hat\rho}"]
            & {\eqz{\anyL{\rho_1}}{\anyL{\rho_2}}} 
                \ar[rr, "\eqz{\anyL{\rho_1}}{\anyL{\rho_2}}"]
                    && {\anyL{\sigma_1}} \ar[r, yshift=3pt, "\anyL{\rho_1}"] 
                        \ar[r, yshift=-3pt, "\anyL{\rho_2}"'] 
                            & {\anyL{\sigma_2}}
    \end{tikzcd} 
    \end{center}
\end{proof}
\begin{example}
    This is an equalizer diagram 
    \begin{center}
    \begin{tikzcd}
        {?_{\vb}} \ar[r, "\pair{?_\va}{\incl_2}"] 
            & {\incl_1^{\va \to \va \oplus \vb}}
                \ar[rr, yshift=3pt, "\pair{\incl_1}{\incl_1 \oplus 1_\vb}"] \ar[rr, yshift=-3pt, "\pair{\incl_2}{\incl_2 \oplus 1_\vb}"'] 
            && {\incl_1^{\va_1 \oplus \va_2 \to \va_1 \oplus \va_2 \oplus \vb}}
    \end{tikzcd}
    \end{center}
\end{example}

\begin{lemma}[Pullbacks in $\tempC$]\label{lem:pullbacks-tempC}
    The category $\tempC$ has all pullbacks and $\concL\dash$, $\anyL\dash$ preserve pullbacks.
\end{lemma}
\begin{proof}
    The pullback of $\sigma_1 \xto{\rho_1} \hat\sigma \xfrom{\rho_2} \sigma_2$ is given componentwise
    \begin{center}
    \begin{tikzcd}
        {\concL\sigma}
            \ar[rrrr, bend left=30, "\concL{\rho_2'}"]
            \ar[ddrr, bend right=30, "\concL{\rho_1'}"']
            \ar[dr, "\concL{\hat\sigma}"]
            \ar[rr, dashed, "\concL\vh"]
                && {\concL{\sigma_1} \pltimes{\concL\sigma} \concL{\sigma_2}}
                \ar[rr, "\ext^{\concL{\rho_1}}{\concL{\rho_2}}"]
                    \ar[dr, dashed]
                    \ar[dd, "\ext^{\concL{\rho_2}}\concL{\rho_1}"' yshift=20pt]
                        && {\concL{\sigma_2}}
                            \ar[dr, "\concL{\sigma_2}"']
                            \ar[dd, "\concL\rho"' yshift=20pt]
        \\
            & {\anyL\sigma}
                \ar[rrrr, bend left=60, "\anyL{\rho_2'}" xshift=20pt]
                \ar[rr, dashed, "\anyL\vh" xshift=-40pt]
                \ar[ddrr, bend right=30, "\anyL{\rho_1'}"']
                    && {\anyL{\sigma_1} \pltimes{\anyL\sigma} \anyL{\sigma_2}}
                        \ar[rr, "\ext^{\anyL{\rho_2}}\anyL{\rho_1}" xshift=-20pt]
                        \ar[dd, "\ext^{\anyL{\rho_1}}\anyL{\rho_2}" yshift=20pt]
                                && {\anyL{\sigma_2}}
                                    \ar[dd, "\anyL{\rho_2}"]
        \\
                && {\concL{\sigma_1}}
                    \ar[rr, "\concL{\rho_1}" xshift=-40pt]
                    \ar[rd, "\ext^{\anyL{\rho_1}}\anyL{\rho_2}"']
                        && {\concL\sigma}
                            \ar[dr, "\concL\sigma"]
        \\
                    &&& {\anyL{\sigma_1}} 
                            \ar[rr, "\anyL{\rho_1}"]
                            && {\anyL\sigma}
    \end{tikzcd}
    \end{center}
\end{proof}

\begin{example}
    This is a pullback square
    \begin{center}
    \begin{tikzcd} 
        {?_\vb} 
            \ar[r, "\pair{?_{\va_1}}{\incl_2}"] 
            \ar[d, "\pair{?_{\va_2}}{\incl_2}"']
            \pbsquare{dr}
            & {\incl_1^{\va_1 \to \va_1 \oplus \vb}}
                \ar[d, "\pair{\incl_1}{1_{\va_1} \oplus \incl_2}"]
        \\
        {1_{\va_2 \oplus \vb}}
            \ar[r, "\pair{\incl_2}{\incl_2}"']
            & {1_{\va_1 \oplus (\va_2 \oplus \vb)}}
    \end{tikzcd}
    \end{center}
\end{example}

\subsubsection{Generated Cospans}

\begin{lemma}[Generated Cospans in $\tempC$]\label{lem:gen-cospan}
    Given cospan $\anyL{\sigma_1} \xto{\tau_1} \hat\vw \xfrom{\tau_2} \anyL{\sigma_2}$ in $\worldC$ 
    we generate a cospan $\sigma_1 \xto{\rho^\tau_1} \sigma^\tau \xfrom{\rho^\tau_2} \sigma_2$ in the least possible way, 
    that is, it satisfies
    \begin{itemize}
        \item Satisfy $\anyL{\sigma^\tau} = \hat\vw$, $\anyL{\rho^\tau_1} = \tau_1$ and $\anyL{\rho^\tau_2} = \tau_2$.
        \item For any other cospan $\sigma_1 \xto{\rho_1'} \hat\sigma \xfrom{\rho_2'} \sigma_2$ that factors through $\hat\vw$
                \begin{tikzcd}
                    {\anyL{\sigma_1}} \ar[dr, "\anyL{\rho_1'}"'] \ar[r, "\tau_1"] & {\hat\vw} \ar[d, dashed, "\exists \tau'"] & {\anyL{\sigma_2}} \ar[l, "\tau_2"'] \ar[dl, "\anyL{\rho_2'}"] \\
                                                                                      & {\anyL{\hat\sigma}}
                \end{tikzcd},
              there exists a unique $\vh : \sigma^\tau \to \hat\sigma$ that satisfies $\rho_1' = \vh \comp \rho^\tau_1$ and $\rho_2' = \vh \comp \rho^\tau_2$.
    \end{itemize}

    We say that $\rho^\tau_1, \rho^\tau_2$ are generated by $\tau_1, \tau_2$.
    Intuitively, a generated cospan selects the smallest $\concL{\sigma^\tau}$ possible, that is, it selects the least concrete locations possible.
\end{lemma}
\begin{proof}
    We define $\anyL{\sigma^\tau} = \hat\vw$, $\anyL{\rho^\tau_1} = \tau_1$, $\anyL{\rho^\tau_2} = \tau_2$.
    To define $\concL{\sigma^\tau}$ we first consider their common concrete locations through the pullback
    \begin{center}
    \begin{tikzcd}[sep=small]
        {\concL{\sigma_1} \pltimes{\hat\vw} \concL{\sigma_2}} 
            \ar[rrrr, "\res_{\concL{\sigma_2};\tau_2}\concL{\sigma_1};\tau_1"] 
            \ar[dd, "\res_{\concL{\sigma_1}; \tau_1}\concL{\sigma_2}; \tau_2"'] 
            \pbsquare{ddrrrr}
                &&&& {\concL{\sigma_2}} \ar[d, "\concL{\sigma_2}"] \\
                &&&& {\anyL{\sigma_2}} \ar[d, "{\tau_2}"] \\
        {\concL{\sigma_1}} \ar[r, "\concL{\sigma_1}"'] 
                & {\anyL{\sigma_1}} \ar[rrr, "{\tau_1}"']
                &&& {\hat\vw}
    \end{tikzcd}
    \end{center}
    and further take the L.I.C. in $\worldC$ of the pullback morphisms
    \begin{center}
    \begin{tikzcd}
        {\concL{\sigma_1} \pltimes{\hat\vw} \concL{\sigma_2}} 
            \ar[rr, "\res_{\concL{\sigma_2}; \tau_2}\concL{\sigma_1}; \tau_1" yshift=5pt] 
            \ar[dd, "\res_{\concL{\sigma_1}; \tau_1}\concL{\sigma_2}; \tau_2"']
                && {\concL{\sigma_2}} \ar[dd, "\concL{\rho^\tau_2}"] 
    \\\\
        {\concL{\sigma_1}} \ar[rr, "\concL{\rho^\tau_1}"']
                && {\concL{\sigma^\tau} := \concL{\sigma_1} \liplus{\concL{\sigma_1} \pltimes{\hat\vw} \concL{\sigma_2}} \concL{\sigma_2}}
                                           
    \end{tikzcd}
    \end{center}
    The morphism $\concL{\sigma^\tau} \xto{\concL{\sigma^\tau}} \anyL{\sigma^\tau}$ is the unique morphism induced by the L.I.C. towards the pullback. 
    The UP gives that the morphisms $\rho^\tau_1$, $\rho^\tau_2$ are commuting squares.

    \paragraph{Least Property}
    Consider a compatible cospan $\rho_1', \rho_2'$.
    We only need to extend $\anyL\vh = \tau'$ to a morphism in $\tempC$.
    Let $\concL\vh$ be the induced morphism by the L.I.C.
    \begin{center}
    \begin{tikzcd}
        {\concL{\sigma_1} \pltimes{\hat\vw} \concL{\sigma_2}} \ar[d, "\res_{\concL{\sigma_1}; \tau_1}\concL{\sigma_2}; \tau_2"']
                                                                              \ar[rr, "\res_{\concL{\sigma_2}; \tau_2}\concL{\sigma_1}; \tau_1"]
            && {\concL{\sigma_2}} \ar[d, "\concL{\rho^\tau_2}"] \ar[rdd, bend left=40, "\concL{\rho_2'}"] & \\
        {\concL{\sigma_1}} \ar[rr, "\concL{\rho^\tau_1}"'] \ar[rrrd, bend right=20, "\concL{\rho_1'}"'] 
            && {\concL{\sigma^\tau}} \ar[dr, dashed, "\concL\vh"] & \\
                                    &&& {\concL{\hat\sigma}}
    \end{tikzcd}
    \end{center}
    The outer square commutes by
    \begin{align*}
        \concL{\rho_2'} \comp (\res_{\concL{\sigma_2};\tau_2}\concL{\sigma_1};\tau_1) & = \concL{\rho_1'} \comp (\res_{\concL{\sigma_1};\tau_1}\concL{\sigma_2};\tau_2) \iff \\
        \concL{\hat\sigma} \comp \concL{\rho_2'} \comp (\res_{\concL{\sigma_2};\tau_2}\concL{\sigma_1};\tau_1) & = \concL{\hat\sigma} \comp \concL{\rho_1'} \comp (\res_{\concL{\sigma_1};\tau_1}\concL{\sigma_2};\tau_2) = \\
        \anyL{\rho_2'} \comp \concL{\sigma_2} \comp (\res_{\concL{\sigma_2};\tau_2}\concL{\sigma_1};\tau_1) & = \anyL{\rho_2'} \comp \concL{\sigma_1} \comp (\res_{\concL{\sigma_1};\tau_1}\concL{\sigma_2};\tau_2) = \\
        \tau' \comp \tau_2 \comp \concL{\sigma_2} \comp (\res_{\concL{\sigma_2};\tau_2}\concL{\sigma_1};\tau_1) & = \tau' \comp \tau_1 \comp \concL{\sigma_1} \comp (\res_{\concL{\sigma_1};\tau_1}\concL{\sigma_2};\tau_2)
    \end{align*}
    and satisfies independence by \cref{lem:pullback-mono-post-comp}.
    Further $\vh$ is a commuting square, that is $\concL{\hat\sigma} \comp \concL\vh = \anyL\vh \comp \concL{\sigma^\tau}$, by the jointly epi property of $\concL{\rho^\tau_1}$ and $\concL{\rho^\tau_2}$.
\end{proof}

\begin{lemma}[Generated Cospans and Jointly Epi Property]\label{lem:gen-cospan-epi-prop}
    A cospan $\sigma_1 \xto{\rho_1} \hat\sigma \xfrom{\rho_2} \sigma_2$ is 
    generated by $\anyL{\sigma_1} \xto{\anyL{\rho_1}} \anyL{\hat\sigma} \xfrom{\anyL{\rho_2}} \anyL{\sigma_2}$ 
    iff $\concL{\rho_1}$ and $\concL{\rho_2}$ are jointly surjective.
\end{lemma}
\begin{proof}\nl
    $\to$) 
    Generated cospans are unique up to isomorphism, and so $\rho_1$, $\rho_2$ is iso to the cospan induced by \cref{lem:gen-cospan} at $\anyL{\rho_1}$, $\anyL{\rho_2}$. 
    Since the cospan of \cref{lem:gen-cospan} at $\concL\dash$ is a L.I.C., we have that $\concL{\rho_1}$, $\concL{\rho_2}$ are jointly epi.

    $\from$)
    We only need to show that the cospan $\rho_1$, $\rho_2$ satisfies the UP of generated cospans, so consider
    \begin{center}
    \begin{tikzcd}
        {\anyL{\sigma_1}} 
            \ar[dr, "\anyL{\rho_1'}"'] \ar[r, "\anyL{\rho_1}"] 
            & {\anyL{\hat\sigma}} 
                \ar[d, dashed, "\exists \anyL\vh"] 
                & {\anyL{\sigma_2}} 
                    \ar[l, "\anyL{\rho_2}"'] \ar[dl, "\anyL{\rho_2'}"] 
        \\
            & {\anyL{\sigma'}}
    \end{tikzcd}
    \end{center}
    Define $\concL\vh$ as 
    \begin{align*}
        \concL\vh & : \concL{\hat\sigma} \to \concL{\sigma'} \\
        \concL\vh(\va : \vC) & \eqdef \begin{cases} \concL{\rho_1'}(\va_1)  & \;\mif\; \concL{\rho_1}(\va_1) = \va \\
                                                    \concL{\rho_2'}(\va_2)  & \;\mif\; \concL{\rho_2}(\va_2) = \va
                                      \end{cases}
    \end{align*}
    It is well-defined for if $\concL{\rho_1}(\va_1) = \va = \concL{\rho_1}(\va_2)$, we have
    \begin{align*}
        \concL{\hat\sigma}\concL{\rho_1}(\va_1) &= \concL{\hat\sigma}\concL{\rho_2}(\va_2) \iff \\
        \anyL{\rho_1}\concL{\sigma_1}(\va_1) &= \anyL{\rho_2}\concL{\sigma_2}(\va_2) \iff \\
        \anyL\vh\anyL{\rho_1}\concL{\sigma_1}(\va_1) &= \anyL\vh\anyL{\rho_2}\concL{\sigma_2}(\va_2) \iff \\
        \anyL{\rho_1'}\concL{\sigma_1}(\va_1) &= \anyL{\rho_2'}\concL{\sigma_2}(\va_2) \iff \\
        \concL{\sigma'}\concL{\rho_1'}(\va_1) &= \concL{\sigma'}\concL{\rho_1'}(\va_2) \iff \\
        \concL{\rho_1'}(\va_1) &= \concL{\rho_1'}(\va_2)
    \end{align*}
    It is easy to show that $\vh$ is a morphism in $\tempC$ and uniqueness of $\vh$ follows from the jointly epi property. 
\end{proof}

\begin{lemma}[Instantiation data for Generated Cospans]\label{lem:gen-cospan-data}
    Given a diagram of the form 
        \begin{tikzcd} 
            & {\anyL\sigma} \ar[dl, "\anyL{\rho_1}"'] \ar[dr, "\anyL{\rho_2}"] & \\ 
            {\anyL{\sigma_1}} \ar[dr, "\tau_1"'] && {\anyL{\sigma_2}} \ar[dl, "\tau_2"] \\
                                                & {\hat\vw} &
        \end{tikzcd} where $\rho_2$ has data $\mvarM : \dheapF(\concL{\sigma_2} \ominus \concL{\rho_2}, \anyL{\sigma_2})$,
        we can equip $\rho^\tau_1$ with data $\eta_{\rho^\tau_1} : \dheapF(\concL{\hat\sigma} \ominus \concL{\rho^\tau_1}, \anyL{\hat\sigma})$.
\end{lemma}
\begin{proof}
    We have a canonical isomorphism
    \[ \concL{\hat\sigma} \ominus \concL{\rho^\tau_1} \iso \concL{\sigma_2} \ominus \res_{(\concL{\sigma_2}; \tau_2)}\concL{\sigma_1};\tau_1 \]
    and further from the pullback property
    \begin{center}
    \begin{tikzcd}[sep=small]
    {\concL\sigma} \ar[rrrrrd, bend left=20, "\concL{\rho_2}"] \ar[dddr, bend right=50, "\concL{\rho_1}"'] \ar[dr, dashed, "\vh"] & 
\\
        & {\concL{\sigma_1} \pltimes{\hat\vw} \concL{\sigma_2}} 
            \ar[rrrr, "\res_{\concL{\sigma_2};\tau_2}\concL{\sigma_1};\tau_1"] 
            \ar[dd, "\res_{\concL{\sigma_1}; \tau_1}\concL{\sigma_2}; \tau_2"] 
            \pbsquare{ddrrrr}
                &&&& {\concL{\sigma_2}} \ar[d, "\concL{\sigma_2}"] 
\\
                &&&&& {\anyL{\sigma_2}} \ar[d, "{\tau_2}"] 
\\
        & {\concL{\sigma_1}} \ar[r, "\concL{\sigma_1}"'] 
            & {\anyL{\sigma_1}} \ar[rrr, "{\tau_1}"']
                &&& {\hat\vw}
    \end{tikzcd}
    \end{center}
    we have 
\[ \concL{\sigma_2} \ominus \concL{\rho_2} \iso (\concL{\sigma_2} \ominus \res_{(\concL{\sigma_2}; \tau_2)}\concL{\sigma_1};\tau_1) \mplus 
                                                (\concL{\sigma_1} \pltimes{\hat\vw} \concL{\sigma_2} \ominus \vh) \]
    Putting these together gives a morphism 
    \[ \concL{\hat\sigma} \ominus \concL{\rho^\tau_1} \xto{\incl_1} \concL{\sigma_2} \ominus \concL{\rho_2} \]
    and so we define $\mvarM_{\rho^\tau_1} = \dheapF(\incl_1, \tau_2)(\mvarM)$.
\end{proof}

\subsubsection{Local Independence Structure}

We equip $\tempC$ with a local independence structure that takes into consideration the instantiations of location variables.

\begin{definition}[Local Independent Squares in $\tempC$]
    \begin{tikzcd}[sep=small] 
                    & {\sigma} \ar[ld, "\rho_1"'] \ar[rd, "\rho_2"] & \\
        {\sigma_1} \ar[rd, "\rho_1'"'] &            & {\sigma_2} \ar[ld, "\rho_2'"] \\
                                &  {\hat\sigma} &
    \end{tikzcd} 
    is a local independent square when the pullback of $\anyL{\sigma_1} \pltimes{\anyL{\hat\sigma}} \anyL{\sigma_2}$ is $\anyL\sigma$
    \begin{center}
        \begin{tikzcd}
            {\anyL\sigma} \ar[r, "\anyL{\rho_2}"] \ar[d, "\anyL{\rho_1}"'] \pbsquare{dr} & {\anyL{\sigma_2}} \ar[d, "\anyL{\rho'_2}"] \\
            {\anyL{\sigma_1}} \ar[r, "\anyL{\rho'_1}"']                                  & {\anyL{\hat\sigma}}
        \end{tikzcd}
    \end{center}
    Thus an independent square conflates two locations $\va \in \anyL{\sigma_1}$ and $\vb \in \anyL{\sigma_2}$ only if they come from $\anyL\sigma$.
\end{definition}

\begin{lemma}[Local Independent Coproducts in $\tempC$]\label{lem:lic-exist}
    For every span $\sigma_1 \xfrom{\rho_1} \sigma \xto{\rho_2} \sigma_2$ the local independent coproduct (L.I.C.) exists.
\end{lemma}
\begin{proof}
    Consider the L.I.C. in $\worldC$
    \begin{center}
    \begin{tikzcd}
        {\anyL\sigma} \ar[d, "{\anyL{\rho_2}}"'] \ar[rr, "{\anyL{\rho_1}}"] && {\anyL{\sigma_1}} \ar[d, "{\ext^{\anyL{\rho_1}}\anyL{\rho_2}}"] \\
        {\anyL{\sigma_2}} \ar[rr, "{\ext^{\anyL{\rho_2}}\anyL{\rho_1}}"'] && {\anyL{\sigma_1} \liplus{\anyL\sigma} \anyL{\sigma_2}}
    \end{tikzcd}
    \end{center}

    Using \cref{lem:gen-cospan} for $\anyL{\sigma_1} \xto{\ext^{\anyL{\rho_1}}\anyL{\rho_2}} \anyL{\sigma_1} \liplus{\anyL\sigma} \anyL{\sigma_2} \xfrom{\ext^{\anyL{\rho_2}}\anyL{\rho_1}} \anyL{\sigma_2}$ can extend this to morphisms in $\tempC$
        \[ \sigma_1 \xto{\ext^{\rho_1}\rho_2} \sigma_1 \liplus\sigma \sigma_2 \xfrom{\ext^{\rho_2}\rho_1} \sigma_2 \]
        We have that $\ext^{\rho_1}\rho_2 \comp \rho_1 = \ext^{\rho_2}\rho_1 \comp \rho_2$ from $\ext^{\anyL{\rho_1}}\anyL{\rho_2} \comp \anyL{\rho_1} = \ext^{\anyL{\rho_2}}\anyL{\rho_1} \comp \anyL{\rho_2}$.

    \paragraph{Induced Morphism}
    Consider the diagram where each square is independent.
    \begin{center}
    \begin{tikzcd}     
                                          & {\sigma} \ar[ld, "{\rho_1}"'] \ar[rd, "{\rho_2}"] & \\
{\sigma_1} \ar[rd, pos=0.30, "{\ext^{\rho_1}\rho_2}"] \ar[rdd, "\rho_1'"'] &   & {\sigma_2} \ar[ld, pos=0.30,"{\ext^{\rho_2}\rho_1}"'] \ar[ldd, "\rho_2'"] \\
                                                             & {\sigma_1 \liplus\sigma \sigma_2} \ar[d, dashed, pos=0.20, "{\hat\rho}"] & \\
                                          & {\hat\sigma} &
    \end{tikzcd}
    \end{center}

    The $\anyL{\hat\rho}$ part of morphism $\hat\rho$ is induced by
    \begin{center}
    \begin{tikzcd}
        {\anyL\sigma} \ar[d, "{\anyL{\rho_2}}"'] \ar[rr, "{\anyL{\rho_1}}"] 
            && {\anyL{\sigma_1}} \ar[d, "{\ext^{\anyL{\rho_1}}\anyL{\rho_2}}"] \ar[rdd, bend left=40, "\anyL{\rho_1'}"] & \\
        {\anyL{\sigma_2}} \ar[rr, "{\ext^{\anyL{\rho_2}}\anyL{\rho_1}}"'] \ar[rrrd, bend right=20, "\anyL{\rho_2'}"'] 
            && {\anyL{\sigma_1} \liplus{\anyL\sigma} \anyL{\sigma_2}} \ar[dr, dashed, "\anyL{\hat\rho}"] & \\
                                                    &&& {\anyL{\hat\sigma}}
    \end{tikzcd}
    \end{center}
    which by \cref{lem:gen-cospan} extends to a unique $\hat\rho$.
    $\mcatI$-neutrality of the injections and composition of independent (product) squares follows from the independent structure on $\worldC$.
\end{proof}

\begin{example}\label{ex:tempC-lic}
The following diagram is a L.I.C. square: 
\begin{center}
    \begin{tikzcd}
                            & {?_{\va,\vb,\vc}} \ar[ld, "\rho_1 = \pair{?_{\va, \vb}}{\incl^{\va,\vb,\vc \inclto \va,\vb,\vc,\vd}}"'] 
                                                \ar[rd, "\rho_2 = \pair{?_{\va,\vc}}{\incl^{\va,\vb,\vc \inclto \va,\vb,\vc,\ve}}"] & \\
    {\incl^{\va, \vb \inclto \va,\vb,\vc,\vd}} \ar[rd, "\ext^{\rho_1}{\rho_2} = \pair{\incl^{\va,\vb \inclto \va,\vb,\vc}}{\incl^{\va,\vb,\vc,\vd \inclto \va,\vb,\vc,\vd,\ve}}"'] && 
    {\incl^{\va,\vc \inclto \va,\vb,\vc,\ve}} \ar[ld, "\ext^{\rho_2}{\rho_1} = \pair{\incl^{\va,\vc \inclto \va,\vb,\vc}}{\incl^{\va,\vb,\vc,\ve \inclto \va,\vb,\vc,\vd,\ve}}"] \\
                                                                            & {\incl^{\va,\vb,\vc \inclto \va,\vb,\vc,\vd,\ve}}
    \end{tikzcd}
\end{center}
Notice that location variables $\vd$ and $\ve$ are not conflated because they are fresh w.r.t. $?_{\va,\vb,\vc}$.
\end{example}

\begin{remark}\label{rem:lic-T-data}
Applying \cref{lem:gen-cospan-data} to the L.I.C. of morphisms $\sigma_1 \xfrom{\rho_1} \sigma \xto{\rho_2} \sigma_2$,
we have that if $\rho_1$ is equipped with instantiation data $\mvarM : \dheapF(\concL{\sigma_1} \ominus \concL{\rho_1},\anyL{\sigma_1})$,
then $\ext^{\rho_2}\rho_1$ can be equipped with instantiation data 
\[ \mvarM_{\ext^{\rho_2}\rho_1} : \dheapF(\concL{\sigma_1 \liplus\sigma \sigma_2} \ominus \concL{\ext^{\rho_2}\rho_1}, \anyL{\sigma_1 \liplus\sigma \sigma_2}) \] 
\end{remark}

\begin{example}
    Consider the square from \cref{ex:tempC-lic}.
    If morphism $\rho_1$ is equipped with instantiation data $\mvarM : \dheapF(\{\va, \vb\}, \{\va,\vb,\vc,\vd\})$,
    we equip $\ext^{\rho_2}\rho_1$ with instantiation data $\dheapF(\{\vb \}, \{\va,\vb,\vc, \vd, \ve\})$ by $\dheapF(\incl^{\vb \inclto \va,\vb}, \incl^{\va,\vb,\vc,\vd \inclto \va,\vb,\vc,\vd,\ve})(\mvarM)$.
\end{example}

\begin{lemma}[Pasting Law for Local Independent Coproducts]\label{lem:lic-pasting}
    In a diagram
    \begin{tikzcd}
        {\sigma} \ar[r, "\rho_1"] \ar[d, "\rho_1'"']
            & {\sigma_1} \ar[r, "\rho_3"] \ar[d, "\rho_2"]
                & {\sigma_3} \ar[d, "\rho_3'"]
        \\
        {\sigma_1'} \ar[r, "\rho_2'"']
            & {\sigma_2} \ar[r, "\rho_3''"']
                & {\sigma_3'}
    \end{tikzcd}
    where the left square is a L.I.C., then the whole square is a L.I.C. iff the right square is.
\end{lemma}
\begin{proof}
    By \cref{lem:gen-cospan-epi-prop}, to show that a commuting square is a L.I.C., it is enough to show that the morphisms out of the span are jointly surjective at the $\anyL\dash$, and $\concL\dash$ level, and further that we have a pullback at the $\anyL\dash$ level.

    For the $\to$ direction, it is clear that if any $\anyL{\rho_3''}$ and $\anyL{\rho_3'}$ are not jointly surjective then the whole square cannot be as well which is a contradiction.
    The same applies to the $\concL\dash$ level, and so the right square must be jointly surjective.
    We must further show that we have a pullback at the $\anyL\dash$ level, and so consider $\anyL{\rho_3''}(\vb) = \anyL{\rho_3'}(\va)$.
    From the left square being a L.I.C. we have that $\anyL{\rho_2'}$ and $\anyL{\rho_2}$ are jointly surjective.
    So if $\anyL{\rho_2}(\vb') = \vb$ then we are done. 
    Otherwise we have $\anyL{\rho_2'}(\vb') = \vb$, 
    and by the outer square being L.I.C., it is a pullback and therefor, 
    $\vb'$ and $\va$ are related by $\anyL{\sigma}$, which implies that they are related by $\anyL{\sigma_1}$ also.
    \nl
    For the $\from$ direction, consider a $\va : \vC \in \anyL{\sigma_3'}$, and we show that it is mapped from $\anyL{\sigma_1'}$ or $\anyL{\sigma_3}$.
    By the right square being jointly surjective, either $\va$ is mapped from $\anyL{\sigma_3}$ or it comes from $\anyL{\sigma_2}$.
    If it comes from $\anyL{\sigma_2}$, then by the left square being L.I.C. either it is mapped from $\anyL{\sigma_1'}$ or it is mapped from $\anyL{\sigma_1}$ and since it is a commuting square, it must also be mapped from $\anyL{\sigma_3}$.
    The same reasoning applies to $\concL\dash$.
    That the outer square is also a pullback at the $\anyL\dash$ level follows from pullback pasting.
\end{proof}

\subsubsection{Variable Preserving Morphisms}

We identify a class of morphisms in $\tempC$: morphisms that do not instantiate any location variables from $\varL{\sigma_1}$.
We call these morphisms variable preserving since every variable in $\varL{\sigma_1}$ is also a variable in $\varL{\sigma_2}$.
\begin{definition}[Variable Preserving Morphisms]
A morphism $\rho : \sigma_1 \to \sigma_2$ is variable preserving if it satisfies $\concL{\sigma_2} \pltimes{\anyL{\sigma_2}} \varL{\sigma_1} = \emptyset$.
\end{definition}
\begin{example}
    The morphism $?_{\va} \xto{\pair{?_{\vb}}{\incl^{\va \inclto \va,\vb,\vc}}} \incl^{\vb \inclto \va,\vb,\vc}$ is a variable preserving morphism since variable $\va$ from the domain is mapped to a variable in the codomain.
    The morphism $?_\va \xto{\pair{?_\va}{1_\va}} 1_\va$ is not a variable preserving morphism since variable $\va$ is instantiated by the morphism.
\end{example}

\begin{lemma}[Variable Preserving Property Equivalence]\label{lem:vpres-prop}
    Let $\rho : \sigma_1 \to \sigma_2$. We have the following equivalences

    \begin{tikzcd}
        {\emptyset} 
            \ar[r] \ar[d] \pbsquare{dr} 
            & {\varL{\sigma_1}} 
                \ar[d, "\varL\rho"'] 
        \\
        {\concL{\sigma_2}} 
            \ar[r, "\concL{\sigma_2}"'] 
            & {\anyL{\sigma_2}}
    \end{tikzcd} $\xLeftrightarrow{(1)}$
    \begin{tikzcd}
        {\anyL{\sigma_1}} 
            \ar[r, "\anyL\rho"]
            & {\anyL{\sigma_2}}
        \\
        {\concL{\sigma_1}} 
            \ar[u, "\concL{\sigma_1}"] \ar[r, "\concL\rho"] \pbsquare{ur}
            & {\concL{\sigma_2}} 
                \ar[u, "\concL{\sigma_2}"]
    \end{tikzcd} $\xLeftrightarrow{(2)}$
    \begin{tikzcd}
        {\varL{\sigma_1}} 
            \ar[dr, "\varL\rho"'] \ar[r, dashed, "\exists \varL\rho\res|{\sigma_2}"] 
            & {\varL{\sigma_2}} 
                \ar[d, "\varL{\sigma_2}"] 
        \\
            & {\anyL{\sigma_2}}
    \end{tikzcd} $\xLeftrightarrow{(3)}$
    \begin{tikzcd}
        {\concL{\sigma_2} \ominus \concL\rho} 
            \ar[d, dashed, "\exists \concL{\sigma_2}\res|\rho"] \ar[r, "\compl{\concL\rho}"] 
            & {\concL{\sigma_2}} 
                \ar[d, "\concL{\sigma_2}"'] 
        \\
        {\anyL{\sigma_2} \ominus \anyL\rho} 
            \ar[r, "\compl{\anyL\rho}"'] 
            & {\anyL{\sigma_2}}
    \end{tikzcd}

\end{lemma}
\begin{proof}\nl
    For (1), consider the canonical isomorphism $\anyL{\sigma_1} \iso \concL{\sigma_1} \mplus \varL{\sigma_1}$.

    For $\xRightarrow{(2)}$ consider $\varL\rho : \varL{\sigma_1} \to \anyL{\sigma_2} \iso \concL{\sigma_2} \mplus \varL{\sigma_2}$.
    If $\varL\rho$ maps to $\concL{\sigma_2}$ we get a contradiction from the pullback.

    For $\xLeftarrow{(2)}$, we have the diagram 
    \begin{tikzcd}
        {\emptyset} \ar[d] \ar[r] & {\varL{\sigma_1}} \ar[d, "\varL\rho"'] \ar[r, dashed, "\varL\rho\res|{\sigma_2}"] & {\varL{\sigma_2}} \ar[dl, "\varL{\sigma_2}"] \\
        {\concL{\sigma_2}} \ar[r, "\concL{\sigma_2}"] & {\anyL{\sigma_2}}
    \end{tikzcd}
    which is a pullback since $\img{\varL{\sigma_2}} \setInt \img{\concL{\sigma_2}} = \emptyset$ which makes the inner square a pullback.

    For $\xRightarrow{(3)}$ the second pullback gives $\concL{\sigma_2}\compl{\concL\rho}(\va : \vC) \not\in \img{\anyL\rho}$.

    For $\xLeftarrow{(3)}$, we show $\concL{\sigma_2} \pltimes{\anyL{\sigma_2}} \varL{\sigma_1} = \emptyset$, and so consider $\varL\rho(\va : \vC) = \concL{\sigma_2}(\vb : \vC)$.
    From $\rho$ being a commuting square we have $\vb : \vC \in \concL{\sigma_2} \ominus \concL\rho$. 
    But then $\concL{\sigma_2}(\vb) \in \anyL{\sigma_2} \ominus \anyL\rho$ which is a contradiction as it must also be equal to $\varL\rho(\va)$.
\end{proof}

\begin{notation}
    For a variable preserving morphism $\sigma_1 \xto\rho \sigma_2$, 
    let $\compl\rho : (\sigma_2 \ominus \rho) \to \sigma_2$ denote the morphism in $\tempC$ given by $\xRightarrow{(3)}$ in \cref{lem:vpres-prop}.
\end{notation}

\begin{lemma}[Variable Preservation Closed Under Composition]\label{lem:comp-vpres}
    Variable preserving morphisms are closed under composition.
\end{lemma}
\begin{proof}
    Follows directly by pullback pasting.
\end{proof}

\begin{lemma}[Variable Preserving Cancellation]\label{lem:vpres-cancl}
    In \begin{tikzcd}
        {\sigma_1} 
            \ar[rr, "\rho"] \ar[dr, "\rho_1"']
                && {\sigma_3}
        \\
            & {\sigma_2}
                \ar[ur, "\rho_2"']
    \end{tikzcd}
    if $\rho$ and $\rho_2$ are variable preserving, so is $\rho_1$.
\end{lemma}
\begin{proof}
    Follows directly by pullback pasting.
\end{proof}

\begin{lemma}[Equalizers are Variable Preserving]\label{lem:eq-vpres-tempC}
   In an equalizer diagram
    \begin{center}
    \begin{tikzcd}
        {\sigma_1} \ar[r, "\ve"] & {\sigma_2} \ar[rr, shift left=1,"\rho_1"] \ar[rr, shift right=1, "\rho_2"'] && {\sigma_3}
    \end{tikzcd}
    \end{center}
    $\ve$ is variable preserving.
\end{lemma}
\begin{proof}\nl
    We want to show $\concL{\sigma_2} \pltimes{\anyL{\sigma_2}} \varL{\sigma_1} = \emptyset$ and so consider
    \begin{center}
        \begin{tikzcd}
            {\vw} \ar[rrd, bend left=30, "\tau_2"] \ar[ddr, bend right=30, "\tau_1"'] && \\
                                                          & {\emptyset} \ar[r] \ar[d]                      & {\varL{\sigma_1}} \ar[d, "\varL\ve"] \\
                                                          & {\concL{\sigma_2}} \ar[r, "\concL{\sigma_2}"'] & {\anyL{\sigma_2}}
        \end{tikzcd}
    \end{center}
    Notice that the following square $\ve'$ commutes
    \begin{center}
        \begin{tikzcd}
            {\concL{\sigma_1} \mplus \varL{\sigma_1} \iso \anyL{\sigma_1}} \ar[r, "\anyL\ve"] & {\anyL{\sigma_2}} \\
            {\concL{\sigma_1} \mplus \vw} \ar[u, "1 \mplus \tau_2"] \ar[r, "\copair{\concL\ve}{\tau_1}"'] & {\concL{\sigma_2}} \ar[u, "\concL{\sigma_2}"']
        \end{tikzcd}
    \end{center}
    and the copairing $\copair{\concL\ve}{\tau_1}$ is mono
    \begin{align*}
        \copair{\concL\ve}{\tau_1} \comp \theta_1 & = \copair{\concL\ve}{\tau_1} \comp \theta_2 \iff \\
        \concL{\sigma_2} \comp \copair{\concL\ve}{\tau_1} \comp \theta_1 & = \concL{\sigma_2} \comp \copair{\concL\ve}{\tau_1} \comp \theta_2 \iff \\
        \anyL\ve \comp 1 \oplus \tau_2 \comp \theta_1 & = \anyL\ve \comp 1 \oplus \tau_2 \comp \theta_2 \iff \\
        1 \oplus \tau_2 \comp \theta_1 & = 1 \oplus \tau_2 \comp \theta_2 \iff \\
        \theta_1 & = \theta_2
    \end{align*}

    Further notice that $\ve'$ commutes with $\rho_1$, $\rho_2$ from $\anyL{\rho_1} \comp \anyL\ve = \anyL{\rho_1} \comp \anyL\ve$
    \begin{align*}
        \concL{\rho_1} \comp \tau_1 &= \concL{\rho_2} \comp \tau_1 \iff \\
        \concL{\sigma_3} \comp \concL{\rho_1} \comp \tau_1   &= \concL{\sigma_3} \comp \concL{\rho_2} \comp \tau_1 \iff \\
        \anyL{\rho_1} \comp \concL{\sigma_2} \comp \tau_1      &= \anyL{\rho_2} \comp \concL{\sigma_2} \comp \tau_1 \iff \\
        \anyL{\rho_1} \comp \varL\ve \comp \tau_2  &= \anyL{\rho_2} \comp \varL\ve \comp \tau_2 =  \\
        \anyL{\rho_1} \comp \anyL\ve \comp \varL{\sigma_1} \comp \tau_2  &= \anyL{\rho_2} \comp \anyL\ve \comp \varL{\sigma_1} \comp \tau_2
    \end{align*}
    and so the universal property of the equalizer gives us an injective morphism $\concL{\sigma_1} \mplus \vw \to \concL{\sigma_1}$, which means $\vw = \emptyset$. 
\end{proof}

\begin{lemma}[Regular Monos in $\tempC$]\label{lem:rmono-vpres-tempC}
   A morphism $\rho : \sigma_1 \to \sigma_2$ in $\tempC$ is regular mono iff $\rho$ is variable preserving. 
\end{lemma}
\begin{proof}
    \nl
    $\to$) Follows directly from \cref{lem:eq-vpres-tempC}.
    \nl
    $\from$)
    We show that this is an equalizer diagram: 
    \begin{center}
    \begin{tikzcd}
        {\sigma_1} \ar[r, "{\rho}"] & {\sigma_2} \ar[rr, shift left=1,"{\ext^\rho\rho}_l"] \ar[rr, shift right=1, "{\ext^\rho\rho}_r"'] && {\sigma_2 \liplus{\sigma_1} \sigma_2}
    \end{tikzcd}
    \end{center}
    
    From \cref{lem:vpres-prop} and \cref{lem:lic-vpres} we have $\concL{\sigma_2}\res|\rho$ and we can simplify the diagram to
    \begin{center}
    \begin{tikzcd}
            & {\anyL{\sigma_1} \mplus (\anyL{\sigma_2} \ominus \anyL\rho)} 
            \ar[rr, shift left=1, "1 \mplus \incl_1"] 
            \ar[rr, shift right=1, "1 \mplus \incl_2"'] 
                && {\anyL{\sigma_1} \mplus (\anyL{\sigma_2} \ominus \anyL\rho) \mplus (\anyL{\sigma_2} \ominus \anyL\rho)} & 
\\
    {\sigma_2} \ar[r, phantom, pos=0.20, "\iso"] & {} && {} \ar[r, phantom, pos=0.95, "\iso"] & {\sigma_2 \liplus{\sigma_1} \sigma_2} 
\\
            & {\concL{\sigma_1} \mplus (\concL{\sigma_2} \ominus \concL\rho)} 
            \ar[uu, "\concL{\sigma_1} \mplus \concL{\sigma_2}\res|\rho"'] 
            \ar[rr, shift left=1, "1 \mplus \incl_1"] 
            \ar[rr, shift right=1, "1 \mplus \incl_2"']    
                                                        && {\concL{\sigma_1} \mplus (\concL{\sigma_2} \ominus \concL\rho) \mplus (\concL{\sigma_2} \ominus \concL\rho)} 
        \ar[uu, "\concL{\sigma_1} \mplus \concL{\sigma_2}\res|\rho \mplus \concL{\sigma_2}\res|\rho"'] & \\
    \end{tikzcd}
    \end{center}
    Consider $(1 \mplus \incl_1) \comp \iso \comp \phi = (1 \mplus \incl_2) \comp \iso \comp \phi$. Then $\phi$ must factor through $\rho$ and uniqueness follows from $\rho$ being mono.
\end{proof}

\begin{lemma}[Isos are Variable Preserving]\label{lem:iso-vpres}
    All isomorphisms are variable preserving morphisms.
\end{lemma}
\begin{proof} Follows from \cref{lem:rmono-vpres-tempC}. \end{proof}

\begin{lemma}[Local Independent Coproducts and Variable Preserving Property]\label{lem:lic-vpres}
    In a L.I.C. where $\rho_2$ is a variable preserving morphism
    \begin{center}
    \begin{tikzcd}[sep=small]     
                                                     & {\sigma} \ar[ld, "{\rho_1}"'] \ar[rd, "{\rho_2}" black]  & \\
        {\sigma_1} \ar[rd, "{\ext^{\rho_1}\rho_2}"'] &                                                                     & {\sigma_2} \ar[ld, "{\ext^{\rho_2}\rho_1}"] \\
                                                     &              {\sigma_1 \liplus\sigma \sigma_2}                      &
    \end{tikzcd}
    \end{center}
    we have that $\sigma_1 \liplus\sigma \sigma_2 \iso \concL{\sigma_1} \liplus{\concL\sigma} \concL{\sigma_2} \xdashto{\concL{\sigma_1} \mplus \concL{\sigma_2}\res|{\rho_2}}
                                                                                                       \anyL{\sigma_1} \liplus{\anyL\sigma} \anyL{\sigma_2}$ and
                 $\ext^{\rho_1}\rho_2$ is variable preserving as well.
\end{lemma}
\begin{proof}\nl
    Notice that this is a pullback square
    \begin{center}
    \begin{tikzcd}[sep=small]
        {\concL\sigma} \ar[rrrr, "\concL{\rho_1}"] \ar[dd, "\concL{\rho_2}"'] \pbsquare{ddrrrr}
                  &&&& {\concL{\sigma_1}} \ar[d, "\concL{\sigma_1}"] \\
                  &&&& {\anyL{\sigma_1}} \ar[d, "\ext^{\anyL{\rho_1}}\anyL{\rho_2}"] \\
        {\concL{\sigma_2}} \ar[r, "\concL{\sigma_2}"'] 
                  & {\anyL{\sigma_2}} \ar[rrr, "\ext^{\anyL{\rho_2}}\anyL{\rho_1}"']
                  &&& {\anyL{\sigma_1} \liplus{\anyL\sigma} \anyL{\sigma_2}}
    \end{tikzcd}
    \end{center}
    By noticing the canonical isomorphism 
    \[ \anyL{\sigma_1} \liplus{\anyL\sigma} \anyL{\sigma_2} \iso \concL{\sigma} \oplus \varL\sigma \oplus (\anyL{\sigma_1} \ominus \anyL{\rho_1})
                                                                                               \oplus (\anyL{\sigma_2} \ominus \anyL{\rho_2}) \]
    and that $\ext^{\anyL{\rho_1}}\anyL{\rho_2} \comp \concL{\sigma_1}$ maps to $\concL\sigma$, $\varL\sigma$ and $\anyL{\sigma_1} \ominus \anyL{\rho_1}$,
    whereas $\ext^{\anyL{\rho_2}}\anyL{\rho_1} \comp \concL{\sigma_2}$ maps to $\concL\sigma$ and $(\anyL{\sigma_2} \ominus \anyL{\rho_2})$ only (via $\concL{\sigma_2}\res|{\rho_2}$).

    That the induced morphism $\concL{\sigma_1} \liplus{\concL{\sigma}} \concL{\sigma_2}$ is $\concL{\sigma_1} \mplus \concL{\sigma_2}\res|{\rho_2}$ follows from uniqueness 
    ($\concL{\sigma_1}$ here has the isomorphic type $\concL\sigma \mplus (\concL{\sigma_1} \ominus \concL{\rho_1}) \to \anyL\sigma \mplus (\anyL{\sigma_1} \ominus \anyL{\rho_1})$). 

    That $\ext^{\rho_1}\rho_2$ is variable preserving follows from
    \begin{center}
    \begin{tikzcd}
        {(\concL{\sigma_1} \liplus{\concL\sigma} \concL{\sigma_2}) \ominus \ext^{\concL{\rho_1}}\concL{\rho_2}} 
            \ar[dd, bend right=40, "\iso"', end anchor=west] \ar[r, "\compl{\ext^{\concL{\rho_1}}\concL{\rho_2}}" yshift=5pt]
            & {\concL{\sigma_1} \liplus{\concL\sigma} \concL{\sigma_2}} 
                \ar[r, "\concL{\sigma_1} \liplus{\concL\sigma} \concL{\sigma_2}"]
                & {\anyL{\sigma_1} \liplus{\anyL\sigma} \anyL{\sigma_2}} 
        \\
        {\concL{\sigma_2}} 
            \ar[r, "\concL{\sigma_2}"] \ar[ur, "\ext^{\concL{\sigma_2}}\concL{\sigma_1}"]
            & {\anyL{\sigma_2}} 
                \ar[ru, "\ext^{\anyL{\rho_2}}\anyL{\rho_1}"']
        \\
        {\concL{\sigma_2} \ominus \concL{\rho_2}} 
            \ar[u, "\compl{\concL{\rho_2}}"'] \ar[r, "\concL{\sigma_2}\res|{\rho_2}"'] 
            & {\anyL{\sigma_2} \ominus \anyL{\rho_2}} 
                \ar[u, "\compl{\anyL{\rho_2}}"'] \ar[r, "\iso"']
                & {(\anyL{\sigma_1} \liplus{\anyL\sigma} \anyL{\sigma_2}) \ominus \ext^{\anyL{\rho_1}}\anyL{\rho_2}} 
                    \ar[uu, "\compl{\ext^{\anyL{\rho_1}}\anyL{\rho_2}}"']
    \end{tikzcd}
    \end{center}
\end{proof}

\begin{lemma}[Pullbacks and Variable Preserving Property]\label{lem:pullback-vpres}
    Let $\rho_2$ be variable preserving in 
    \begin{center}
    \begin{tikzcd}[sep=small]     
            & {\sigma} \ar[ld, "\rho_1"'] \ar[rd, "\rho_2"] 
        \\
        {\sigma_1} \ar[rd, "\rho_1'"'] 
                && {\sigma_2} \ar[ld, "\rho_2'"] 
        \\
            & {\hat\sigma}
    \end{tikzcd}
    \end{center}
    Then the above diagram is a pullback iff its $\anyL\dash$-image is a pullback.
\end{lemma}
\begin{proof}
    The $\to$ direction is immediate by \cref{lem:pullbacks-tempC}.
    For $\from$, by \cref{lem:pullback-mono-post-comp} it is enough to show that
    \begin{tikzcd}[sep=small]
        {\concL\sigma} 
            \ar[rr, "\concL{\rho_1}"] \ar[dd, "\concL{\rho_2}"]
            && {\concL{\sigma_1}}
                \ar[d, "\concL{\rho_1'}"]
        \\
            && {\concL{\hat\sigma}}
                \ar[d, "\concL{\hat\sigma}"]
        \\
        {\concL{\sigma_2}}
            \ar[r, "\concL{\rho_2'}"]
                & {\concL{\hat\sigma}}
                    \ar[r, "\concL{\hat\sigma}"]
                    & {\anyL{\hat\sigma}}
    \end{tikzcd}
    is a pullback which is the same as
    \begin{tikzcd}[sep=small]
        {\concL\sigma} 
            \ar[rr, "\concL{\rho_1}"] \ar[dd, "\concL{\rho_2}"]
            && {\concL{\sigma_1}}
                \ar[d, "\concL{\sigma_1}"]
        \\
            && {\anyL{\sigma_1}}
                \ar[d, "\anyL{\rho_1'}"]
        \\
        {\concL{\sigma_2}}
            \ar[r, "\concL{\sigma_2}"]
                & {\anyL{\sigma_2}}
                    \ar[r, "\anyL{\rho_2'}"]
                    & {\anyL{\hat\sigma}}
    \end{tikzcd}
    which is given by the pullback
    \begin{tikzcd}
        {\concL\sigma} 
            \ar[r, "\concL\sigma"] \ar[d, "\concL{\rho_2}"'] \pbsquare{dr}
            & {\anyL\sigma}
                \ar[r, "\anyL{\rho_1'}"] \ar[d, "\anyL{\rho_2}"] \pbsquare{dr}
                    & {\anyL{\sigma_1}}
                        \ar[d, "\anyL{\rho_1'}"]
        \\
        {\concL{\sigma_2}}
            \ar[r, "\concL{\sigma_2}"] 
                & {\anyL{\sigma_2}}
                    \ar[r, "\anyL{\rho_2'}"]
                    & {\anyL{\hat\sigma}}
    \end{tikzcd} and \cref{lem:pullback-factoring}.
\end{proof}

\subsubsection{Variable Reflecting Morphisms}

We identify another class of morphisms in $\tempC$: morphisms that do not introduce location variables.
We call these morphisms variable reflecting since every variable in $\varL{\sigma_2}$ comes from $\varL{\sigma_1}$.
\begin{definition}[Variable Reflecting Morphisms]
A morphism $\rho : \sigma_1 \to \sigma_2$ is variable reflecting if it satisfies $(\anyL{\sigma_2} \ominus \varL\rho) \pltimes{\anyL{\sigma_2}} \varL{\sigma_2} = \emptyset$.
\end{definition}
\begin{example}
    The morphism $?_{\va,\vb} \xto{\pair{?_\va}{1_{\va,\vb}}} \incl^{\va \inclto \va,\vb}$ is a variable reflecting morphism since variable $\vb$ from the codomain comes from the domain.
    The morphism $?_\va \xto{\pair{?_\va}{\incl^{\va \inclto \va,\vb}}} \incl^{\va \inclto \va, \vb}$ is not a variable reflecting morphism since variable $\vb$ is fresh.
\end{example}

\begin{lemma}[Variable Reflecting Property Equivalence]\label{lem:vrefl-prop}
    Let $\rho : \sigma_1 \to \sigma_2$. We have the following equivalences
    \begin{center}
    \begin{tikzcd}
        {\emptyset} \ar[r] \ar[d] \pbsquare{dr} & {\anyL{\sigma_2} \ominus \varL\rho} \ar[d, "\compl{\varL\rho}"'] \\
        {\varL{\sigma_2}} \ar[r, "\varL{\sigma_2}"'] & {\anyL{\sigma_2}}
    \end{tikzcd} $\xLeftrightarrow{(1)}$
    \begin{tikzcd}
        {\varL{\sigma_2}} \ar[rr, dashed, "\exists \inv^{\varL\rho}\res|{\sigma_2}"] \ar[drr, "\varL{\sigma_2}"'] && {\varL{\sigma_1}} \ar[d, "\varL\rho"'] \\
                                                                                                                   && {\anyL{\sigma_2}}
    \end{tikzcd} $\xLeftrightarrow{(2)}$
    \begin{tikzcd}
        {\concL{\sigma_2} \ominus \concL\rho} \ar[r, "\compl{\concL\rho}"] & {\concL{\sigma_2}} \ar[d, "\concL{\sigma_2}"] \\
        {\anyL{\sigma_2} \ominus \anyL\rho} \ar[u, dashed, "\exists \inv^{\concL{\sigma_2}}\res|\rho"] \ar[r, "\compl{\anyL\rho}"'] & {\anyL{\sigma_2}}
    \end{tikzcd}
    \end{center}
\end{lemma}
\begin{proof}
    For $\xRightarrow{(1)}$, consider $\varL\rho : \varL{\sigma_1} \to \anyL{\sigma_2} \iso \concL{\sigma_2} \mplus \varL{\sigma_2}$.
    If there exists $\va : \vC \in \varL{\sigma_2}$ that $\varL\rho$ does not map to, it must also be an element of $\anyL{\sigma_2} \ominus \varL\rho$ by definition. 
    But then the pullback gives us a contradiction.
    For $\xLeftarrow{(1)}$, we have a diagram
    \begin{center}
    \begin{tikzcd}
        {\emptyset} \ar[d] \ar[r] & {\anyL{\sigma_2} \ominus \varL\rho} \ar[d, "\compl{\varL\rho}"] \\
        {\varL{\sigma_2}} \ar[r, "\varL{\sigma_2}"] \ar[dr, "\inv^{\varL\rho}\res|{\sigma_2}"'] & {\anyL{\sigma_2}} \\
                                                                                                & {\varL{\sigma_1}} \ar[u, "\varL\rho"']
    \end{tikzcd}
    \end{center}
    which is a pullback since $\img{\varL\rho} \setInt \img{\compl{\varL\rho}} = \emptyset$.
    For $\xRightarrow{(2)}$, first notice that this is a pullback
    \begin{center}
    \begin{tikzcd}
        {\emptyset} \ar[d] \ar[r] \pbsquare{dr} & {\anyL{\sigma_2} \ominus \varL\rho} \ar[d, "\compl{\anyL\rho}"] \\
        {\varL{\sigma_2}} \ar[r, "\varL{\sigma_2}"'] & {\anyL{\sigma_2}}
    \end{tikzcd}
    \end{center}
    given by the inclusion $\anyL{\sigma_2} \ominus \anyL\rho \inclto \anyL{\sigma_2} \ominus \varL\rho$ and pullback pasting.
    Now consider the morphism
    \[ \anyL{\sigma_2} \ominus \anyL\rho \xto{\compl{\anyL\rho}} \anyL{\sigma_2} \iso \concL{\sigma_1} \mplus (\concL{\sigma_2} \ominus \concL\rho) \mplus \varL{\sigma_2} \]
    From $\rho$ being a commuting square it cannot map to $\concL{\sigma_1}$, and from the pullback it can also not map to $\varL{\sigma_2}$.
    Lastly for $\xLeftarrow{(2)}$, consider $\varL\rho : \varL{\sigma_1} \to \anyL{\sigma_2} \iso \concL{\sigma_2} \mplus \varL{\sigma_2}$, and a $\va : \vC \in \varL{\sigma_2}$ that is not mapped to from $\varL\rho$.
    This gives us that either $\va \in \anyL{\sigma_2} \ominus \anyL\rho$ which is a contradiction from $\inv^{\concL{\sigma_2}}\res|\rho$
    or it comes from $\concL{\sigma_1}$ which is a contradiction from $\rho$ being a commutative square.

\end{proof}

\begin{lemma}[Composition of Variable Reflecting Morphisms]\label{lem:comp-vrefl}
    Variable reflecting morphisms are closed under composition.
\end{lemma}
\begin{proof}
    For variable reflecting morphisms $\sigma_1 \xto{\rho_1} \sigma_2 \xto{\rho_2} \sigma_3$ we have
    \begin{center}
    \begin{tikzcd}
    {\varL{\sigma_3}} \ar[rd, "\varL{\sigma_3}"'] \ar[r, dashed, "\inv^{\varL{\rho_2}}\res|{\sigma_3}"] 
        & {\varL{\sigma_2}} \ar[d, "\varL{\rho_2}"'] \ar[rd, "\varL{\sigma_2}"'] \ar[r, dashed, "\inv^{\varL{\rho_1}}\res|{\sigma_2}"] 
            & {\varL{\sigma_1}} \ar[d, "\varL{\rho_1}"] \ar[dd, bend left=70, "\varL{\rho_2 \comp \rho_1}"]
    \\
        & {\anyL{\sigma_3}} \ar[rd, Rightarrow]
            & {\anyL{\sigma_2}} \ar[d, "\anyL{\rho_2}"]
    \\
           && {\anyL{\sigma_3}}
    \end{tikzcd}
    \end{center}\vspace{-1.5\baselineskip}
\end{proof}

\begin{lemma}[Variable Reflecting Cancellation]\label{lem:vrefl-cancl}
    In \begin{tikzcd}
        {\sigma_1} 
            \ar[rr, "\rho"] \ar[dr, "\rho_1"']
                && {\sigma_3}
        \\
            & {\sigma_2}
                \ar[ur, "\rho_2"']
    \end{tikzcd}
    if $\rho$ is variable reflecting, so is $\rho_2$.
\end{lemma}
\begin{proof}
    We show
    \begin{tikzcd}
        {\concL{\sigma_3} \ominus \concL{\rho_2}} \ar[r, "\compl{\concL{\rho_2}}"] & {\concL{\sigma_3}} \ar[d, "\concL{\sigma_3}"] \\
        {\anyL{\sigma_3} \ominus \anyL{\rho_2}} \ar[u, dashed, "\exists \inv^{\concL{\sigma_3}}\res|{\rho_2}"] \ar[r, "\compl{\anyL{\rho_2}}"'] & {\anyL{\sigma_3}}
    \end{tikzcd}.
    The variable reflecting property of $\rho$ gives us a $(\anyL{\sigma_3} \ominus \anyL\rho) \dashto (\concL{\sigma_3} \ominus \concL\rho)$
    and using the canonical isos we have
    \[ \anyL{\sigma_3} \ominus \anyL{\rho_2} \to 
       \anyL{\sigma_3} \ominus \anyL\rho     \to 
       \concL{\sigma_3} \ominus \concL\rho   \iso 
       (\concL{\sigma_3} \ominus \concL{\rho_2}) \oplus (\concL{\sigma_2} \ominus \concL{\rho_1}) \] 
    and its image must lie in $\concL{\sigma_3} \ominus \concL{\rho_2}$ by the commuting square of $\rho_2$.
\end{proof}

\begin{lemma}[Variable Preserving with Reflecting Cancellation]\label{lem:vpres-vrefl-cancl}
    In \begin{tikzcd}
        {\sigma_1} 
            \ar[rr, "\rho"] \ar[dr, "\rho_1"']
                && {\sigma_3}
        \\
            & {\sigma_2}
                \ar[ur, "\rho_2"']
    \end{tikzcd}
    if $\rho$ is variable preserving and $\rho_1$ is reflecting, then $\rho_2$ is variable preserving.
\end{lemma}
\begin{proof}
The variable preserving property of $\rho$ gives us a $\concL{\sigma_3} \ominus \concL\rho \dashto \anyL{\sigma_3} \ominus \anyL\rho$,
the variable reflecting property of $\rho_1$ gives us a $\anyL{\sigma_2} \ominus \anyL{\rho_1} \dashto \concL{\sigma_2} \ominus \concL{\rho_1}$,
and together with the canonical isos we have
\[ 
    \concL{\sigma_3} \ominus \concL{\rho_2} \to
        \concL{\sigma_3} \ominus \concL\rho \to
            \anyL{\sigma_3} \ominus \anyL\rho \iso
                (\anyL{\sigma_3} \ominus \anyL{\rho_2}) \oplus (\anyL{\sigma_2} \ominus \anyL{\rho_1}) \to
                    (\anyL{\sigma_3} \ominus \anyL{\rho_2}) \oplus (\concL{\sigma_2} \ominus \concL{\rho_1})
\]
and its image must lie in $\anyL{\sigma_3} \ominus \anyL{\rho_2}$.
\end{proof}

\begin{lemma}[Variable Reflecting with Preserving Cancellation]\label{lem:vrefl-vpres-cancl}
    In \begin{tikzcd}
        {\sigma_1} 
            \ar[rr, "\rho"] \ar[dr, "\rho_1"']
                && {\sigma_3}
        \\
            & {\sigma_2}
                \ar[ur, "\rho_2"']
    \end{tikzcd}
    if $\rho$ is variable reflecting and $\rho_2$ is variable preserving, then $\rho_1$ is variable reflecting.
\end{lemma}
\begin{proof}
    The variable reflecting property of $\rho$ gives us a $(\anyL{\sigma_3} \ominus \anyL\rho) \dashto (\concL{\sigma_3} \ominus \concL\rho)$,
    and the variable preserving property of $\rho_2$ gives us a $\concL{\sigma_3} \ominus \concL{\rho_2} \dashto \anyL{\sigma_3} \ominus \anyL{\rho_2}$,
    and together with the canonical isos we have
    \[ 
        \anyL{\sigma_2} \ominus \anyL{\rho_1} \to 
            \anyL{\sigma_3} \ominus \anyL\rho \to 
                \concL{\sigma_3} \ominus \concL\rho \iso
                    (\concL{\sigma_3} \ominus \concL{\rho_2}) \oplus (\concL{\sigma_2} \ominus \concL{\rho_1}) \to
                        (\anyL{\sigma_3} \ominus \anyL{\rho_2}) \oplus (\concL{\sigma_2} \ominus \concL{\rho_1})
    \]
    and its image must lie in $\concL{\sigma_2} \ominus \concL{\rho_1}$.
\end{proof}

\begin{lemma}[Local Independent Coproducts and Variable Reflecting Morphisms]\label{lem:lic-vrefl}
    In a L.I.C. square where $\rho_2$ is variable reflecting: 
    \begin{center}
    \begin{tikzcd}[sep=small]     
                                                     & {\sigma} \ar[ld, "{\rho_1}"'] \ar[rd, "{\rho_2}" black]  & \\
        {\sigma_1} \ar[rd, "{\ext^{\rho_1}\rho_2}"'] &                                                          & {\sigma_2} \ar[ld, "{\ext^{\rho_2}\rho_1}"] \\
                                                     &              {\sigma_1 \liplus\sigma \sigma_2}               &
    \end{tikzcd}
    \end{center}
    we have that $\ext^{\rho_1}\rho_2$ is variable reflecting as well.
\end{lemma}
\begin{proof}
    Consider
    \begin{center}
    \begin{tikzcd}[sep=small]
    {\concL\sigma} 
        \ar[rrrrd, bend left=15, "\concL{\rho_2}"] 
        \ar[dddr, bend right=60, "\concL{\rho_1}"'] 
        \ar[rd, dashed, "\vh"]
    \\ 
        & {\concL{\sigma_1} \pltimes{\anyL{\sigma_1} \liplus{\anyL\sigma} \anyL{\sigma_2}} \concL{\sigma_2}} 
            \ar[rrr, "\pi_2"] 
            \ar[dd, "\pi_1"'] 
            \ar[dr, dashed, "\vh'"]
                &&& {\concL{\sigma_2}}
                    \ar[d, "\concL{\sigma_2}"]
    \\
            && {\anyL{\sigma}} 
                    \ar[rr, "\anyL{\rho_2}"] \ar[d, "\anyL{\rho_1}"'] \pbsquare{drr}
            && {\anyL{\sigma_2}} 
                    \ar[d, "\ext^{\anyL{\rho_2}}\anyL{\rho_1}"] 
    \\
        & {\concL{\sigma_1}} 
            \ar[r, "\concL{\sigma_1}"'] 
            & {\anyL{\sigma_1}} 
                \ar[rr, "\ext^{\anyL{\rho_1}}\anyL{\rho_2}"'] 
                && {\anyL{\sigma_1} \liplus{\anyL\sigma} \anyL{\sigma_2}}
    \end{tikzcd}
    \end{center}

    We first show that $\inv^{\concL{\sigma_2}}\res|{\rho_2}$ factors through $\concL{\sigma_2} \ominus \pi_2$
    \begin{center}
    \begin{tikzcd}[sep=small]
    {\anyL{\sigma_2} \ominus \anyL{\rho_2}} 
        \ar[rr, "\inv^{\concL{\sigma_2}}\res|{\rho_2}"] 
        \ar[dr, dashed]
            && {\concL{\sigma_2} \ominus \concL{\rho_2}} 
                \ar[dr, "\compl{\concL{\rho_2}}"]
    \\
        & {\concL{\sigma_2} \ominus \pi_2} 
            \ar[rr, "\compl{\pi_2}"']
                && {\concL{\sigma_2}}
    \end{tikzcd}
    \end{center}
    and so consider $\va : \vC \in \anyL{\sigma_2} \ominus \anyL{\rho_2}$.
    From the diagram we have that $\va : \vC \in \anyL{\sigma_2} \ominus (\concL{\sigma_2} \comp \proj_2)$
    and consider the canonical isomorphism 
    \[\anyL{\sigma_2} \ominus (\concL{\sigma_2} \comp \proj_2) \iso 
        (\anyL{\sigma_2} \ominus \concL{\sigma_2}) \mplus (\concL{\sigma_2} \ominus \proj_2) \]
    Morphism $\inv^{\concL{\sigma_2}}\res|{\rho_2}$ tells us that $\va$ is in the image of $\concL{\sigma_2}$, and so
    it must come from $\concL{\sigma_2} \ominus \proj_2$.

    We then show that $\ext^{\rho_1}\rho_2$ is variable reflecting by observing that the following diagram commutes
    \begin{center}
    \begin{tikzcd}[sep=small]
        {\concL{\sigma_1 \liplus{\sigma} \sigma_2} \ominus \concL{\ext^{\rho_1}\rho_2}}
            \ar[drrr, "\compl{\concL{\ext^{\rho_1}\rho_2}}"]
    \\
        {\concL{\sigma_2} \ominus \pi_2} 
            \ar[u, "\iso"] \ar[r, "\compl{\pi_2}"]
            & {\concL{\sigma_2}} 
                \ar[rr, "\concL{\ext^{\rho_2}\rho_1}"'] \ar[dd, "\concL{\sigma_2}"']
                && {\concL{\sigma_1 \liplus{\sigma} \sigma_2}}
                        \ar[dd, "\concL{\sigma_1 \liplus\sigma \sigma_2}"]
    \\\\
        {\anyL{\sigma_2} \ominus \anyL{\rho_2}} 
            \ar[r, "\compl{\anyL{\rho_2}}"] \ar[uu, dashed]
            & {\anyL{\sigma_2}} 
                \ar[rr, "\anyL{\ext^{\rho_2}\rho_1}"]
                    && {\anyL{\sigma_1 \liplus\sigma \sigma_2}}
    \\
        {\anyL{\sigma_1 \liplus\sigma \sigma_2} \ominus \anyL{\ext^{\rho_1}\rho_2}} 
            \ar[u, "\iso"] \ar[urrr, "\compl{\anyL{\ext^{\rho_1}\rho_2}}"']
    \end{tikzcd}
    \end{center}\vspace{-1.5\baselineskip}
\end{proof}

\subsubsection{Factorisation System}

\begin{lemma}[Epimorphisms in $\tempC$]\label{lem:epis-T}
   A morphism $\rho : \sigma_1 \to \sigma_2$ in $\tempC$ is epi iff $\anyL\rho$ is surjective.
\end{lemma}
\begin{proof}\nl
    $\to$) Consider $\rho : \sigma_1 \to \sigma_2$ such that for every $\vl\,\vr$, $\vl \comp \rho = \vr \comp \rho \implies \vl = \vr$.
           Take the L.I.C. of $\sigma_2 \xfrom\rho \sigma_1 \xto\rho \sigma_2$, which gives us two morphisms ${\ext^\rho\rho}_\vl, \ext^\rho\rho_\vr : \sigma_2 \to \sigma_2 \liplus{\sigma_1} \sigma_2$
           that satisfy ${\ext^\rho\rho}_l \comp \rho = {\ext^\rho\rho}_r \comp \rho$.
           From the epi property of $\rho$, we have that $\anyL{{\ext^\rho\rho}_l} = \anyL{{\ext^\rho\rho}_r}$ which implies that $\anyL\rho$ is surjective by definition of $\anyL{\sigma_1} \liplus{\anyL\sigma} \anyL{\sigma_2}$.
    \nl
    $\from$) Let $\vl, \vr$ such that $\vl \comp \rho = \vr \comp \rho$.
             This means that $\anyL\vl \comp \anyL\rho = \anyL\vr \comp \anyL\rho$ and so $\anyL\vl = \anyL\vr$ which also implies $\concL\vl = \concL\vr$.
\end{proof}

\begin{notation}
    We denote epis using $\eto$ and regular monos with $\rmto$.
\end{notation}

\begin{lemma}[Orthogonal Factorisation System in $\tempC$]
    The classes of epis and regular monos (variable preserving) morphisms gives rise to an orthogonal factorisation system in $\tempC$.
\end{lemma}
\begin{proof}
The intuition here is that the epi part will represent the instantiations of the location variables coming from its domain
and the regular mono part will be the introduction of all other locations.

A morphism $\rho : \sigma_1 \to \sigma_2$ is factored as follows
\begin{center}
\begin{tikzcd}
    {\anyL{\sigma_1}} \ar[r, Rightarrow]
        & {\anyL{\sigma_1}} \ar[r, "\anyL\rho"] 
            & {\anyL{\sigma_2}}  
    \\
    {\concL{\sigma_1}} \ar[u, "\concL{\sigma_1}"] \ar[r,dashed] \ar[rr, bend right=30, "\concL\rho"'] 
        & {\concL{\sigma_2} \pltimes{\anyL{\sigma_2}} \anyL{\sigma_1}} \ar[u] \ar[r] \pbsquare{ur}
            & {\concL{\sigma_2}} \ar[u, "\concL{\sigma_2}"']
\end{tikzcd}
\end{center}
\end{proof}
\begin{example}
    We factor the morphism $?_\va \xto{\pair{?_{\va,\vb}}{\incl^{\va \inclto \va,\vb,\vc}}} \incl^{\va,\vb \inclto \va,\vb,\vc}$ as
    \[ ?_\va \xto{\pair{?_\va}{1_\va}} 1_\va \xto{\pair{\incl^{\va \inclto \va,\vb}}{\incl^{\va \inclto \va,\vb,\vc}}} \incl^{\va,\vb \inclto \va,\vb,\vc} \]
\end{example}


\subsection[Category of Template Instantiations]{The category of template instantiations $\tinstC$}

We extend category $\tempC$ with instantiation data. 
This is where our template heap functor will be indexed over.
\begin{itemize}
    \item Objects: Sort-preserving injective renamings between worlds $\Minus\vw \injto \Plus\vw$.
    \item Morphisms: Injective, sort-preserving $\concL\me$ and $\anyL\me$ such that
                    \begin{center}
                    \begin{tikzcd}{\Plus\vw} \ar[r, tail, "{\anyL\me}"] & {\Plus{\hat\vw}} \\
                                  {\Minus\vw} \ar[u, tail, "\sigma"] \ar[r, tail, "{\concL\me}"'] & {\Minus{\hat\vw}} \ar[u, tail, "\hat\sigma"']
                     \end{tikzcd}
                    \end{center}
                    and instantiation data $\mvarM_\me : \dheapF(\Minus{\hat\vw} \ominus \concL\me, \Plus{\hat\vw})$.
\end{itemize}
We have a forgetful functor $\tinstU : \tinstC \to \tempC$, and an inclusion functor
$\inclITI\dash : \instC \inclto \tinstC$ which maps to its identity morphism.
We inherit the terminology of $\tempC$ and talk about concrete locations and location variables, 
we overload functors $\concL\dash,\anyL\dash : \tinstC \to \worldC$ by precomposing them with $\tinstU$,
and talk about variable preserving and reflecting morphisms by noticing that $\tinstU$ reflects these properties.

\begin{definition}[Composition of morphisms in $\tinstC$]
Composition of $\sigma_1 \xto{\me_1} \sigma_2 \xto{\me_2} \sigma_3$ is defined as composition of the underlying injections and composition of the instantiation data is given by
\begin{align*}
        \mvarM_{\me_2 \comp \me_1} &: \dheapF(\concL{\sigma_3} \ominus \concL{\me_2 \comp \me_1}, \anyL{\sigma_3}) \\
        \va : \vC & \mapsto \begin{cases}\sem{\phi(\vC)}\anyL{\me_2}(\proj_{\va'}\mvarM_{\me_1}) & \mif \concL{\me_2}(\va') = \va \\
                                         \proj_\va{\mvarM_{\me_2}} & \mif \va \in \concL{\sigma_3} \ominus \concL{\me_2}
                            \end{cases}
\end{align*}
\end{definition}

\begin{remark}
    Every $\sigma, \mvarM \in \dheapF(\concL\sigma, \anyL\sigma)$ induces a morphism $?_\mvarM$
    \begin{center}
        \begin{tikzcd}
            {\emptyset} \ar[rr, "?_{\anyL\sigma}"] && {\anyL\sigma} \\
            {\emptyset} \ar[u, "?_\emptyset"] \ar[rr, "?_{\concL\sigma}"'] && {\concL\sigma} \ar[u, "\concL\sigma"']
        \end{tikzcd}
    \end{center}
\end{remark}

\begin{definition}[Template Heap Functor]
We define the template heap functor $\theapF : \tinstC \to \setC$ as
\begin{align*}
    \theapF(\sigma_1 \xto\me \sigma_2) &: \dheapF(\concL{\sigma_1}, \anyL{\sigma_1}) \to \dheapF(\concL{\sigma_2}, \anyL{\sigma_2}) \\
    \theapF(\me)(\mvarM') &= \mvarM_{\me \comp ?_{\mvarM'}}
\end{align*}
Note that $\mvarM_\me : \theapF(\concL{\sigma_2} \ominus \concL\me \xto{\compl{\concL{\me}}} \concL{\sigma_2} \xto{\concL{\sigma_2}} \anyL{\sigma_2})$ 
and we have that $\theapF \iso \homF\tinstC{?_\emptyset}\dash$.
\end{definition}

\begin{definition}[Variable Extension Functor]
We lift the variable extension functor from $\tempC$ to $\tinstC$
\[\begin{array}{r@{\,}c@{\,}l}
    \extF\dash\vw & : & \tinstC \to \tinstC \\
    \extF\sigma\vw & = & \incl_1 \comp \concL\sigma \\
    \extF\me\vw & = & \Pair{\extF{\tinstU\me}\vw}{\dheapF\pair1{\incl_1}(\mvarM_\me)}
\end{array}\]
\end{definition}

\begin{lemma}[Inhabitation in $\tinstC$]\label{def:inhabited-sorts-tinstC}
    A set of sorts is inhabited (\cref{def:inhabited-sorts}) iff
    for every template $\sigma$, there exists a concrete template $\inclITI{\vw'}$ and a morphism $\sigma \xto\me \inclITI{\vw'}$.
\end{lemma}
\begin{proof}
    For the $\to$ direction, consider $\varL\sigma$ and so for every $\va : \vC \in \varL\sigma$ we have a $\va : \vC \to \vw$.
    By induction on world $\varL\sigma$, this gives us a $\varL\sigma \to \hat\vw$ and a heap $\mvarM : \heapF(\hat\vw)$ by taking the disjoint sum of the worlds and heaps.
    We use this data to construct a morphism from $\sigma$ to $1_{\concL\sigma \oplus \hat\vw}$.

    For the $\from$ direction, we have a morphism $?_{\va : \vC} \xto\me 1_\vw$ for every $\va : \vC$ which is exactly the data needed.
\end{proof}


\begin{lemma}[Monomorphisms in $\tinstC$]\label{lem:monos-tinistC}
    Every morphism in $\tinstC$ is mono.
\end{lemma}
\begin{proof}
    We show $\hat\me : \sigma_2 \to \hat\sigma$ is mono.
    Consider $\hat\me \comp \me_1 = \hat\me \comp \me_2$. 
    We have that $\tinstU\me_1 = \tinstU\me_2$ from \cref{lem:monos-T}.
    It remains to show $\mvarM_{\me_1} = \mvarM_{\me_2}$, so take $\va : \vC \in \concL{\sigma_2} \ominus \concL{\me_1}$.
    \begin{align*}
        \proj_\va\mvarM_{\me_1} & = \proj_\va\mvarM_{\me_2} \iff \\
        \sem{\phi(\vC)}\anyL{\hat\me}(\proj_\va\mvarM_{\me_1}) & = \sem{\phi(\vC)}\anyL{\hat\me}(\proj_\va\mvarM_{\me_2}) \iff \\
        \proj_{\concL{\hat\me}(\va)}(\mvarM_{\me \comp \me_1}) & = \proj_{\concL{\hat\me}(\va)}(\mvarM_{\me \comp \me_2})
    \end{align*}
\end{proof}

\begin{lemma}[Forgetful Functor Preserves Equalizers]\label{lem:tinstU-eq-pres}
    $\tinstU : \tinstC \to \tempC$ preserves equalizers. 
\end{lemma}
\begin{proof}
    Let $\ve : \sigma_1 \to \sigma_2$ be an equalizer of $\me_1$, $\me_2$.
    We show that $\sigma_1$ is canonically isomorphic to $\eqz{\tinstU{\me_1}}{\tinstU{\me_2}}$ when extended with appropriate instantiation data.

    \begin{center}
        \begin{tikzcd}
            {\eqz{\tinstU{\me_1}}{\tinstU{\me_2}}} 
                \ar[rr, "{\eqz{\tinstU{\me_1}}{\tinstU{\me_2}}}"]
                && {\sigma_2} 
                    \ar[r, yshift=3pt, "\tinstU{\me_1}"] \ar[r, yshift=-3pt, "\tinstU{\me_2}"']
                & {\sigma_3}
            \\
            {\sigma_1} 
                \ar[urr, "\tinstU\ve"'] 
                \ar[u, dashed, "\vh"]
        \end{tikzcd}
    \end{center}
    We equip $\vh$ and $\eqz{\tinstU\me_1}{\tinstU\me_2}$ by splitting $\mvarM_\ve$ using the canonical isomorphism
    \[ \concL{\sigma_2} \ominus \concL{\tinstU\me} \iso (\concL{\sigma_2} \ominus \eqz{\concL{\tinstU\me_1}}{\concL{\tinstU\me_2}}) \oplus
                                                        (\eqz{\concL{\tinstU\me_1}}{\concL{\tinstU\me_2}} \ominus \concL\vh) \]
    and we show that the $(\eqz{\concL{\tinstU\me_1}}{\concL{\tinstU\me_2}} \ominus \concL\vh)$ part of data $\mvarM_\ve$ has locations in $\eqz{\anyL{\tinstU\me_1}}{\anyL{\tinstU\me_2}}$.

    So consider $\va : \vC \in \eqz{\concL{\tinstU\me_1}}{\concL{\tinstU\me_2}} \ominus \concL\vh$, and $\pi_{\eqz{\concL{\tinstU\me_1}}{\concL{\tinstU\me_2}}(\va)} \mvarM_\ve$. 
    Notice that
    \begin{align*}
        \mvarM_{\me_1 \comp \ve} &= \mvarM_{\me_2 \comp \ve} \implies  & \tag{def of $\eqz{\concL{\tinstU\me_1}}{\concL{\tinstU\me_2}}$} \\
        \pi_{\concL{\me_1}(\eqz{\concL{\tinstU\me_1}}{\concL{\tinstU\me_2}}(\va))} \mvarM_{\me_1 \comp \ve} & = 
        \pi_{\concL{\me_2}(\eqz{\concL{\tinstU\me_1}}{\concL{\tinstU\me_2}}(\va))} \mvarM_{\me_2 \comp \ve} \implies      & \tag{def of $\mvarM_{\me' \comp \me}$} \\
        \sem{\phi(\vC)}\anyL{\me_1}(\pi_{\eqz{\concL{\tinstU\me_1}}{\concL{\tinstU\me_2}}(\va)} \mvarM_\ve) & = 
        \sem{\phi(\vC)}\anyL{\me_2}(\pi_{\eqz{\concL{\tinstU\me_1}}{\concL{\tinstU\me_2}}(\va)} \mvarM_\ve)
    \end{align*}
    and so for every $\vb : \vC' \in \anyL{\sigma_2}$, we have 
    \[ \vb \isIn \pi_{\eqz{\concL{\tinstU\me_1}}{\concL{\tinstU\me_2}}(\va)} \mvarM_\ve \implies \anyL{\me_1}(\vb) = \anyL{\me_2}(\vb) \]

    We show that $\mvarM_{\me_1 \comp \eqz{\tinstU\me_1}{\tinstU\me_2}} = \mvarM_{\me_2 \comp \eqz{\tinstU\me_1}{\tinstU\me_2}}$
    and so consider $\va : \vC \in \concL{\sigma_3} \ominus \concL{\me_1 \comp \eqz{\tinstU\me_1}{\tinstU\me_2}}$ and proceed by cases
    \begin{itemize}
        \item $\va \in \concL{\sigma_3} \ominus \concL{\tinstU\me_1}$ and $\va \in \concL{\sigma_3} \ominus \concL{\tinstU\me_2}$)
            \begin{align*}
                \proj_\va \mvarM_{\me_1 \comp \eqz{\tinstU\me_1}{\tinstU\me_2}} & = \proj_\va \mvarM_{\me_2 \comp \eqz{\tinstU\me_1}{\tinstU\me_2}} \iff \\
                \proj_\va \mvarM_{\me_1} & = \proj_\va \mvarM_{\me_2} \iff \\
                \proj_\va \mvarM_{\me_1 \comp \ve} & = \proj_\va \mvarM_{\me_2 \comp \ve}
            \end{align*}
        \item $\va \in \concL{\sigma_3} \ominus \concL{\tinstU\me_1}$, $\va' \in \concL{\sigma_2} \ominus \eqz{\concL{\tinstU\me_1}}{\concL{\tinstU\me_2}}$,
              $\va = \concL{\me_2}(\va')$)
            \begin{align*}
                \proj_\va \mvarM_{\me_1 \comp \eqz{\tinstU\me_1}{\tinstU\me_2}} & = \proj_\va \mvarM_{\me_2 \comp \eqz{\tinstU\me_1}{\tinstU\me_2}} \iff \\
                \proj_\va \mvarM_{\me_1} & = \sem{\phi(\vC)}\anyL{\me_2}(\proj_{\va'} \mvarM_{\eqz{\tinstU\me_1}{\tinstU\me_2}}) \iff \\
                \proj_\va \mvarM_{\me_1} & = \sem{\phi(\vC)}\anyL{\me_2}(\proj_{\va'} \mvarM_\ve) \iff \\
                \proj_\va \mvarM_{\me_1 \comp \ve} & = \proj_{\va} \mvarM_{\me_2 \comp \ve}
            \end{align*}
        \item $\va' \in \concL{\sigma_2} \ominus \eqz{\concL{\tinstU\me_1}}{\concL{\tinstU\me_2}}$,
              $\va \in \concL{\sigma_3} \ominus \concL{\tinstU\me_2}$, $\va = \concL{\me_1}(\va')$)
              Symmetric to case above.
        \item $\va_1 \in \concL{\sigma_2} \ominus \eqz{\concL{\tinstU\me_1}}{\concL{\tinstU\me_2}}$,
              $\va_2 \in \concL{\sigma_2} \ominus \eqz{\concL{\tinstU\me_1}}{\concL{\tinstU\me_2}}$,
              $\concL{\me_1}(\va_1) = \va = \concL{\me_2}(\va_2)$)
            \begin{align*}
                \proj_\va \mvarM_{\me_1 \comp \eqz{\tinstU\me_1}{\tinstU\me_2}} & = \proj_\va \mvarM_{\me_2 \comp \eqz{\tinstU\me_1}{\tinstU\me_2}} \iff \\
                \sem{\phi(\vC)}\anyL{\me_1}(\proj_{\va_1} \mvarM_{\eqz{\tinstU\me_1}{\tinstU\me_2}}) & = 
                \sem{\phi(\vC)}\anyL{\me_2}(\proj_{\va_2} \mvarM_{\eqz{\tinstU\me_1}{\tinstU\me_2}}) \iff \\
                \sem{\phi(\vC)}\anyL{\me_1}(\proj_{\va_1} \mvarM_\ve) & = 
                \sem{\phi(\vC)}\anyL{\me_2}(\proj_{\va_2} \mvarM_\ve) \iff \\
                \proj_\va \mvarM_{\me_1 \comp \ve} & = \proj_\va \mvarM_{\me_2 \comp \ve}
            \end{align*}
    \end{itemize}


    As a result, by the UP of $\ve$ we have a $\inv^\vh : \eqz{\tinstU\me_1}{\tinstU\me_2} \to \sigma_1$ in $\tinstC$.
    It remains to show that $\inv^\vh$ is the inverse of $\vh$, that is
    \begin{center}
    \begin{tikzcd}
        {\eqz{\tinstU{\me_1}}{\tinstU{\me_2}}} 
            \ar[rr, "{\eqz{\tinstU{\me_1}}{\tinstU{\me_2}}}"]
            \ar[d, bend left=30, dashed, "\inv^\vh" yshift=5pt]
            && {\sigma_2} 
        \\
        {\sigma_1}
            \ar[urr, "\ve"']
            \ar[u, bend left=30, dashed, "\vh" yshift=5pt, "\iso"' yshift=-5pt]
    \end{tikzcd}
    \end{center}
    which follows from $\ve$ and $\eqz{\tinstU\me_1}{\tinstU\me_2}$ being mono.
\end{proof}
% u does not reflect all equalizers necessarily:
% consider \emptyset -{data: True, True}-> 1_{a1,a2 : Bool} -{id, swap}-> 1_{a1,a2 : Bool}
% in T this is an equalizing diagram but in TI we make a choice on the data to be True, True
% which breaks the UP. We could also have given {data: False, False}.

\begin{lemma}[Forgetful Functor Reflects Pullbacks]\label{lem:tinstU-pls-refl}
    $\tinstU : \tinstC \to \tempC$ reflects pullbacks.
\end{lemma}
\begin{proof}
    Consider a diagram 
        \begin{tikzcd}[sep=small]
            & {\sigma} \ar[dl, "\vp_1"'] \ar[dr, "\vp_2"]
        \\
        {\sigma_1} \ar[dr, "\me_1"'] 
                && {\sigma_2} \ar[dl, "\me_2"]
        \\
            & {\hat\sigma}
        \end{tikzcd}
    s.t. its $\tinstU$-image is a pullback.
    Given $\me_1 : \sigma' \to \sigma_1$, $\me_2 : \sigma' \to \sigma_2$ s.t. $\me_1 \comp \me_1' = \me_2 \comp \me_2'$,
    we get a $\vh : \sigma' \to \sigma$ in $\tempC$.

    We equip $\vh$ with data using the canonical iso
    \[ \concL{\sigma_1} \ominus \concL{\me_1'} \iso (\concL{\sigma_1} \ominus \concL{\vp_1}) \oplus (\concL\sigma \ominus \concL\vh) \]
    by giving $\va : \vC \in \concL\sigma \ominus \concL\vh$, the cell $\proj_{\concL{\vp_1}(\va)} \mvarM_{\me_1'}$.
    Notice that we can factor its locations through $\anyL\sigma$ since
    \begin{align*}
        \sem{\phi(\vC)}\anyL{\me_1}(\proj_{\concL{\vp_1}(\va)} \mvarM_{\me_1'}) & = 
        \proj_{\concL{\me_1 \comp \vp_1}(\va)} \mvarM_{\me_1 \comp \me_1'} \\
         & = \proj_{\concL{\me_2 \comp \vp_2}(\va)} \mvarM_{\me_2 \comp \me_2'} \\
         & = \sem{\phi(\vC)}\anyL{\me_2}(\proj_{\concL{\vp_2}(\va)} \mvarM_{\me_2'})
    \end{align*}
    and so for every $\vb : \vC' \in \anyL{\sigma_1}$, we have 
    \[ \vb \isIn \proj_{\concL{\vp_1}(\va)} \mvarM_{\me_1'} \implies \anyL{\me_1}(\vb) = \anyL{\me_2}(\vb) \]
    and so we can factor its locations through $\inv^{\anyL{\vp_1}}$ since its $\tinstU$-image is a pullback. 

    To show that $\me_1' = \vp_1 \comp \vh$, consider $\va : \vC \in \concL{\sigma_1} \ominus \concL{\vp_1}$.
    Since the $\concL\dash$ image is a pullback, it must be that $\concL{\me_1}(\va) \in \concL{\hat\sigma} \ominus \concL{\me_2}$.
    This implies that 
        \[ \sem{\phi(\vC)}\anyL{\me_1}(\proj_\va \mvarM_{\vp_1}) = \proj_\va \mvarM_{\me_1'}
                                                                 = \sem{\phi(\vC)}\anyL{\me_1}(\proj_\va \mvarM_{\me_1'}) \] 
    which implies $\proj_\va \mvarM_{\vp_1} = \proj_\va \mvarM_{\me_1'}$.
    The case for $\me_2' = \vp_1 \comp \vh$ is symmetric and uniqueness of $\vh$ follows from $\proj_1$ being mono.
\end{proof}
% u does not generally preserve pullbacks.
% Consider cospan 1_{a : B} -{a \mapsto a, False}-> 1_{a,b : Bool} <-{a \mapsto a, True} 1_{a : B} 
% In T, the vertex of the pullback is 1_a
% but in TI the data would not match for vertex 1_a.
% So its pullback becomes 1_a <-{a \mapsto a, True}- ?_a -{a \mapsto a, False}-> 1_a
% This is a counterexample that u preserves pullbacks

\begin{lemma}[Location Functor Preserves Pullbacks]\label{lem:anyL-pls-pres}
    $\anyL\dash : \tinstC \to \worldC$ preserves pullbacks.
\end{lemma}
\begin{proof}
    Consider 
    \begin{tikzcd}
        {\sigma} \ar[r, "\vp_2"] \ar[d, "\vp_1"'] \pbsquare{dr}
            & {\sigma_2} \ar[d, "\me_2"]
        \\
        {\sigma_1} \ar[r, "\me_1"']
            & {\hat\sigma}
    \end{tikzcd}.
    If there are $\va_1 : \vC \in \anyL{\sigma_1}$ and $\va_2 : \vC \in \anyL{\sigma_2}$ s.t. $\anyL{\me_1}(\va_1) = \anyL{\me_2}(\va_2)$,
    we can form a new $\sigma' = \concL{\sigma} \to {\anyL\sigma \oplus \va : \vC}$ that maps $\va$ to $\va_1$ and $\va_2$ and notice that the data of $\vp_1$ and $\vp_2$ will be the same.
    As a result, the new square with root $\sigma'$ will commute and thus it must be that $\va : \vC$ is already in $\sigma$.
\end{proof}

\subsubsection{Reachability}

A concept that becomes relevant when we introduce instantiation data with location variables is the notion of reachable and unreachable locations.
We talk about (un)reachable concrete locations and (un)reachable location variables, where reachability refers to whether we can reach the location by traversing the heap starting from a given concrete location.

\begin{definition}[Reachable Locations]\label{def:reach-loc}
    The reachable locations from $\vw \xto\tau \anyL\sigma$ w.r.t. a template heap $\mvarM : \theapF(\sigma)$
    is the least $\vw^\reach \xto{\tau^\ureach} \anyL\sigma$ such that
    \begin{enumerate}[(i)]
        \item World $\vw$ factors through $\vw^\reach$
            \begin{center}
                \begin{tikzcd} {\vw} \ar[dr, "\exists{\tau^\reach}"'] \ar[r, "\tau"] & {\anyL\sigma} \\ 
                                    & {\vw^\reach} \ar[u, "\tau^\ureach"'] 
             k   \end{tikzcd}
            \end{center}
            \item For every $\vb : \vC \in \vw^\reach \pltimes{\anyL\sigma} \concL\sigma$ \ie every reachable concrete location, the locations in cell $\vd = \proj_{\res_{\concL\sigma}\tau^\ureach(\vb)} \mvarM$ are contained in $\vw^\reach$
            \begin{center}
                \begin{tikzcd}
                    {1} \ar[r, "\vd"] \ar[dr, "\exists \vd^\reach"'] &  {\sem{\phi(\vC)}{\anyL\sigma}} \\
                                                                  & {\sem{\phi(\vC)}\vw^\reach} \ar[u, "\sem{\phi(\vC)}{\tau^\ureach}"']
                \end{tikzcd}
            \end{center}
    \end{enumerate}
    Least here means that for any other compatible $\tau' : \vw' \to \anyL\sigma$,
    there exists a unique $\vh : \vw^\reach \to \vw'$ that satisfies $\tau' \comp \vh = \tau^\ureach$.

    We refer to $\vw^\reach$ as the reachable locations and $\anyL\sigma \ominus \tau^\ureach$ as the unreachable ones.
\end{definition}

\begin{example}\label{ex:reach-loc}
    Consider $\tau = \incl^{\va \inclto \va,\vb,\vc}$ and template heap over $1_{\va,\vb,\vc}$ being: 
    \begin{center}
    \renewcommand\ndsize{15}
    \begin{tikzpicture}
        \node[ccell=$\vc$] (c) {};
        \node[ccell=$\va$] (a) [below=0pt of c] {};
        \node[ccell=$\vb$] (b) [below=0pt of a] {};
        \draw[ptr] (c.center) to ++\east to ++\south to (a.east);
        \draw[ptr] ($ (a.center) + (0,-\gap pt) $) to ++\east to ++\south to ($ (b.east) + (0,-\gap pt) $);
    \end{tikzpicture}
    \end{center}
    The least $\tau^\ureach$ is $\incl^{\va,\vb \inclto \va,\vb,\vc}$ since $\vc$ is not reachable from $\va$.
    Notice that $\incl^{\va \inclto \va,\vb,\vc}$ is not the least as it does not satisfy condition (ii) for $\va$.
\end{example}

\paragraph{Reachability in terms of Heap Walks}

We give an alternative, algorithmic description of reachability (\cref{def:reach-loc}) in terms of \emph{heap walks} which describe finite, sequential traversals of the heap.
\begin{definition}[Heap Walks]
    For a heap $\mvarM \in \theapF(\sigma)$, a sequence 
    \[ \va_0 : \vC_0, \hdots, \va_\vn : \vC_\vn, \vc : \vC \in \concL\sigma \times \klstar{\concL\sigma} \times \anyL\sigma \] 
    is called a \emph{heap walk} if
    \[ \All{\vi < \vn} \concL\sigma(\va_{\vi+1}) \isIn \proj_{\va_\vi} \mvarM  \Land \vc \isIn \proj_{\va_\vn} \mvarM \]
    and it contains \emph{no duplicate locations}
    \[ \All{\vi, \vj \le \vn} \vi \neq \vj \implies \va_\vi \neq \va_\vj \Land \All{\vi \le \vn} \concL\sigma(\va_\vi) \neq \vc \]
\end{definition}
We informally use the proposition $(\va \isIn \vd)$ to state that location $\va$ is accessed inside cell $\vd$.
For brevity, we avoid defining this directly but the intuition should be clear.
\begin{example}
    Consider the heap from \cref{ex:reach-loc}. Its heap walks are 
    (i) $\vc, \va, \vb$, 
    (ii) $\vc, \va$, 
    (iii) $\va, \vb$
\end{example}

\begin{definition}[Reachable Location via a Heap Walk]
    A location $\vc : \vC \in \anyL\sigma$ is said to be reachable from $\va : \vC' \in \anyL\sigma$ in heap $\mvarM : \theapF(\sigma)$ if 
    $\va = \vc$ or there exists a heap walk $\va' : \vC', \vec\vb, \vc : \vC$ via $\mvarM$ with $\concL\sigma(\va') = \va$.
\end{definition}

\begin{definition}[Reachable Locations via Heap Walks]\label{def:reach-loc-hwalks}
    The reachable locations from $\vw \xto\tau \anyL\sigma$ w.r.t. a template heap $\mvarM : \theapF(\sigma)$
    is a $\vw^\reach \xto{\tau^\ureach} \anyL\sigma$ that satisfies
    \begin{enumerate}[(i)]
        \item World $\vw$ factors through $\vw^\reach$
            \begin{center}
                \begin{tikzcd} {\vw} \ar[dr, "\exists{\tau^\reach}"'] \ar[r, "\tau"] & {\anyL\sigma} \\ 
                                    & {\vw^\reach} \ar[u, "\tau^\ureach"'] 
                \end{tikzcd}
            \end{center}
        \item For every $\vc : \vC \in \anyL\sigma$, $\vc \in \img{\tau^\ureach} \iff \vc \text{ is reachable from a } \tau(\va)$
        \end{enumerate}
\end{definition}

\begin{lemma}[Reachable Locations Property Equivalence]\label{lem:reach-loc-prop-equiv}
    The two definitions of reachability, \cref{def:reach-loc} and \cref{def:reach-loc-hwalks}, are equivalent.
\end{lemma}
\begin{proof}
    Property (i) of both definitions is identical, so we only show that property (ii) in the two definitions is equivalent.

    \cref{def:reach-loc} $\implies$ \cref{def:reach-loc-hwalks}) 

    Direction $\To$) 
    Let $\vc : \vC \in \vw^\reach$.
    If $\tau^\ureach(\vc) \in \img\tau$ we are done as then $\vc$ is reachable from itself. 
    If $\tau^\ureach(\vc) \not\in \img\tau$, we enumerate all heap walks $\va_0 : \vC_0, \hdots, \va_\vn : \vC_\vn, \tau^\ureach(\vc) : \vC$ of at most $\card{\vw^\reach}$ length, 
    that further satisfy $\All\vi \sigma(\va_\vi) \in \img{\tau^\ureach}$.
    There must exist at least one such heap walk as then $\vw^\reach \setDiff \{\vc\}$ would be the least world.
    If any of these heap walks has $\sigma(\va_0) \in \img\tau$ we are done.  
    Otherwise collect all $\va_\vi$ in all the enumerated heap walks and remove them from $\vw^\reach$, call this $\hat\vw^\reach$.
    We show that $\hat\vw^\reach$ would be the least world which would be a contradiction.
    Take a $\vb \in \hat\vw^\reach \pltimes{\anyL\sigma} \concL\sigma$, and consider the locations in $\vd_\vb$.
    If $\vd_\vb$ contains a location $\vb'$ which is not in $\hat\vw^\reach$ it would mean that $\vb'$ is part of a heap walk generated above. 
    Since the heap walks do not contain duplicates it must be less than $\card{\vw^\reach}$ and thus we can extend that heap walk with $\vb$.
    So that would be a contradiction that $\vb \in \hat\vw^\reach$.
    And so $\hat\vw^\reach$ is the least world which is itself a contradiction and so there must exist at least one heap walk with $\concL\sigma(\va_0) \in \img\tau$.

    Direction $\From$) We have that $\vc : \vC$ is reachable from some $\tau(\va)$. If $\va = \vc$ we are done.
    Otherwise we have a heap walk $\va_0 : \vC_0, \hdots, \va_\vn : \vC_\vn, \vc : \vC$ with $\concL\sigma(\va_0) \in \img\tau$.
    Proceed by induction on $\vn$. If $\vn = 0$, then $\concL\sigma(\va_0) \in \img\tau$, and so it is also in $\img{\tau^\ureach}$.
    If $\va_\vn \in \vw^\reach$ and $\va_{\vn+1} \in \vd_{\va_\vn}$, then $\va_{\vn + 1}$ must also be in $\vw^\reach$ by the  
    fact that $\vd_{\va_\vn}$ factors through $\vw^\reach$.
    So every $\va_\vi$ is in $\vw^\reach$, and thus $\vc$ must also be in $\vw^\reach$.
    \nl
    \cref{def:reach-loc-hwalks} $\implies$ \cref{def:reach-loc}) 
    Consider a $\va : \vC \in \vw^\reach \pltimes{\anyL\sigma} \concL\sigma$, and $\vd_\va \in \sem{\phi(\vC)}\anyL\sigma$.
    Since $\va$ is in $\vw^\reach$, $\va$ is reachable from a $\tau(\va')$ by assumption, and any location in $\vd_\va$ will also be reachable 
    by extending the heap walk with that new location.
    As a result, by the bi-implication, the new location will also be in $\vw^\reach$.
\end{proof}

\begin{lemma}[Decidability of Reachability]\label{lem:reach-loc-dec}
    There exists an algorithm that computes $\vw^\reach \xto{\tau^\ureach} \anyL\sigma$ for every $\vw \xto\tau \anyL\sigma$.
\end{lemma}
\begin{proof} 
    Intuitively, the algorithm will traverse the heap from $\vw \pltimes{\anyL\sigma} \concL\sigma$ and return all the locations found including $\img\tau$. 
\end{proof}

\subsubsection{Connectivity}

We extend the notion of reachability to morphisms in $\tinstC$ to characterise morphisms that satisfy 
(i) all cells in the instantiation data are reachable from locations in its domain, 
or (ii) none of the cells in the instantiation data are reachable from locations in its domain. 
We appropriately call morphisms that satisfy (i) as connected, and (ii) as disconnected.

We first give a factorisation of morphisms in $\tinstC$, by factoring out the reachable and unreachable components of a morphism.

\begin{definition}[Factorisation in $\tinstC$]\label{def:con-fact}
    For a morphism $\me : \sigma_1 \to \sigma_2$, we consider the reachable locations from $\anyL\me$ w.r.t. $\mvarM_\me : \theapF(\concL{\sigma_2} \ominus \concL\me \xto{\compl{\concL\me}} \concL{\sigma_2} \xto{\concL{\sigma_2}} \anyL{\sigma_2})$.
    By \cref{lem:reach-loc-dec} we get 
    \[ \anyL{\sigma_1} \xto{{\anyL\me}^\reach} \vw^\reach \xto{{\anyL\me}^\ureach} \anyL{\sigma_2} \]

    We extend ${\anyL\me}^\reach$ and ${\anyL\me}^\ureach$ into morphisms in $\tinstC$ by taking the \emph{largest possible choice} for the concrete locations,
    that is, we consider the pullback
    \begin{center}
    \begin{tikzcd}
        {\anyL{\sigma_1}} 
            \ar[r, "{\anyL\me}^\reach"] \ar[rrr, bend left=40, "\anyL\me"] 
            & {\vw^\reach} 
                \ar[rr, "{\anyL\me}^\ureach"] 
                && {\anyL{\sigma_2}} 
        \\
        {\concL{\sigma_1}} 
            \ar[u, "\concL{\sigma_1}"] \ar[r, dashed, "\vh"] \ar[rrr, bend right=40, "\concL\me"] 
            & {\vw^\reach \pltimes{\anyL{\sigma_2}} \concL{\sigma_2}} 
                \ar[u] \ar[rr] \pbsquare{urr} 
                && {\concL{\sigma_2}} 
                    \ar[u, "\concL{\sigma_2}"']
    \end{tikzcd}
    \end{center}
    where the instantiation data are given by the canonical isomorphism
    \[ \concL{\sigma_2} \ominus \concL\me \iso (\concL{\sigma_2} \ominus \res_{\concL{\sigma_2}}{\anyL\me}^\ureach) 
                                          \mplus (\vw^\reach \pltimes{\anyL{\sigma_2}} \concL{\sigma_2} \ominus \vh) \]
    
    We show that locations $(\vw^\reach \pltimes{\anyL{\sigma_2}} \concL{\sigma_2}) \ominus \vh$ have cells that access only $\vw^\reach$ from
    \[ (\vw^\reach \pltimes{\anyL{\sigma_2}} \concL{\sigma_2}) \ominus \vh \iso \vw^\reach \pltimes{\anyL{\sigma_2}} (\concL{\sigma_2} \ominus \concL\me) \]
    which follows from pullback pasting
    
    \begin{center}
    \begin{tikzcd}
        {\vw^\reach} 
            \ar[r, "{\anyL\me}^\ureach"] 
            & {\anyL{\sigma_2}} 
        \\
        {\vw^\reach \pltimes{\anyL{\sigma_2}} \concL{\sigma_2}} \pbsquare{ur}
            \ar[u, "\res_{{\anyL\me}^\ureach}\concL{\sigma_2}"]
            \ar[r, "\res_{\concL{\sigma_2}}{\anyL\me}^\ureach"]
            & {\concL{\sigma_2}}
                \ar[u, "\concL{\sigma_2}"']
        \\
        {(\vw^\reach \pltimes{\anyL{\sigma_2}} \concL{\sigma_2}) \ominus \vh} \pbsquare{ur}
            \ar[u, "\compl\vh"]
            \ar[r, "\res_{\concL{\sigma_2}}{\anyL\me}^\ureach"' yshift=-2pt]
            & {\concL{\sigma_2} \ominus \concL\me}
                \ar[u, "\compl{\concL\me}"']
    \end{tikzcd}
    \end{center}

    and so \cref{def:reach-loc} gives us the factoring of the cells through $\vw^\reach$.
\end{definition}

\begin{notation}
    For a morphism $\me : \sigma_1 \to \sigma_2$, we write $\con\me : \sigma_1 \to \sigma_{\con\me}$ and $\uncon\me : \sigma_{\con\me} \to \sigma_2$ for its factorisation given by \cref{def:con-fact}.
    We also write $\sigma_{\uncon\me} \eqdef \sigma_2 \ominus \uncon\me$.
\end{notation}

\begin{example}
    \renewcommand\ndsize{15}

    Consider the morphism $?_\va \xto{\pair{?_{\{\va,\vb\}}}{\incl^{\va \inclto \va,\vb,\vc}}} \incl^{\va,\vb \inclto \va,\vb,\vc}$ with data
    \begin{tikzpicture}[baseline={(a.base)}]
        \node[vcell=$\vc$] (c) {};
        \node[ccell=$\va$] (a) [below=0pt of c] {};
        \node[ccell=$\vb$] (b) [below=0pt of a] {};
        \draw[ptr] (c.center) to ++\east to ++\south to (a.east);
        \draw[ptr] ($ (a.center) + (0,-\gap pt) $) to ++\east to ++\south to ($ (b.east) + (0,-\gap pt) $);
    \end{tikzpicture}

    Its reachable/unreachable factoring given by \cref{def:con-fact} is
    \[ ?_\va \xto{\pair{?_{\{\va,\vb\}}}{\incl^{\va \inclto \va,\vb}}} 1_{\va,\vb} \xto{\pair{\incl^{\va,\vb \inclto \va,\vb,\vc}}{\incl^{\va,\vb \inclto \va,\vb,\vc}}} \incl^{\va,\vb \inclto \va,\vb,\vc} \]
    with data
    \begin{tikzpicture}
        \node[ccell=$\va$] (a) {};
        \node[ccell=$\vb$] (b) [below=0pt of a] {};
        \draw[ptr] ($ (a.center) + (0,-\gap pt) $) to ++\east to ++\south to ($ (b.east) + (0,-\gap pt) $);
    \end{tikzpicture} for the reachable morphism.
\end{example}

\begin{definition}[Connected morphisms]\label{def:connected-morph}
    A morphism $\me : \sigma_1 \to \sigma_2$ is said to be connected if $\uncon\me$ is an isomorphism.
\end{definition}
\begin{definition}[Disconnected morphisms]
    A morphism $\me : \sigma_1 \to \sigma_2$ is said to be disconnected if $\con\me$ is an isomorphism.
\end{definition}

\begin{lemma}[Factoring Connectedness]\label{lem:con-fact-is-con}
    For a morphism $\me : \sigma_1 \to \sigma_2$, its $\con\me$ part is connected.
\end{lemma}
\begin{proof}
    We have the diagram
    \begin{tikzcd}
        {\sigma_1} \ar[rr, "\con\me"] \ar[dr, "\con{\con\me}"'] 
                && {\sigma_{\con\me}} \ar[r, "\uncon\me"] 
                    & {\sigma_2}
        \\
                & {\sigma_{\con{\con\me}}} \ar[ur, "\uncon{\con\me}"']
    \end{tikzcd} and we want to show that $\uncon{\con\me}$ is iso.

    It is enough to show that $\anyL{\uncon{\con\me}}$ is iso, as $\concL{\uncon{\con\me}}$ is given by the pullback
    \begin{center}
    \begin{tikzcd}
        {\anyL{\sigma_{\con{\con\me}}}} 
            \ar[r, "\anyL{\uncon{\con\me}}"] 
            & {\anyL{\sigma_{\con\me}}}
        \\
        {\concL{\sigma_{\con{\con\me}}}} \pbsquare{ur}
            \ar[u] \ar[r, "\concL{\uncon{\con\me}}"']
            & {\concL{\sigma_{\con\me}}} 
                \ar[u, "\concL{\sigma_{\con\me}}"']
    \end{tikzcd}
    \end{center}
    and so if $\anyL{\uncon{\con\me}}$ is iso, so is $\concL{\uncon{\con\me}}$.

    We proceed by showing that 
    \[ \anyL{\sigma_1} \xto{\anyL{\con{\con\me}}} \anyL{\sigma_{\con{\con\me}}} (\xto{\anyL{\uncon{\con\me}}} \anyL{\sigma_{\con\me}} \xto{\anyL{\uncon\me}} \anyL{\sigma_2}) \] 
    satisfies \cref{def:reach-loc} for $\anyL{\sigma_1} \xto{\anyL\me} \anyL{\sigma_2}$ and $\mvarM_\me$.
    As a result, by the least property of \cref{def:reach-loc} for $\anyL{\con\me}$, we get that $\anyL{\uncon{\con\me}}$ must be iso.

    Property (i) is immediate, and so consider property (ii).
    Consider $\va : \vC \in (\concL{\sigma_2} \ominus \concL\me) \pltimes{\anyL{\sigma_2}} \anyL{\sigma_{\con{\con\me}}}$,
    where the pullback is given by
    \begin{center}
        \begin{tikzcd}
        {\anyL{\sigma_{\con{\con\me}}}}
            \ar[r, "\anyL{\uncon{\con\me}}"] 
            & {\anyL{\sigma_{\con\me}}} 
                \ar[r, "\anyL{\uncon\me}"] 
                & {\anyL{\sigma_2}}
        \\
        {\concL{\sigma_{\con{\con\me}}}} \pbsquare{ur}
            \ar[u]
            \ar[r, "\concL{\uncon{\con\me}}"'] 
            & {\concL{\sigma_{\con\me}}} \pbsquare{ur}
                \ar[u]
                \ar[r, "\concL{\uncon\me}"'] 
                & {\concL{\sigma_2}}
                    \ar[u, "\concL{\sigma_2}"']
        \\
        {\concL{\sigma_{\con{\con\me}}} \ominus \concL{\con{\con\me}}} \pbsquare{ur}
            \ar[u, "\compl{\concL{\con{\con\me}}}"]
            \ar[r, "\concL{\uncon{\con\me}}"]
                & {\concL{\sigma_{\con\me}} \ominus \concL{\con\me}} \pbsquare{ur}
                    \ar[u, "\compl{\concL{\con\me}}"]
                    \ar[r, "\concL{\uncon\me}"]
                & {\concL{\sigma_2} \ominus \concL\me}
                    \ar[u, "\compl{\concL\me}"]
        \end{tikzcd}
    \end{center}
    which we use to prove that cell $\va$ has locations in $\anyL{\sigma_{\con\con\me}}$
    \begin{align*}
        \proj_{\concL{\uncon\me}\concL{\uncon{\con\me}}(\va)} \mvarM_\me
            & = \sem{\phi(\vC)}\anyL{\uncon\me}(\proj_{\concL{\uncon{\con\me}}(\va)} \mvarM_{\con\me}) \\
            & = \sem{\phi(\vC)}\anyL{\uncon\me}(\sem{\phi(\vC)}\anyL{\uncon{\con\me}}(\proj_{\va} \mvarM_{\con{\con\me}}))
    \end{align*}
\end{proof}

\begin{lemma}[Factoring and Variable Preservation]\label{lem:con-fact-is-vpres}
    For a morphism $\me : \sigma_1 \to \sigma_2$, its $\uncon\me$ part is variable preserving.
\end{lemma}
\begin{proof}
    Morphism $\uncon\me$ is given by the following diagram
    \begin{center}
    \begin{tikzcd}
        {\anyL{\sigma_1}} \ar[r, "\anyL{\con\me}"] \ar[rr, bend left=40, "\anyL\me"] & {\anyL{\sigma_\me}} \ar[r, "\anyL{\uncon\me}"] & {\anyL{\sigma_2}} \\
        {\concL{\sigma_1}} \ar[u, "\concL{\sigma_1}"] \ar[r, dashed, "\concL{\con\me}"'] \ar[rr, bend right=40, "\concL\me"] & {\concL{\sigma_\me}} \ar[u, "\concL{\sigma_\me}"] \ar[r, "\concL{\uncon\me}"'] \pbsquare{ur} & {\concL{\sigma_2}} \ar[u, "\concL{\sigma_2}"']
    \end{tikzcd}
    \end{center}
    which by \cref{lem:vpres-prop} since $\uncon\me$ is a pullback it is also variable preserving.
\end{proof}

\begin{lemma}[Epimorphisms in $\tinstC$]\label{lem:epis-tinstC}
   A morphism $\me : \sigma_1 \to \sigma_2$ in $\tinstC$ is epi iff $\me$ is connected.
\end{lemma}
\begin{proof}\nl
    $\to$)  
    From \cref{def:con-fact} we can factor $\me$ into $\sigma_1 \xto{\con\me} \sigma_\me \xto{\uncon\me} \sigma_2$ and
    to show that $\me$ is connected it is enough to show that $\anyL{\uncon\me}$ is surjective.

    By \cref{lem:con-fact-is-vpres} we form morphisms $1 \mplus \incl_1$ and $1 \mplus \incl_2$ in $\tinstC$
    \begin{center}
    \begin{tikzcd}
            & {\anyL{\sigma_\me} \mplus (\anyL{\sigma_2} \ominus \anyL{\uncon\me})} 
                \ar[rr, shift left=1, "1 \mplus \incl_1"] \ar[rr, shift right=1, "1 \mplus \incl_2"'] 
                && {\anyL{\sigma_\me} \mplus \Exp{(\anyL{\sigma_2} \ominus \anyL{\uncon\me})}2}
            \\
        {\sigma_2} 
            \ar[r, phantom, "\iso_{\mathsf{r}}" xshift=-28pt] 
            & {}
            \\
            & {\concL{\sigma_\me} \mplus (\concL{\sigma_2} \ominus \concL{\uncon\me})} 
                \ar[uu, "\concL{\sigma_\me} \mplus \concL{\sigma_2}\res|{\uncon\me}"] 
                \ar[rr, shift left=1, "1 \mplus \incl_1"] 
                \ar[rr, shift right=1, "1 \mplus \incl_2"']    
                && {\concL{\sigma_\me} \mplus \Exp{(\concL{\sigma_2} \ominus \concL{\uncon\me})}2}
                    \ar[uu, "\concL{\sigma_\me} \mplus \Exp{\concL{\sigma_2}\res|{\uncon\me}}2"']
    \end{tikzcd}
    \end{center}
    with data $\mvarM_{1 \mplus \incl_\vi} = \dheapF(\iso, 1 \oplus \incl_{\vi \mod 2 + 1} \comp \anyL{\iso_{\mathsf{r}}})(\mvarM_{\uncon\me})$.
    The data $\mvarM_{(1 \oplus \incl_\vi) \comp \uncon\me}$ results in two copies of $\mvarM_{\uncon\me}$ and so
    $\mvarM_{(1 \oplus \incl_1) \comp \uncon\me} = \mvarM_{(1 \oplus \incl_2) \comp \uncon\me}$. 
    As a result 
    \begin{align*}
        (1 \oplus \incl_1) \comp \uncon\me & = (1 \oplus \incl_2) \comp \uncon\me \implies \\
        (1 \oplus \incl_1) \comp \me        & = (1 \oplus \incl_2) \comp \me \implies \\
        1 \oplus \incl_1                    & = 1 \oplus \incl_2 \implies \\
        \anyL{\uncon\me}                   &\text{ is surjective} \implies \\
        \uncon\me                          &\text{ is iso} \\
    \end{align*}

    $\from$) Consider $\me_1 \comp \me = \me_2 \comp \me$.
             We show that $\anyL{\me_1} = \anyL{\me_2}$ which in turn implies $\me_1 = \me_2$.

             Consider $\va' : \vC \in \anyL{\sigma_2}$, which by connectedness gives us a $\va : \vC \in \anyL{\sigma_\me}$ s.t. $\anyL{\uncon\me}(\va) = \va'$.
             If $\va$ is in the image of $\anyL{\con\me}$ we are done.
             Otherwise by \cref{lem:reach-loc-prop-equiv} we have a heap walk $\vb_0 : \vC_0, \hdots, \vb_\vn : \vC_\vn, \va : \vC$
             along $\mvarM_{\sigma_\me}$ where $\concL{\sigma_\me}(\vb_0) \in \img{\anyL{\con\me}}$.

             We proceed by induction on the length of the sequence to show that 
             \[ \All{\vi \le \vn} \anyL{\me_1}\concL{\sigma_\me}(\vb_\vi) = \anyL{\me_2}\concL{\sigma_\me}(\vb_\vi) \]
             $\vn = 0$) Then $\vb_0$ is in the image of $\anyL{\con\me}$ and so $\anyL{\me_1}\concL{\sigma_\me}(\vb_0) = \anyL{\me_2}\concL{\sigma_\me}(\vb_0)$
                        by $\me \comp \me_1 = \me \comp \me_2$.\nl
             $\vn = \vn' + 1$) 
                        We have
                        \begin{align*}
                            \anyL{\me_1}\concL{\sigma_\me}(\vb_{\vn'}) &= \anyL{\me_2}\concL{\sigma_\me}(\vb_{\vn'}) = \tag{$\mathsf{IH}$} \\
                            \concL{\sigma_2}\concL{\me_1}(\vb_{\vn'})  &= \concL{\sigma_2}\concL{\me_2}(\vb_{\vn'}) \implies \tag{by $\concL{\sigma_2}$ is injective} \\
                            \concL{\me_1}(\vb_{\vn'})                  &= \concL{\me_2}(\vb_{\vn'}) \implies        \tag{by $\mvarM_{\me_1 \comp \me} = \mvarM_{\me_2 \comp \me}$} \\
                            \sem{\phi(\vC_{\vn'})}\anyL{\me_1}(\vd_{\vn'})  &= \sem{\phi(\vC_{\vn'})}\anyL{\me_2}(\vd_{\vn'}) \implies  \tag{and since $\vb_\vn \isIn \vd_{\vn'}$} \\
                            \anyL{\me_1}\concL{\sigma_\me}(\vb_\vn)    &= \anyL{\me_2}\concL{\sigma_\me}(\vb_\vn)                       
                        \end{align*}
             Since $\va \text{ is in } \vd_\vn$ and $\anyL{\me_1}\concL{\sigma_\me}(\vb_\vn) = \anyL{\me_2}\concL{\sigma_\me}(\vb_\vn)$,
             we have that $\anyL{\me_1}(\va) = \anyL{\me_2}(\va)$.
\end{proof}

\begin{lemma}[Composition of Connected Morphisms]
    Connected morphisms are closed under composition.
\end{lemma}
\begin{proof}
    Follows immediately from \cref{lem:epis-tinstC} and epis being closed under composition.
\end{proof}

\begin{lemma}[Regular Monos in $\tinstC$]\label{lem:rmono-dcon-tinstC}
    A morphism $\me : \sigma_1 \to \sigma_2$ in $\tinstC$ is regular mono iff $\me$ is variable preserving.
\end{lemma}
\begin{proof}\nl
    $\to$) If $\me$ is part of an equalizer diagram, by \cref{lem:tinstU-eq-pres}, if we forget the data we still have an equalizer diagram in $\tempC$.
           By \cref{lem:rmono-vpres-tempC}, $\tinstU\me$ will be variable preserving and since variable preservation is reflected by $\tinstU$, $\me$ is variable preserving.\nl
    \def\ciso{\iso^{\mathsf{c}}}
    \def\isoO{\iso^1}
    \def\isoT{\iso^2}
    $\from$) 
    We show that this is an equalizer diagram (\cref{lem:vpres-prop} gives us $\concL{\sigma_2}\res|\me$, and $\ciso$ is the canonical isomorphism induced)
    \begin{center}
    \begin{tikzcd}
            & {\anyL{\sigma_1} \mplus (\anyL{\sigma_2} \ominus \anyL\me)} 
                \ar[rr, shift left=1, "1 \mplus \incl_1"] 
                \ar[rr, shift right=1, "1 \mplus \incl_2"'] 
                    && {\anyL{\sigma_1} \mplus (\anyL{\sigma_2} \ominus \anyL\me) \mplus (\anyL{\sigma_2} \ominus \anyL\me)}
\\
        {\sigma_2} \ar[r, phantom, pos=0.20, "\ciso"] & {}
\\
            & {\concL{\sigma_1} \mplus (\concL{\sigma_2} \ominus \concL\me)} 
                \ar[uu, "\concL{\sigma_1} \mplus \concL{\sigma_2}\res|\me"'] 
                \ar[rr, shift left=1, "1 \mplus \incl_1"] 
                \ar[rr, shift right=1, "1 \mplus \incl_2"']    
                    && {\concL{\sigma_1} \mplus (\concL{\sigma_2} \ominus \concL\me) \mplus (\concL{\sigma_2} \ominus \concL\me)} 
                        \ar[uu, "\concL{\sigma_1} \mplus \concL{\sigma_2}\res|\me \mplus \concL{\sigma_2}\res|\me"']
    \end{tikzcd}
    \end{center}
    We equip $1 \oplus \incl_1$ with instantiation data $\dheapF(\isoO, 1 \oplus \incl_2 \comp \anyL\ciso)(\mvarM_\me)$ by noticing the isomorphism
    \[ (\concL{\sigma_1} \mplus (\concL{\sigma_2} \ominus \concL\me) \mplus (\concL{\sigma_2} \ominus \concL\me)) \ominus (1 \oplus \incl_1)
        \isoO \concL{\sigma_2} \ominus \concL\me \]
    We do the same for $1 \oplus \incl_2$, \ie we use $\dheapF(\isoT, 1 \oplus \incl_1 \comp \anyL\ciso)(\mvarM_\me)$.

    To show that $1 \oplus \incl_1 \comp \ciso \comp \me = 1 \oplus \incl_2 \comp \ciso \comp \me$, we need to show that the data match.
    Notice that 
    \[ (\concL{\sigma_1} \mplus \Exp{(\concL{\sigma_2} \ominus \concL\me)}2) \ominus (1 \oplus \incl_1 \comp \concL\me) \iso \Exp{(\concL{\sigma_2} \ominus \concL\me)}2 \]
    and $\mvarM_{1 \oplus \incl_1 \comp \ciso \me} = \mvarM_{1 \oplus \incl_2 \comp \ciso \me}$ as they have two identical copies of $\dheapF(1, \ciso)(\mvarM_\me)$. 

    Now consider $(1 \mplus \incl_1) \comp \ciso \comp \me' = (1 \mplus \incl_2) \comp \ciso \comp \me'$. 
    Then $\tinstU\me'$ must factor through $\tinstU\me$ and $\mvarM_{\me'}$ can be split between the factoring through the canonical isomorphism.
    Uniqueness of this factoring follows from $\me$ being mono.
\end{proof}

\begin{lemma}[Regular Monos are closed under composition in $\tinstC$]\label{lem:rmono-comp}
    Regular monos in $\tinstC$ are closed under composition.
\end{lemma}
\begin{proof}
    Follows from \cref{lem:rmono-dcon-tinstC} and \cref{lem:vpres-dcon-equiv}, $\vu$ preserving and reflecting variable preserving morphisms and \cref{lem:comp-vpres}.
\end{proof}

\begin{lemma}[Variable Preservation and Disconnectedness]\label{lem:vpres-dcon-equiv}
    A morphism $\me : \sigma_1 \to \sigma_2$ is variable preserving iff it is disconnected.
\end{lemma}
\begin{proof}
    The $\from$ direction follows from \cref{lem:con-fact-is-vpres} and since $\con\me$ is iso it is also variable preserving, and variable preservation is closed under composition.
    For the $\to$ direction, by \cref{lem:con-fact-is-vpres} and \cref{lem:vpres-cancl} we have that $\con\me$ is variable preserving.
    Since $\con\me$ is epi and regular mono it is also iso.
\end{proof}

\begin{lemma}[Factoring Disconnectedness]\label{lem:con-fact-is-dcon}
    For a morphism $\me : \sigma_1 \to \sigma_2$, its $\uncon\me$ part is disconnected.
\end{lemma}
\begin{proof}
    By \cref{lem:con-fact-is-vpres}, $\uncon\me$ is variable preserving which by \cref{lem:vpres-dcon-equiv} means it disconnected.
\end{proof}

\subsubsection{Factorisation System}\label{sec:fact-tinstC}

\begin{lemma}[Factorisation is Unique]\label{lem:fact-uniq-tinstC}
    The connectivity (epi-regular mono) factorisation in $\tinstC$ is unique up to unique isomorphism.
\end{lemma}
\begin{proof}
    We show that the $\sigma_1 \xto{\con\me} \sigma_{\con\me} \xto{\uncon\me} \sigma_2$ factorisation is unique up to unique iso, 
    so consider 
    \begin{center}
        \begin{tikzcd}
            {\sigma_1} 
                \ar[r, twoheadrightarrow,  "\con\me"] \ar[rr, twoheadrightarrow, bend left=35, "\me"] \ar[drr, twoheadrightarrow, bend right=70,  "\me_1"'] \ar[dr, "\con{\me_1}"' yshift=2pt]
                & {\sigma_{\con\me}}
                    \ar[-rmonotip, r, "\uncon\me"] \ar[d, dashed, "\vh"', "\iso"]
                    & {\sigma_2}
            \\
                & {\sigma_{\con{\me_1}}}
                    \ar[-rmonotip, r, "\uncon{\me_1}"', "\iso"]
                    & {\sigma'}
                        \ar[-rmonotip, u, "\me_2"']
        \end{tikzcd}
    \end{center}
    and we show that there exists a unique iso $\vh : \sigma_{\con\me} \to \sigma_{\con{\me_1}}$ that commutes in the diagram above.

    We first show that 
    \[ \anyL{\sigma_1} \xto{\anyL{\con{\me_1}}} \anyL{\sigma_{\con{\me_1}}} (\xto{\uncon{\me_1}} \sigma' \xto{\anyL{\me_2}} \sigma_2) \]
    satisfies \cref{def:reach-loc} for $\anyL{\sigma_1} \xto{\anyL\me} \anyL{\sigma_2}$ and $\mvarM_\me$.

    Property (i) is obvious, and for property (ii) consider $(\concL{\sigma_2} \ominus \concL\me) \pltimes{\anyL{\sigma_2}} \anyL{\sigma_{\con{\me_1}}}$, given by
    \begin{center}
        \begin{tikzcd}
            {\anyL{\sigma_{\con{\me_1}}}}
                \ar[r, "\anyL{\uncon{\me_1}}"]
                & {\anyL{\sigma'}}
                    \ar[r, "\anyL{\me_2}"]
                    & {\anyL{\sigma_2}}
            \\
            {\concL{\sigma_{\con{\me_1}}}}
                \pbsquare{ur} \ar[u, "\concL{\sigma_{\con{\me_1}}}"] \ar[r, "\concL{\uncon{\me_1}}"]
                & {\concL{\sigma'}}
                    \pbsquare{ur} \ar[u, "\concL{\sigma'}"] \ar[r, "\concL{\me_2}"]
                    & {\concL{\sigma_2}} 
                        \ar[u, "\concL{\sigma_2}"]
            \\
            {\concL{\sigma_{\con{\me_1}}} \ominus \concL{\con{\me_1}}}
                \pbsquare{ur} \ar[u, "\compl{\concL{\con{\me_1}}}"] \ar[r, "\concL{\uncon{\me_1}}"]
                & {\concL{\sigma'} \ominus \concL{\me_1}}
                    \pbsquare{ur} \ar[u, "\compl{\concL{\me_1}}"] \ar[r, "\concL{\me_2}"]
                & {\concL{\sigma_2} \ominus \concL\me}
                    \ar[u, "\compl{\concL{\me}}"]
        \end{tikzcd}
    \end{center}
    and we show that locations in every $\proj_{\concL{\me_2}\concL{\uncon{\me_1}}(\va)} \mvarM_\me$ factor through $\anyL{\sigma_{\con{\me_1}}}$ 
    \begin{align*}
        \proj_{\concL{\me_2}\concL{\uncon{\me_1}}(\va)} \mvarM_\me & = \proj_{\concL{\me_2}\concL{\uncon{\me_1}}(\va)} \mvarM_{\me_2 \comp \me_1} \\
                                                                    & = \sem{\phi(\vC)}\anyL{\me_2}(\proj_{\concL{\uncon{\me_1}}(\va)} \mvarM_{\me_1}) \\
                                                                    & = \sem{\phi(\vC)}\anyL{\me_2 \comp \uncon{\me_1}}(\proj_\va \mvarM_{\con{\me_1}})
    \end{align*}

    By the UP of reachable locations we have $\anyL\vh : \anyL{\sigma_{\con\me}} \to \anyL{\sigma_{\con{\me_1}}}$, which we extend to $\tinstC$ by
    \begin{center}
    \begin{tikzcd}
        {\anyL{\sigma_{\con\me}}}
            \ar[r, "\anyL\vh"] \ar[rrr, bend left=30, "\anyL{\uncon\me}"]
            & {\anyL{\sigma_{\con{\me_1}}}}
                \ar[r, "\anyL{\uncon{\me_1}}"]
                & {\anyL{\sigma'}}
                    \ar[r, "\anyL{\me_2}"]
                    & {\anyL{\sigma_2}}
        \\
        {\concL{\sigma_{\con\me}}}
            \pbsquare{ur} \ar[r, dashed, "\concL\vh"] \ar[u, "\concL{\sigma_{\con\me}}"] \ar[rrr, bend right=30, "\concL{\uncon\me}"]
            & {\concL{\sigma_{\con{\me_1}}}}
                \pbsquare{ur} \ar[r, "\concL{\uncon{\me_1}}"] \ar[u, "\concL{\sigma_{\con{\me_1}}}"]
                & {\concL{\sigma'}}
                    \pbsquare{ur} \ar[r, "\anyL{\me_2}"] \ar[u, "\concL{\sigma'}"]
                    & {\concL{\sigma_2}}
                        \ar[u, "\concL{\sigma_2}"]
    \end{tikzcd}
    \end{center}
    Notice that since the outer square and the right (combined) square are pullbacks, so is the left square,
    and $\con{\me_1} = \vh \comp \con\me$ and $\uncon\me = \me_2 \comp \uncon{\me_1} \comp \vh$.

    We further show that
    \[ \anyL{\sigma_1} \xto{\anyL{\con\me}} \anyL{\sigma_{\con\me}} (\xto{\anyL\vh} \anyL{\sigma_{\con{\me_1}}} \xto{\anyL{\uncon{\me_2}}} \anyL{\sigma'}) \]
    satisfies \cref{def:reach-loc} for $\anyL{\sigma_1} \xto{\anyL{\me_1}} \anyL{\sigma'}$ and $\mvarM_{\me_1}$.

    Property (i) is obvious, and for property (ii) consider $(\concL{\sigma'} \ominus \concL{\me_1}) \pltimes{\anyL{\sigma'}} \anyL{\sigma_{\con\me}}$, given by 
    \begin{center}
        \begin{tikzcd}
            {\anyL{\sigma_{\con\me}}} \ar[r, "\anyL\vh"] 
                & {\anyL{\sigma_{\con{\me_1}}}}
                    \ar[r, "\anyL{\uncon{\me_2}}"]
                    & {\anyL{\sigma'}}
            \\
            {\concL{\sigma_{\con\me}}}
                \pbsquare{ur} \ar[u, "\concL{\sigma_{\con\me}}"] \ar[r, "\concL\vh"]
                & {\anyL{\sigma_{\con{\me_1}}}}
                    \pbsquare{ur} \ar[u, "\concL{\sigma_{\con{\me_1}}}"] \ar[r, "\concL{\uncon{\me_2}}"]
                    & {\concL{\sigma'}}
                        \ar[u, "\concL{\sigma'}"]
            \\
            {\concL{\sigma_{\con\me}} \ominus \concL{\con\me}}
                \pbsquare{ur} \ar[u, "\compl{\concL{\con\me}}"] \ar[r, "\concL\vh"]
                & {\concL{\sigma_{\con{\me_1}}} \ominus \concL{\con{\me_1}}}
                    \pbsquare{ur} \ar[u, "\compl{\concL{\con{\me_1}}}"] \ar[r, "\concL{\uncon{\me_2}}"]
                    & {\concL{\sigma'} \ominus \concL{\me_1}}
                        \ar[u, "\compl{\concL{\me_1}}"]
        \end{tikzcd}
    \end{center}
    and we show that locations in every $\proj_{\concL{\uncon{\me_1}}\concL\vh(\va)} \mvarM_{\me_1}$ factor through $\anyL{\sigma_{\con\me}}$ 
    \begin{align*}
        \proj_{\concL{\uncon{\me_1}}\concL\vh(\va)} \mvarM_{\me_1} & = \proj_{\concL{\uncon{\me_1}}\concL\vh(\va)} \mvarM_{\uncon{\me_1} \comp \con{\me_1}} \\
                                                                    & = \sem{\phi(\vC)}\anyL{\uncon{\me_1}}(\proj_{\concL\vh(\va)} \mvarM_{\con{\me_1}}) & \tag{def of $\mvarM_{\uncon{\me_1} \comp \con{\me_1}}$}\\
                                                                    & = \sem{\phi(\vC)}\anyL{\uncon{\me_1} \comp \anyL\vh}(\proj_\va \mvarM_{\con\me}) & \tag{UP of $\anyL{\sigma_{\con\me}}$}
    \end{align*}

    As a result, we get a $\anyL{\inv^\vh} : \anyL{\sigma_{\con{\me_1}}} \to \anyL{\sigma_{\con\me}}$ which is the inverse of $\anyL\vh$ by
    \begin{align*}
        \anyL\vh \comp \anyL{\inv^\vh} & = 1 \iff \\
        \anyL{\uncon\me_1} \comp \anyL\vh \comp \anyL{\inv^\vh} & = \anyL{\uncon{\me_1}} \iff \\
        \anyL{\me_2} \comp \anyL{\uncon\me_1} \comp \anyL\vh \comp \anyL{\inv^\vh} & = \anyL{\me_2} \comp \anyL{\uncon{\me_1}} \iff \\
        \anyL{\uncon\me} \comp \anyL{\inv^\vh} & = \anyL{\me_2} \comp \anyL{\uncon{\me_1}}
    \end{align*}
    and $\anyL\vh \comp \anyL{\inv^\vh} \comp \anyL\vh = \anyL\vh \iff \anyL{\inv^\vh} \comp \anyL\vh = 1$.

    We extend $\anyL{\inv^\vh}$ into a morphism in $\tempC$ by
    \begin{center}
    \begin{tikzcd}
        {\anyL{\sigma_{\con{\me_1}}}} 
            \ar[r, "\anyL{\inv^\vh}"] \ar[rr, bend left=30, "\anyL{\me_2} \comp \anyL{\uncon{\me_1}}"]
            & {\anyL{\sigma_{\con\me}}}
                \ar[r, "\anyL{\uncon\me}"]
                    & {\anyL{\sigma_2}}
        \\
        {\concL{\sigma_{\con{\me_1}}}} 
            \pbsquare{ur} \ar[u, "\concL{\sigma_{\con{\me_1}}}"] \ar[r, dashed, "\concL{\inv^\vh}"] \ar[rr, bend right=30, "\concL{\me_2} \comp \concL{\uncon{\me_1}}"']
            & {\concL{\sigma_{\con\me}}}
                \pbsquare{ur} \ar[u, "\concL{\sigma_{\con\me}}"] \ar[r, "\concL{\uncon\me}"]
                    & {\concL{\sigma_2}}
                        \ar[u, "\concL{\sigma_2}"]
    \end{tikzcd}
    \end{center}
    where the left square is a pullback because the outer and right square are pullbacks, and regular monos compose (\cref{lem:rmono-comp}).
    As a result, $\concL{\inv^\vh}$ is also iso.
    
    Uniqueness of iso $\vh$ follows from uniqueness of the underlying components.  
\end{proof}

\begin{lemma}[Orthogonal Factorisation System in $\tinstC$]\label{lem:ofs-tinstC}
    The classes of epis (connected) and regular monos (disconnected) morphisms gives rise to an orthogonal factorisation system in $\tinstC$.
\end{lemma}
\begin{proof}
    Follows from \cref{lem:fact-uniq-tinstC}, 
    epis being closed under composition and including all isomorphisms,
    and regular monos in $\tinstC$ being closed under composition (\cref{lem:rmono-comp}) and include all isomorphisms.
\end{proof}
 
\subsubsection{Unification}
\label{sec:unification-theaps}

A natural question to ask when dealing with two template heaps that have a common set of public locations is whether they unify,
that is, if there exist instantiations that we can apply to both heaps respectively such that they result in the same heap.
Through the isomorphism $\theapF \iso \homF\tinstC{?_\emptyset}\dash$, we simply phrase this question as whether  
a span $\sigma_1 \xfrom{\me_1} \sigma \xto{\me_2} \sigma_2$ has a cospan $\sigma_1 \xto{\me_1'} \hat\sigma \xfrom{\me_2'} \sigma_2$ such that
the obvious square commutes.

\begin{definition}[Unifier of Morphisms in $\tinstC$]
    We say that two morphisms $\sigma_1 \xfrom{\me_1} \sigma \xto{\me_2} \sigma_2$ unify if there exists a diagram
    \begin{center}
    \begin{tikzcd}
                                      & {\sigma} \ar[dl, "\me_1"'] \ar[dr, "\me_2"] & \\
        {\sigma_1} \ar[dr, "\me_1'"'] &                                             & {\sigma_2} \ar[dl, "\me_2'"] \\
                                      & {\hat\sigma}                                &
    \end{tikzcd}
    \end{center}
    We call $\me_1', \me_2'$ a unifier of $\me_1, \me_2$.
\end{definition}
Notice that some spans do not have any compatible cospan, that is, some pair of heaps do not unify as the following example shows.
\begin{example}
    Consider the following diagram where both morphisms instantiate the variable $\va$ but with different values (true (T) and false (F)).
    \begin{center}
    \begin{tikzcd}[remember picture]
        & {?_\va} \ar[dl] \ar[dr] & \\ 
        |[alias=s1]|{1_\va} & & 
        |[alias=s2]|{1_\va} & & 
    \end{tikzcd}
    \renewcommand\ndsize{9}
    \begin{tikzpicture}[overlay, remember picture] 
        \node[ccell=$\va$] (d1) [above=15pt of s1, xshift=15pt] {\tiny T};
        \node[ccell=$\va$] (d2) [above=15pt of s2, xshift=-15pt] {\tiny F};
    \end{tikzpicture}
    \end{center}
    There cannot be any unifiers that would make the diagram commute as the resulting heaps will always be different.
\end{example}

\paragraph{General Unifiers}

We seek to identify the most general unifers of two morphisms $\sigma_1 \xfrom{\me_1} \sigma \xto{\me_2} \sigma_2$,
that is, the least cospan $\sigma_1 \xto{\me_1'} \hat\sigma \xfrom{\me_2'} \sigma_2$ such that every other cospan factors uniquely thorough it. 
Unfortunately such M.G.U. does not always exist as the following counterexample shows.
\begin{example}\label{ex:mgu-counterexample}
    Consider the following two unifers\nl\nl

    \renewcommand\ndsize{9}
    \hspace{-10pt}
    \begin{tikzcd}[remember picture]
        & {?_\va} \ar[dl] \ar[dr] & \\ 
        |[alias=s1]|{\incl^{\va,\vb,\vd \inclto \va,\vb,\vc,\vd}} \ar[dr] & & 
        |[alias=s2]|{\incl^{\va,\vc,\vd \inclto \va,\vb,\vc,\vd}} \ar[dl] & & \\
        & |[alias=hs]|{1_{\va,\vb,\vc,\vd}} 
    \end{tikzcd}\hspace{-50pt}
    \begin{tikzcd}[remember picture]
        & {?_\va} \ar[dl] \ar[dr] & \\ 
        |[alias=t1]|{\incl^{\va,\vb,\vd \inclto \va,\vb,\vc,\vd}} \ar[dr] & & 
        |[alias=t2]|{\incl^{\va,\vc,\vd \inclto \va,\vb,\vc,\vd}} \ar[dl] & & \\
        & |[alias=ht]|{1_{\va,\vb,\vc,\vd_1, \vd_2}}
    \end{tikzcd}

    \begin{tikzpicture}[overlay, remember picture] 
        \node[ccell=$\vd$] (sd1) [above=35pt of s1, xshift=30pt] {};
        \node[ccell=$\vb$] (sb1) [below=0pt of sd1] {};
        \node[ccell=$\va$] (sa1) [below=0pt of sb1] {};
        \node[vcell=$\vc$] (sc1) [below=0pt of sa1] {};
        \draw[ptr] (sb1.center) to ++\east to ++\north to (sd1.east);
        \draw[ptr] ($ (sa1.center) + (0,\gap pt) $) to ++\east to ++\nnnorth to ($ (sb1.east) + (0,-\gap pt) $);
        \draw[ptr] ($ (sa1.center) + (0,-\gap pt) $) to ++\east to ++\ssouth to (sc1.east);
    
        \node[vcell=$\vb$] (sb2) [above=35pt of s2, xshift=-30pt] {};
        \node[ccell=$\va$] (sa2) [below=0pt of sb2] {};
        \node[ccell=$\vc$] (sc2) [below=0pt of sa2] {};
        \node[ccell=$\vd$] (sd2) [below=0pt of sc2] {};
        \draw[ptr] ($ (sa2.center) + (0,\gap pt) $) to ++\east to ++\nnorth to (sb2.east);
        \draw[ptr] ($ (sa2.center) + (0,-\gap pt) $) to ++\east to ++\sssouth to ($ (sc2.east) + (0,\gap pt) $);
        \draw[ptr] (sc2.center) to ++\east to ++\south to (sd2.east);

        \node[ccell=$\vc$] (sc1p) [left=15pt of hs,yshift=2pt] {};
        \draw[ptr] (sc1p.center) to ++\east to node[midway,xshift=\gap pt, scale=\shpscale] {$\vd$} ++\north;

        \node[ccell=$\vb$] (sb2p) [right=15pt of hs,yshift=5pt] {};
        \draw[ptr] (sb2p.center) to ++\east to node[midway,xshift=\gap pt, scale=\shpscale] {$\vd$} ++\south;

        \node[ccell=$\vd$] (td1) [above=35pt of t1, xshift=30pt] {};
        \node[ccell=$\vb$] (tb1) [below=0pt of td1] {};
        \node[ccell=$\va$] (ta1) [below=0pt of tb1] {};
        \node[vcell=$\vc$] (tc1) [below=0pt of ta1] {};
        \draw[ptr] (tb1.center) to ++\east to ++\north to (td1.east);
        \draw[ptr] ($ (ta1.center) + (0,\gap pt) $) to ++\east to ++\nnnorth to ($ (tb1.east) + (0,-\gap pt) $);
        \draw[ptr] ($ (ta1.center) + (0,-\gap pt) $) to ++\east to ++\ssouth to (tc1.east);
    
        \node[vcell=$\vb$] (tb2) [above=35pt of t2, xshift=-30pt] {};
        \node[ccell=$\va$] (ta2) [below=0pt of tb2] {};
        \node[ccell=$\vc$] (tc2) [below=0pt of ta2] {};
        \node[ccell=$\vd$] (td2) [below=0pt of tc2] {};
        \draw[ptr] ($ (ta2.center) + (0,\gap pt) $) to ++\east to ++\nnorth to (tb2.east);
        \draw[ptr] ($ (ta2.center) + (0,-\gap pt) $) to ++\east to ++\sssouth to ($ (tc2.east) + (0,\gap pt) $);
        \draw[ptr] (tc2.center) to ++\east to ++\south to (td2.east);

        \node[ccell=$\vc$] (tc1p) [left=15pt of ht,yshift=5pt] {};
        \node[ccell=$\vd_2$] (td12p) [below=0pt of tc1p] {};
        \draw[ptr] (tc1p.center) to ++\east to ++\south to (td12p.east);

        \node[ccell=$\vd_1$] (td21p) [right=15pt of ht,yshift=5pt] {};
        \node[ccell=$\vb$] (tb2p) [below=0pt of td21p] {};
        \draw[ptr] (tb2p.center) to ++\east to ++\north to (td21p.east);
    \end{tikzpicture}

    \renewcommand\ndsize{15}
    The unifier on the left results in the heap
    \begin{tikzpicture}[baseline={(b1.base)}]
        \node[ccell=$\vd$] (d1) [] {};
        \node[ccell=$\vb$] (b1) [below=0pt of d1] {};
        \node[ccell=$\va$] (a1) [below=0pt of b1] {};
        \node[ccell=$\vc$] (c1) [below=0pt of a1] {};
        \draw[ptr] (b1.center) to ++\east to ++\north to (d1.east);
        \draw[ptr] ($ (a1.center) + (0,\gap pt) $) to ++\east to ++\nnnorth to ($ (b1.east) + (0,-\gap pt) $);
        \draw[ptr] ($ (a1.center) + (0,-\gap pt) $) to ++\east to ++\sssouth to ($ (c1.east) + (0,\gap pt) $);
        \draw[ptr] (c1.center) to ++\eeast to ++\north to ++\north to ++\north to ++(0,\gap pt) to ($ (d1.east) + (0,\gap pt) $);
    \end{tikzpicture}
    whereas the one on the right results in
    \begin{tikzpicture}[baseline={(b1.base)}]
        \node[ccell=$\vd_1$] (d1) [] {};
        \node[ccell=$\vb$] (b1) [below=0pt of d1] {};
        \node[ccell=$\va$] (a1) [below=0pt of b1] {};
        \node[ccell=$\vc$] (c1) [below=0pt of a1] {};
        \node[ccell=$\vd_2$] (d2) [below=0pt of c1] {};
        \draw[ptr] (b1.center) to ++\east to ++\north to (d1.east);
        \draw[ptr] ($ (a1.center) + (0,\gap pt) $) to ++\east to ++\nnnorth to ($ (b1.east) + (0,-\gap pt) $);
        \draw[ptr] ($ (a1.center) + (0,-\gap pt) $) to ++\east to ++\sssouth to ($ (c1.east) + (0,\gap pt) $);
        \draw[ptr] ($ (c1.center) + (0,-\gap pt) $) to ++\east to ++\sssouth to ($ (d2.east) + (0,\gap pt) $);
    \end{tikzpicture}
    which are mutually incomparable.
\end{example}

Even though most general unifiers do not always exist, we do have that every unifier can be factored though "special" unifiers that are themselves mutually incomparable. 
We call these special unifiers \emph{general unifiers}.

\begin{definition}[General Unifiers]\label{def:general-unifiers}
    We call a unifier $\sigma_1 \xto{\me_1'} \hat\sigma \xfrom{\me_2'} \sigma_2$ general if 
    $\concL{\sigma_1} \xto{\concL{\me_1'}} \concL{\hat\sigma} \xfrom{\concL{\me_2'}} \concL{\sigma_2}$ are jointly surjective and
    $\anyL{\sigma_1} \xto{\anyL{\me_1'}} \anyL{\hat\sigma} \xfrom{\anyL{\me_2'}} \anyL{\sigma_2}$ are jointly surjective.
\end{definition}
\begin{remark}
    By \cref{lem:gen-cospan-epi-prop}, this is equivalent to saying that a unifier is general if it is generated by a jointly surjective cospan
    $\anyL{\sigma_1} \xto{\anyL{\me_1'}} \anyL{\hat\sigma} \xfrom{\anyL{\me_2'}} \anyL{\sigma_2}$.
\end{remark}

\begin{example}
    Both unifiers from \cref{ex:mgu-counterexample} are general unifiers.
    The unifier on the left is generated by 
        \[ \va,\vb,\vc,\vd \xto{1} \va,\vb,\vc,\vd \xfrom{1} \va,\vb,\vc,\vd \]
    whereas the one on the right by 
        \[ \va,\vb,\vc \mplus \vd \xto{1 \mplus \incl_1} \va,\vb,\vc,\vd_1,\vd_2 \xfrom{1 \mplus \incl_2} \va,\vb,\vc \mplus \vd \]
\end{example}

\begin{lemma}[Data of General Unifiers]\label{lem:general-unifier-data}
    If 
    \begin{tikzcd}[sep=small]
        & {\sigma} 
            \ar[dl, "\me_1"'] \ar[dr, "\me_2"] 
    \\ 
    {\sigma_1} 
        \ar[dr, "\me_1'"']
            && {\sigma_2} 
                \ar[dl, "\me_2'"] 
    \\ 
        & {\hat\sigma}
    \end{tikzcd}
    is a general unifier then $\mvarM_{\me_1'}$ and $\mvarM_{\me_2'}$
    are completely determined by $\mvarM_{\me_1}$ and $\mvarM_{\me_2}$.
\end{lemma}
\begin{proof}
    W.l.g. we show that $\mvarM_{\me_2'}$ is completely determined by $\mvarM_{\me_1}$ and the other direction is symmetric.
    Since $\concL{\me_1'}$ and $\concL{\me_2'}$ are jointly surjective, every $\va : \vC \in \concL{\hat\sigma} \ominus \concL{\me_2'}$
    satisfies that there exists a unique $\va' : \vC \in \concL{\sigma_1} \ominus \concL{\me_1}$ s.t. $\concL{\me_1'}(\va') = \va$. 

    This implies by definition of $\mvarM_{\me_2' \comp \me_2}$ and $\mvarM_{\me_2' \comp \me_2} = \mvarM_{\me_1' \comp \me_1}$ that
    \[ \proj_\va \mvarM_{\me_2'} = \proj_\va \mvarM_{\me_2' \comp \me_2} = \proj_{\concL{\me_1'}(\va')} \mvarM_{\me_1' \comp \me_1} = \sem{\phi(\vC)}\anyL{\me_1'}(\proj_{\va'} \mvarM_{\me_1}) \]
\end{proof}

\begin{lemma}[General Unifier Property]\label{lem:general-unifier-prop}
    Every unifier of $\sigma_1 \xfrom{\me_1} \sigma \xto{\me_2} \sigma_2$ factors uniquely (up to iso) through a general unifier.
\end{lemma}
\begin{proof}
    For a unifier $\sigma_1 \xto{\me_1'} \sigma' \xfrom{\me_2'} \sigma_2$ we show the following diagram:
    \begin{center}
    \begin{tikzcd}
            & \sigma \ar[dl, "\me_1"'] \ar[dr, "\me_2"] & 
        \\
        {\sigma_1} \ar[ddr, "\me_1'"'] \ar[dr, "\hat\me_1"]
            && {\sigma_2} \ar[ddl, "\me_2'"] \ar[dl, "\hat\me_2"']
        \\
            & {\hat\sigma} \ar[d, dashed, "\vh" yshift=5pt]
        \\
            & {\sigma'}
    \end{tikzcd}
    \end{center}
    We define the general unifier $\sigma_1 \xto{\hat\me_1} \hat\sigma \xfrom{\hat\me_2} \sigma_2$ as follows.
    \begin{enumerate}
        \item We define $\anyL{\hat\me_1}$ and $\anyL{\hat\me_2}$ as
              \begin{center}
              \begin{tikzcd}
                  {\anyL{\sigma_1}} \ar[dr, dashed, "\anyL{\hat\me_1}"'] \ar[r, "\anyL{\me_1'}"] 
                      & {\anyL{\sigma'}} 
                          & {\anyL{\sigma_2}} 
                              \ar[dl, dashed, "\anyL{\hat\me_2}"] \ar[l, "\anyL{\me_2'}"'] \\ 
                          & {\anyL{\hat\sigma} \eqdef \img{\anyL{\me_1'}} \setUnion \img{\anyL{\me_2'}}} \ar[u, hook, "\anyL\vh"']
              \end{tikzcd}
              \end{center}
              It is clear that they are jointly surjective.
        \item We generate morphisms $\tinstU\hat\me_1$ and $\tinstU\hat\me_2$ from $\anyL{\hat\me_1}$ and $\anyL{\hat\me_2}$ using \cref{lem:gen-cospan},
              and so by the universal property of generated morphisms we have the following diagram in $\tempC$
              \begin{center}
              \begin{tikzcd}
                      & \sigma \ar[dl, "\tinstU\me_1"'] \ar[dr, "\tinstU\me_2"] & 
                  \\
                  {\sigma_1} \ar[ddr, "\tinstU\me_1'"'] \ar[dr, "\tinstU\hat\me_1"]
                      && {\sigma_2} \ar[ddl, "\tinstU\me_2'"] \ar[dl, "\tinstU\hat\me_2"']
                  \\
                      & {\hat\sigma} \ar[d, dashed, "\tinstU\vh" yshift=5pt]
                  \\
                      & {\sigma'}
              \end{tikzcd}
              \end{center}
              That $\tinstU\hat\me_1 \comp \tinstU\me_1 = \tinstU\hat\me_2 \comp \tinstU\me_2$ follows from the outer square and inner triangles commuting.
        \item For the instantiation data we show that $\mvarM_{\me_\vi'}$ can be factored into $\mvarM_\vh$ and $\mvarM_{\hat\me_\vi}$. 
              W.l.g. let $\vi = 1$.
              Through the canonical isomorphism
              \[ \concL{\sigma'} \ominus \concL{\vu\me_1'} \iso (\concL{\sigma'} \ominus \concL{\vu\vh}) \mplus (\concL{\hat\sigma} \ominus \concL{\vu\hat\me_1}) \]
              we have heaps
\[ \mvarM_{\hat\me_1} : \dheapF(\concL{\hat\sigma} \ominus \concL{\hat\me_1}, \anyL{\sigma'}) \;\text{ and }\; \mvarM_\vh : \dheapF(\concL{\sigma'} \ominus \concL\vh, \anyL{\sigma'}) \]
              We further show for every $\va : \vC \in \concL{\hat\sigma} \ominus \hat\me_1$ that $\vd = \proj_\va \mvarM_{\hat\me_1}$ satisfies
              \begin{center}
                  \begin{tikzcd}
                      {1} \ar[dr, "\exists \hat\vd"'] \ar[r, "\vd"] & {\sem{\phi(\vC)}\anyL{\sigma'}} \\
                                                                    &  {\sem{\phi(\vC)}\anyL{\hat\sigma}} \ar[u, "\anyL\vh"']
                  \end{tikzcd}
              \end{center}
              By \cref{lem:gen-cospan-epi-prop} we have that $\concL{\hat\me_1}$ and $\concL{\hat\me_2}$ are jointly surjective, 
              and since $\va : \vC \in \concL{\hat\sigma} \ominus \concL{\hat\me_1}$, and $\vu\hat\me_1 \comp \vu\me_1 = \vu\hat\me_2 \comp \vu\me_2$ we must have a $\va' : \vC \in \concL{\sigma_2} \ominus \concL{\me_2}$ s.t. $\concL{\hat\me_2}(\va') = \va$.
              And so
              \begin{align*}
                  \proj_\va \mvarM_{\hat\me_1} &= \proj_{\concL\vh(\va)} \mvarM_{\me_1'} \\
                                               &= \proj_{\concL\vh(\va)} \mvarM_{\me_1' \comp \me_1} \\
                                               &= \proj_{\concL\vh(\va)} \mvarM_{\me_2' \comp \me_2} \\
                                               &= \proj_{\concL\vh(\concL{\hat\me_2}(\va'))} \mvarM_{\me_2' \comp \me_2} \\
                                               &= \proj_{\concL{\me_2'}(\va')} \mvarM_{\me_2' \comp \me_2} \\
                                               &= \sem{\phi(\vC)}\anyL{\me_2'}(\proj_{\va'} \mvarM_{\me_2}) \\
                                               &= \sem{\phi(\vC)}\anyL\vh(\sem{\phi(\vC)}\anyL{\hat\me_2}(\proj_{\va'} \mvarM_{\me_2}))
              \end{align*}

              Lastly, $\mvarM_{\hat\me_1 \comp \me_1} = \mvarM_{\hat\me_2 \comp \me_2}$ follows from $\vh$ being mono (\cref{lem:monos-tinistC}).
    \end{enumerate}
    To show uniqueness, consider another general unifier $\sigma_1 \xto{\check\me_1} \check\sigma \xfrom{\check\me_2} \sigma_2$.

    Notice that the induced morphism from $\check\sigma \to \sigma'$ must factor through $\hat\sigma$ in an isomorphic way
    \begin{center}
        \begin{tikzcd}
            {\sigma_1}
                \ar[r, "\hat\me_1"]
                    \ar[dr, "\check\me_1"']
                    & {\hat\sigma}
                        & {\sigma_2}
                            \ar[l, "\hat\me_2"']
                                \ar[dl, "\check\me_2"]
            \\
                    & {\check\sigma}
                        \ar[u, dashed, "\iso"]
        \end{tikzcd}
    \end{center}
    It remains to show that their data match which follows from \cref{lem:general-unifier-data}.
\end{proof}

\begin{lemma}[General Unifiers of Identities]\label{lem:general-unifier-id}
    The general unifiers of a span $\sigma \xfrom{1} \sigma \xto\me \sigma'$ is only cospan $\sigma \xto\me \sigma' \xfrom{1} \sigma'$ (up to iso).
\end{lemma}
\begin{proof}
    To show that cospan $\sigma \xto\me \sigma' \xfrom{1} \sigma'$ is a general unifier, we need it to be jointly surjective at the level $\concL\dash$ and $\anyL\dash$, which is immediate from the identity morphism.
    That every other unifier factors through this (and so it is unique up to iso), is immediate.
\end{proof}

\begin{lemma}[Pasting Law for General Unifiers]\label{lem:general-unifier-pasting-law}
    In a diagram
    \begin{tikzcd}
        {\sigma} \ar[r, "\me_1"] \ar[d, "\me_1'"']
            & {\sigma_1} \ar[r, "\me_3"] \ar[d, "\me_2"]
                & {\sigma_3} \ar[d, "\me_3'"]
        \\
        {\sigma_1'} \ar[r, "\me_2'"']
            & {\sigma_2} \ar[r, "\me_3''"']
                & {\sigma_3'}
    \end{tikzcd}
    where the left square is a general unifier,  
    then the whole square is a general unifier iff the right square is.
\end{lemma}
\begin{proof}
    For the $\to$ direction, it is clear that if any $\anyL{\me_3''}$ and $\anyL{\me_3'}$ are not jointly surjective then the whole square cannot be as well which is a contradiction.
    The same applies to the $\concL\dash$ level, and so the right square must be a general unifier.
    \nl
    For the $\from$ direction, consider a $\va : \vC \in \anyL{\sigma_3'}$.
    Then by the right square being a general unifier, either $\va$ is mapped from $\anyL{\sigma_3}$ or it comes from $\anyL{\sigma_2}$.
    If it comes from $\anyL{\sigma_2}$, then by the left square being a general unifier either it is mapped from $\sigma_1'$ or it is mapped from $\sigma_1$ and since it is a commuting square, it must also be mapped from $\sigma_3$.
    The same reasoning applies to $\concL\dash$.
\end{proof}

\begin{lemma}[General Unifiers Preserve Connectivity]\label{lem:general-unifier-con-prop}
    If 
    \begin{tikzcd}[sep=small]
        & {\sigma} 
            \ar[dl, twoheadrightarrow, "\me_1"'] \ar[dr, "\me_2"] 
    \\ 
    {\sigma_1} 
        \ar[dr, "\me_1'"']
            && {\sigma_2} 
                \ar[dl, "\me_2'"] 
    \\ 
        & {\hat\sigma}
    \end{tikzcd}
    is a general unifier and $\me_1$ is connected, then so is $\me_2'$.
\end{lemma}
\begin{proof}
    Consider a $\va : \vC \in \anyL{\hat\sigma} \ominus \anyL{\uncon{\me_2'}}$.
    This means that $\va : \vC \in \anyL{\hat\sigma} \ominus \anyL{\me_2'}$, and so there exists $\va' : \vC \in \anyL{\sigma_1}$ s.t.
    $\anyL{\me_1'}(\va') = \va$ by definition of general unifiers (\cref{def:general-unifiers}).
    From the factorisation we have a diagram which leads to a contradiction that it commutes unless $\anyL{\hat\sigma} \ominus \anyL{\uncon{\me_2'}} = \emptyset$.
    \begin{center}
        \begin{tikzcd}
            & {\sigma} 
                \ar[dl, twoheadrightarrow, "\me_1"'] \ar[d, twoheadrightarrow, "\con{\me_2' \comp \me_2}" yshift=-3pt] \ar[dr, "\me_2"]
        \\
        {\sigma_1}
            \ar[ddr, "\me_1'"'] \ar[r, twoheadrightarrow, dashed]
            & {\sigma_{\con{\me_2' \comp \me_2}}}
                \ar[dd, -rmonotip] \ar[rd, -rmonotip, dashed]
                & {\sigma_2}
                    \ar[d, twoheadrightarrow, "\con{\me_2'}"]
        \\
                && {\sigma_{\con{\me_2'}}} 
                    \ar[dl, -rmonotip, "\uncon{\me_2'}"]
        \\
            & {\hat\sigma}
        \end{tikzcd}
    \end{center}
    Since $\anyL{\uncon{\me_2'}}$ is iso and $\uncon{\me_2'}$ is a regular mono, we have that $\concL{\uncon{\me_2'}}$ is iso and thus connected.
\end{proof}


\begin{lemma}[Decidability of Unification]\label{lem:unification-dec}
    There exists an algorithm that decides for a span $\sigma_1 \xfrom{\me_1} \sigma \xto{\me_2} \sigma_2$ if they unify or not.
    If they unify, it returns a non-empty finite set of general unifiers
    $\GU{\me_1}{\me_2} \subseteqfin \Coprod{\hat\sigma \in \objC\tempC}{\homF\tinstC{\sigma_1}{\hat\sigma} \times \homF\tinstC{\sigma_2}{\hat\sigma}}$
    that contains all its general unifiers up to isomorphism.
\end{lemma}
\begin{proof}
    By \cref{lem:general-unifier-prop} we have that every unifier factors through a general one and there are only a finite number of isomorphic classes of general unifiers by definition of their jointly surjective property.
    As a result, the algorithm will enumerate every isomorphic class of jointly surjective cospans $\anyL{\sigma_1} \xto{\tau_1} \hat\vw \xfrom{\tau_2} \anyL{\sigma_2}$
    by considering a representative for each class (note that $\objC\worldC$ is countable when $\typeSort$ is, which makes it possible to constructively pick one). 

    It will then equip every representative cospan with data given by \cref{lem:gen-cospan-data} for every side which by \cref{lem:general-unifier-data} tells us that it should be the only data that we have to consider. 
    All that is left to check for it to be a general unifier is that it forms a commuting square in $\tinstC$.
    If the algorithm finds no such commuting generated cospans then it must be that the span does not unify.
\end{proof}

\paragraph{Unification and Local Independence Structure}\nl

\begin{definition}[Local Independent Unifiers]
    A unifier is called a local independent unifier (L.I.U.) if its $\tinstU$ image is a local independent square.
\end{definition}

\begin{definition}[Local Independent General Unifiers]
    A unifier is called a local independent general unifier (L.I.G.U.) if its $\tinstU$ image is a local independent coproduct square.
\end{definition}

\begin{lemma}[Local Independent General Unifier is General]
    A local independent general unifier is a general unifier.
\end{lemma}
\begin{proof}
    To show that a L.I.G.U. $\sigma_1 \xto{\me_1'} \hat\sigma \xto{\me_2'} \sigma_2$ 
    is a general unifier it is enough to show that it is generated by
    $\anyL{\sigma_1} \xto{\anyL{\me_1'}} \anyL{\hat\sigma} \xfrom{\anyL{\me_2'}} \anyL{\sigma_2}$ 
    and moreover this cospan is jointly surjective.
    This is the case as L.I.C. in $\tempC$ is generated from the L.I.C. in $\worldC$ which is itself a jointly surjective cospan. 
\end{proof}

\begin{lemma}[Uniqueness of Local Independent General Unifiers]\label{lem:ligu-uniq}
    L.I.G.U. are unique up to iso.
\end{lemma}
\begin{proof}
    Any two L.I.G.U. are iso in $\tempC$, by L.I.C. being unique up to iso.
    That their data must also be match follows from \cref{lem:general-unifier-data}.
\end{proof}

\begin{lemma}[Local Independent General Unifiers Property]\label{lem:ligu-prop}
    The general unifier of a local independent unifier is a local independent general unifier.
\end{lemma}
\begin{proof}
    The general unifier of a L.I.U. $\sigma_1 \xto{\me_1'} {\sigma'} \xfrom{\me_2'} \sigma_2$ is given by
    the generated cospan over
    \begin{center}
    \begin{tikzcd}
        {\anyL{\sigma_1}} \ar[dr, dashed] \ar[r, "\anyL{\me_1'}"]
            & {\anyL{\sigma'}}
                & {\anyL{\sigma_2}}
                    \ar[dl, dashed] \ar[l, "\anyL{\me_2'}"'] \\
                & {\img{\anyL{\me_1'}} \setUnion \img{\anyL{\me_2'}}} \ar[u, hook]
    \end{tikzcd}
    \end{center}
    And since it is a local independent unifier we have
    \begin{center}
    \begin{tikzcd}
        {\anyL{\sigma_1}} \ar[dr, dashed] \ar[r, "\ext^{\anyL{\me_1}}\anyL{\me_2}"]
            & {\anyL{\sigma_1} \liplus{\anyL\sigma} \anyL{\sigma_2}}
                & {\anyL{\sigma_2}}
                    \ar[dl, dashed] \ar[l, "\ext^{\anyL{\me_2}}\anyL{\me_1}"'] \\
                & {\img{\anyL{\me_1'}} \setUnion \img{\anyL{\me_2'}}} \ar[u, "\iso"]
    \end{tikzcd}
    \end{center}
    which is exactly the cospan that generates L.I.G.U.
\end{proof}

\begin{lemma}[Pasting Law for Local Independent General Unifiers]\label{lem:ligu-pasting}
    In a diagram
    \begin{tikzcd}
        {\sigma} \ar[r, "\me_1"] \ar[d, "\me_1'"']
            & {\sigma_1} \ar[r, "\me_3"] \ar[d, "\me_2"]
                & {\sigma_3} \ar[d, "\me_3'"]
        \\
        {\sigma_1'} \ar[r, "\me_2'"']
            & {\sigma_2} \ar[r, "\me_3''"']
                & {\sigma_3'}
    \end{tikzcd}
    where the left square is a L.I.G.U., then the whole square is a L.I.G.U. iff the right square is.
\end{lemma}
\begin{proof}
    Immediate from \cref{lem:lic-pasting}.
\end{proof}

\begin{lemma}\label{lem:ligu-rmono-exist}
    Every span of the form $\sigma_1 \xfrom{\me_1} \sigma \xrto{\me_2} \sigma_2$ has a local independent general unifier.
\end{lemma}

\begin{proof}
    We have that the L.I.C. in $\tempC$ commutes
    \begin{tikzcd}[sep=small]
        & {\sigma} 
            \ar[dl, "\tinstU\me_1"'] \ar[dr, -rmonotip, "\tinstU\me_2"] 
    \\ 
    {\sigma_1} 
        \ar[dr, "\ext^{\tinstU\me_1}\tinstU\me_2"']
            && {\sigma_2}
                \ar[dl, "\ext^{\tinstU\me_2}\tinstU\me_1"] 
    \\ 
        & {\sigma_1 \liplus\sigma \sigma_2}
    \end{tikzcd}.
    We need to show that it also commutes in $\tinstC$ with data given by \cref{rem:lic-T-data}.

    Notice that by \cref{lem:lic-vpres}, we have that $\concL{\sigma_1 \liplus\sigma \sigma_2}$ is given by
    \begin{center}
        \begin{tikzcd}
            {\concL\sigma} 
                \ar[r, "\concL{\tinstU\me_2}"] \ar[d, "\concL{\tinstU\me_1}"']
                & {\concL{\sigma_2}}
                    \ar[d, "\concL{\ext^{\tinstU\me_2}\tinstU\me_1}"]
            \\
            {\concL{\sigma_1}}
                \ar[r, "\concL{\ext^{\tinstU\me_1}\tinstU\me_2}"']
                    & {\concL{\sigma_1} \liplus{\concL\sigma} \concL{\sigma_2}}
        \end{tikzcd}
    \end{center}
    and $\ext^{\tinstU\me_1}\tinstU\me_2$ is variable preserving.

    We equip $\ext^{\tinstU\me_1}\tinstU\me_2$ with data $\dheapF(\iso, \anyL{\ext^{\tinstU\me_2}\tinstU\me_1})(\mvarM_{\me_2})$ with the canonical iso
    \[ (\concL{\sigma_1} \liplus{\concL\sigma} \concL{\sigma_2}) \ominus \concL{\ext^{\tinstU\me_1}\tinstU\me_2} \iso \concL{\sigma_2} \ominus \concL{\tinstU\me_2} \]
    We equip $\ext^{\tinstU\me_2}\tinstU\me_1$ with data in the same way.
    This gives that $\mvarM_{\ext^{\tinstU\me_1}\me_2 \comp \me_1} = \mvarM_{\ext^{\tinstU\me_2}\me_1 \comp \me_2}$.
\end{proof}

\subsubsection{Preservation of Connectivity}

Given $\sigma_1 \xsurjto\me \sigma_2 \xto{\me'} \sigma_3$, we ask whether the connectivity of $\me$ is preserved,
that is, whether every location introduced by $\me$ will remain reachable in the composite $\me' \comp \me$.

\begin{definition}[Connectedness Preserving Morphisms]
    A morphism $\me_2 : \sigma_2 \to \sigma_3$ preserves connectivity of a morphism $\me_1 : \sigma_1 \surjto \sigma_2$
    if it satisfies $\anyL{\sigma_2} \pltimes{\anyL{\sigma_3}} (\anyL{\sigma_3} \ominus \anyL{\uncon{\me_2 \comp \me_1}}) = \emptyset$.
\end{definition}

\begin{lemma}[Connectedness Preserving Morphisms Property Equivalence]\label{lem:con-pres-prop-equiv}
    Let $\sigma_1 \xsurjto{\me_1} \sigma_2 \xto{\me_2} \sigma_3$.
    We have the following equivalence.

    \begin{center}
    \begin{tikzcd}
        {\emptyset} 
            \ar[r] \ar[d] \pbsquare{dr} 
                & {\anyL{\sigma_3} \ominus \anyL{\uncon{\me_2 \comp \me_1}}}
                    \ar[d, "\compl{\anyL{\uncon{\me_2 \comp \me_1}}}"] 
        \\
        {\anyL{\sigma_2}} 
            \ar[r, "\anyL{\me_2}"']
                & {\anyL{\sigma_3}}
    \end{tikzcd} $\Leftrightarrow$
    \begin{tikzcd}
        {\anyL{\sigma_2}}
            \ar[d, dashed, "\exists \reach_{\me_1}^{\me_2 \comp \me_1}"'] \ar[rd, "\anyL{\me_2}"]
        \\
        {\anyL{\sigma_{\con{\me_2 \comp \me_1}}}}
            \ar[r, "\anyL{\uncon{\me_2 \comp \me_1}}"' yshift=-5pt] 
            & {\anyL{\sigma_3}}
    \end{tikzcd}
    \end{center}
\end{lemma}
\begin{proof}\nl
    $\Rightarrow$) 
        Consider the canonical iso $\anyL{\sigma_3} \iso \anyL{\sigma_{\con{\me_2 \comp \me_1}}} \oplus (\anyL{\sigma_3} \ominus \anyL{\uncon{\me_2 \comp \me_1}})$.
        Now if any $\va : \vC \in \anyL{\me_2}$ maps to $\anyL{\sigma_{\con{\me_2 \comp \me_1}}}$, we would get a contradiction $\va : \vC \to \emptyset$ and so
        it must factor through $\anyL{\sigma_{\con{\me_2 \comp \me_1}}}$.
    \nl
    $\Leftarrow$)
        We $\img{\anyL{\me_2}} \setInt \img{\compl{\anyL{\uncon{\me_2 \comp \me_1}}}}$ that must be empty as the image of $\anyL{\me_2}$ factors through $\anyL{\sigma_{\con{\me_2 \comp \me_1}}}$. 
\end{proof}

\Cref{lem:con-pres-prop-equiv} tells us that connectivity preservation is exactly uniqueness of factorisation (\cref{lem:fact-uniq-tinstC}),
and so every morphism preserves connectivity as the following lemma shows.

\begin{lemma}[All Morphisms are Connectedness Preserving]
    Every $\me_2 : \sigma_2 \to \sigma_3$ preserves the connectivity of every $\me_1 : \sigma_1 \surjto \sigma_2$.
\end{lemma}
\begin{proof}
    By \cref{lem:fact-uniq-tinstC}, we have
    \begin{center}
        \begin{tikzcd}
            {\sigma_1} 
                \ar[r, twoheadrightarrow, "\me_1"] \ar[drr, twoheadrightarrow, "\con{\me_2 \comp \me_1}"'] \ar[rrr, bend left=30, "\me_2 \comp \me_1"]
                & {\sigma_2}
                    \ar[r, twoheadrightarrow, "\con{\me_2}"]
                    & {\sigma_{\con{\me_2}}}
                        \ar[d, "\iso"] \ar[r, -rmonotip, "\uncon{\me_2}"]
                        & {\sigma_3}
            \\
                    && {\sigma_{\con{\me_2 \comp \me_1}}}
                        \ar[ur, -rmonotip, "\uncon{\me_2 \comp \me_1}"']
        \end{tikzcd}
    \end{center}
    which proves preservation of connectivity by \cref{lem:con-pres-prop-equiv}.
\end{proof}

\subsubsection{Preservation of Disconnectivity}

Similarly, we ask if a morphism $\me' : \sigma_2 \surjto \sigma_3$ preserves a disconnected morphism $\me : \sigma_1 \rto \sigma_2$,
that is, if every location in $\anyL{\sigma_2} \ominus \anyL\me$ remains unreachable in $\me' \comp \me$. 

Such morphisms can only instantiate variables coming from $\sigma_1$ with data that do not point to locations $\anyL{\sigma_2} \ominus \anyL\me$ in order to 
preserve the unreachability of $\anyL{\sigma_2} \ominus \anyL\me$ in the composite $\me_2 \comp \me_1$.

\begin{definition}[Disconnectedness Preserving Morphisms]\label{def:dcon-pres-morph}
    A morphism $\me_2 : \sigma_2 \to \sigma_3$ preserves the disconnectivity of a morphism $\me_1 : \sigma_1 \rto \sigma_2$
    if it satisfies $(\anyL{\sigma_2} \ominus \anyL{\me_1}) \pltimes{\anyL{\sigma_3}} \anyL{\sigma_{\con{\me_2 \comp \me_1}}} = \emptyset$.
\end{definition}

\begin{lemma}[Disconnectedness Preserving Morphisms Property Equivalence]\label{lem:dcon-pres-prop-equiv}
    Let $\sigma_1 \xrto{\me_1} \sigma_2 \xto{\me_2} \sigma_3$.
    We have the following equivalence.

    \begin{tikzcd}[sep=small]
        {\emptyset} 
            \ar[r] \ar[dd] \pbsquare{ddr} 
            & {\anyL{\sigma_2} \ominus \anyL{\me_1}}
                \ar[d, "\compl{\anyL{\me_1}}"] 
        \\
            & {\anyL{\sigma_2}}
                \ar[d, "\anyL{\me_2}"]
        \\
        {\anyL{\sigma_{\con{\me_2 \comp \me_1}}}} 
            \ar[r, -rmonotip, "\anyL{\uncon{\me_2 \comp \me_1}}"' yshift=-3pt] 
            & {\anyL{\sigma_3}}
    \end{tikzcd}$\mathclap{\xLeftrightarrow{(1)}}$
    \begin{tikzcd}
        {\anyL{\sigma_1}}
            \ar[d, "\anyL{\con{\me_2 \comp \me_1}}"] \ar[r, -rmonotip, "\anyL{\me_1}"] \pbsquare{dr}
            & {\anyL{\sigma_2}}
                \ar[d, "\anyL{\me_2}"'] 
        \\
        {\anyL{\sigma_{\con{\me_2 \comp \me_1}}}}
            \ar[r, -rmonotip, "\anyL{\uncon{\me_2 \comp \me_1}}"' yshift=-5pt] 
            & {\anyL{\sigma_3}}
    \end{tikzcd}$\hspace{5pt}\mathclap{\xLeftrightarrow{(2)}}$
    \begin{tikzcd}
        {\sigma_1}
            \ar[d, twoheadrightarrow, "\con{\me_2 \comp \me_1}"] \ar[r, -rmonotip, "\me_1"] \pbsquare{dr}
            & {\sigma_2}
                \ar[d, "\me_2"'] 
        \\
        {\sigma_{\con{\me_2 \comp \me_1}}}
            \ar[r, -rmonotip, "\uncon{\me_2 \comp \me_1}"' yshift=-5pt] 
            & {\sigma_3}
    \end{tikzcd}$\hspace{5pt}\mathclap{\xLeftrightarrow{(3)}}$
    \begin{tikzcd}[sep=small]
        {\anyL{\sigma_2} \ominus \anyL{\me_1}}
            \ar[d, dashed, "\exists \anyL{\me_2}\res|{\me_1}"'] \ar[r, "\compl{\anyL{\me_1}}"] 
            & {\anyL{\sigma_2}}
                \ar[d, "\anyL{\me_2}"'] 
        \\
        {\anyL{\sigma_3} \ominus \anyL{\uncon{\me_2 \comp \me_1}}}
            \ar[r, "\compl{\anyL{\uncon{\me_2 \comp \me_1}}}"' yshift=-5pt] 
            & {\anyL{\sigma_3}}
    \end{tikzcd}
\end{lemma}
\begin{proof}
    For (1), we consider the canonical isomorphism $\anyL{\sigma_2} \iso \anyL{\sigma_1} \oplus (\anyL{\sigma_2} \ominus \anyL{\me_1})$.
    For $\xRightarrow{(1)}$ we cannot have anything factoring through $(\anyL{\sigma_2} \ominus \anyL{\me_1})$ as we get a contradiction 
    and for $\xLeftarrow{(1)}$, if $\anyL{\sigma_{\con{\me_2 \comp \me_1}}} \xfrom{\pi_1} \vw \xto{\pi_2} (\anyL{\sigma_2} \ominus \anyL{\me_1})$ is the pullback
    on the left, we can form a square in the right one to get a unique injection $\anyL{\sigma_{\con{\me_1}}} \oplus \vw \dashto \anyL{\sigma_{\con{\me_1}}}$
    and so $\vw$ must be empty.
    For $\xRightarrow{(2)}$, by \cref{lem:pullback-vpres}, the diagram in $\tempC$ is a pullback and by 
    \cref{lem:tinstU-pls-refl} it is also a pullback in $\tinstC$.
    For $\xLeftarrow{(3)}$, by \cref{lem:anyL-pls-pres} the $\anyL\dash$-image is a pullback.

    For $\xLeftarrow{(3)}$, we have that 
    \begin{tikzcd}[sep=small]
        {\emptyset} 
            \ar[r] \ar[dd]
            & {\anyL{\sigma_2} \ominus \anyL{\me_1}}
                \ar[r, dashed, "\anyL{\me_2}\res|{\me_1}" yshift=3pt] \ar[d, "\compl{\anyL{\me_1}}"']
                & {\anyL{\sigma_3} \ominus \anyL{\uncon{\me_2 \comp \me_1}}}
                    \ar[ddl, "\compl{\anyL{\uncon{\me_2 \comp \me_1}}}"]
        \\
            & {\anyL{\sigma_2}}
                \ar[d, "\anyL{\me_2}"']
        \\
        {\anyL{\sigma_{\con{\me_2 \comp \me_1}}}}
            \ar[r, "\anyL{\uncon{\me_2 \comp \me_1}}"' yshift=-3pt]
            & {\anyL{\sigma_3}}
    \end{tikzcd}
    is a pullback as $\img{\anyL{\uncon{\me_2 \comp \me_1}}} \setInt \img{\compl{\anyL{\uncon{\me_2 \comp \me_1}}}} = \emptyset$ and so the inner square must be a pullback.

    For $\xRightarrow{(3)}$, the second pullback gives that $\anyL{\me_2}(\compl{\anyL{\me_1}}(\va)) \not\in \img{\anyL{\uncon{\me_2 \comp \me_1}}}$.
    
\end{proof}

\begin{lemma}[Regular Monos are Disconnectedness Preserving]\label{lem:dcon-pres-rmono}
    For $\sigma_1 \xrto{\me_1} \sigma_2 \xrto{\me_2} \sigma_3$, $\me_2$ preserves disconnectivity of $\me_1$.
\end{lemma}
\begin{proof}
    Notice that $\con{\me_2 \comp \me_1}$ is epi and a regular mono and thus iso.
    As a result,
    \begin{center}
        \begin{tikzcd}
            {\anyL{\sigma_1}}
            \ar[r, -rmonotip, "\anyL{\me_1}"] \ar[d, "\anyL{\con{\me_2 \comp \me_1}}"', "\iso"] \pbsquare{dr}
                    & {\anyL{\sigma_2}}
                        \ar[d, -rmonotip, "\anyL{\me_2}"]
            \\
            {\anyL{\sigma_{\con{\me_2 \comp \me_1}}}}
                \ar[r, -rmonotip, "\anyL{\uncon{\me_2 \comp \me_1}}"' yshift=-3pt]
                    & {\anyL{\sigma_3}}
        \end{tikzcd}
    \end{center}
    is a pullback by \cref{lem:pullback-iso-mono}.
\end{proof}

\begin{lemma}[Disconnectedness Preservation is Downwards Closed]\label{lem:dcon-pres-dclosed}
    For $\sigma_1 \xrto{\me_1} \sigma_2 \xto{\me_2} \sigma_3 \xto{\me_3} \sigma_4$, if $\me_3 \comp \me_2$ preserves disconnectivity of $\me_1$ then
    $\me_2$ preserves disconnectivity of $\me_1$.
\end{lemma}
\begin{proof}
    We have a diagram
    \begin{center}
        \begin{tikzcd}
        {\emptyset}
            \ar[r] \ar[dd]
            & {\anyL{\sigma_2} \ominus \anyL{\me_1}} 
                \ar[d, "\compl{\anyL{\me_1}}"]
        \\
            & {\anyL{\sigma_2}}
                \ar[d, "\anyL{\me_2}"]
        \\
        {\anyL{\sigma_{\con{\me_2 \comp \me_1}}}}
            \ar[r, "\anyL{\uncon{\me_2 \comp \me_1}}"] \ar[d, dashed, "\anyL{\lt{\me_3 \comp \uncon{\me_2 \comp \me_1}}}"']
            & {\anyL{\sigma_3}}
                \ar[d, "\anyL{\me_3}"]
        \\
        {\anyL{\sigma_{\con{\me_3 \comp \me_2 \comp \me_1}}}}
            \ar[r, "\anyL{\uncon{\me_3 \comp \me_2 \comp \me_1}}"' yshift=-10pt]
            & {\anyL{\sigma_4}}
        \end{tikzcd}
    \end{center}
    Where the dashed morphism is given by \cref{lem:ofs-tinstC}.
    That $\me_3 \comp \me_2$ preserves disconnectivity of $\me_1$ gives us that the outer square is a pullback,
    and by \cref{lem:pullback-initial-obj}, the inner square is also pullback which translates to 
    $\me_2$ preserves disconnectivity of $\me_1$.
\end{proof}

\begin{lemma}[L.I.G.U. and Disconnectedness Preservation]\label{lem:ligu-dcon}
    Every L.I.G.U. diagram of the form 
    \begin{tikzcd}
        & {\sigma_1} 
            \ar[dl, "\me_1'"'] \ar[dr, -rmonotip, "\me_1"]
    \\
    {\sigma_2'}
        \ar[dr, "\me_2'"']
            && {\sigma_2}
                    \ar[dl, "\me_2"]
    \\
        & {\sigma_3}
    \end{tikzcd}
    satisfies $\me_2$ preserves the disconnectivity of $\me_1$.
\end{lemma}
\begin{proof}
    By \cref{lem:pullback-vpres} we have that the $\tinstU$-image is a pullback in $\tempC$,
    which by \cref{lem:tinstU-pls-refl} it is a pullback in $\tinstC$.

    So we have a pullback
    \begin{center}
        \begin{tikzcd}
            {\sigma_1}
                \ar[r, "\me_1'"] \ar[d, -rmonotip, "\me_1"'] \pbsquare{dr}
                    & {\sigma_2'}
                        \ar[d, -rmonotip, "\me_2'"]
            \\
            {\sigma_2} \ar[r, "\me_2"']
                    & {\sigma_3}
        \end{tikzcd}
    \end{center}
    and notice $\me_2'$ is r.mono by \cref{lem:lic-vpres} and $\tinstU$ reflecting variable preserving morphisms. 
    
    By \cref{lem:pullback-factoring}, we have that 
    \begin{center}
        \begin{tikzcd}
            {\sigma_1}
                \ar[r, twoheadrightarrow, "\con{\me_1'}"] \ar[dd, -rmonotip, "\me_1"'] \pbsquare{ddr}
                    & {\sigma_{\con{\me_2' \comp \me_1'}}}
                        \ar[d, -rmonotip, "\uncon{\me_1'}"]
            \\
                    & {\sigma_2} 
                        \ar[d, -rmonotip, "\me_2'"]
            \\
            {\sigma_2} \ar[r, "\me_2"']
                    & {\sigma_3}
        \end{tikzcd}
    \end{center}
    which proves that $\me_2$ preserves disconnectivity of $\me_1$ by \cref{lem:dcon-pres-prop-equiv}. 
\end{proof}

\begin{lemma}[Disconnectedness Preserving Morphisms are Data Agnostic]\label{lem:dcon-pres-data-agn}
    For 
    \begin{tikzcd}[sep=small]
        {\sigma_1} 
            \ar[r, -rmonotip, "\me_1"]
            & {\sigma_2} \ar[r, "\me_2"]
                & {\sigma_3}
    \end{tikzcd}
    if $\me_2$ preserves disconnectivity of $\me_1$, then for all $\mvarM' : \dheapF(\concL{\sigma_2} \ominus \concL{\me_1}, \anyL{\sigma_2})$,  $\me_2$ preserves disconnectivity of $\me_1' = \pair{\instU\me_1}{\mvarM'}$.
\end{lemma}
\begin{proof}
    We first show that $\anyL{\sigma_1} \xto{\anyL{\con{\me_2 \comp \me_1}}} \anyL{\sigma_{\con{\me_2 \comp \me_1}}} \xto{\anyL{\uncon{\me_2 \comp \me_1}}} \anyL{\sigma_3}$
    satisfies \cref{def:reach-loc} for $\mvarM_{\me_2 \comp \me_1'}$.
    This would give us a $\vh$ that satisfies
    \begin{center}
        \begin{tikzcd}
            {\anyL{\sigma_1}} 
                \ar[r, "\anyL{\con{\me_2 \comp \me_1}}"] \ar[dr, "\anyL{\con{\me_2 \comp \me_1'}}"']
                & {\anyL{\sigma_{\con{\me_2 \comp \me_1}}}}
                    \ar[r, "\uncon{\me_2 \comp \me_1}"]
                        & {\anyL{\sigma_3}}
            \\
                & {\anyL{\sigma_{\con{\me_2 \comp \me_1'}}}}
                    \ar[ur, "\uncon{\me_2 \comp \me_1'}"'] \ar[u, dashed, "\vh"]
        \end{tikzcd}
    \end{center}
    Property (i) follows immediately.
    For Property (ii), first by the lifting property we have
    \begin{tikzcd}
        {\sigma_{\con{\me_2 \comp \me_1}}} 
            \ar[rr, -rmonotip, "\uncon{\me_2 \comp \me_1}"] \ar[dr, dashed, -rmonotip, "\lt{\uncon{\me_2 \comp \me_1}}"']
                && {\sigma_3}
        \\
            & {\sigma_{\con{\me_2}}}
                \ar[ur, -rmonotip, "\uncon{\me_2}"']
    \end{tikzcd},
    and the pullback $\anyL{\sigma_{\con{\me_2 \comp \me_1}}} \pltimes{\anyL{\sigma_3}} (\concL{\sigma_3} \ominus \concL{\me_2 \comp \me_1})$
    is given by the pullback
    ${\concL{\sigma_{\con{\me_2 \comp \me_1}}} \pltimes{\concL{\sigma_{\con{\me_2}}}} (\concL{\sigma_{\con{\me_2}}} \ominus \concL{\con{\me_2} \comp \me_1})}$
    from
    \begin{center}
        \begin{tikzcd}
            {\anyL{\sigma_{\con{\me_2 \comp \me_1}}}}
                \ar[r, "\anyL{\lt{\uncon{\me_2 \comp \me_1}}}"] \ar[rr, bend left=15, "\anyL{\uncon{\me_2 \comp \me_1}}"] 
                    & {\anyL{\sigma_{\me_2}}}
                        \ar[r, "\anyL{\uncon{\me_2}}"]
                        & {\anyL{\sigma_3}}
            \\
            {\concL{\sigma_{\con{\me_2 \comp \me_1}}}}
                \ar[r, "\concL{\lt{\uncon{\me_2 \comp \me_1}}}"] \ar[u, "\concL{\sigma_{\con{\me_2 \comp \me_1}}}"] \pbsquare{ur}
                    & {\concL{\sigma_{\con{\me_2}}}}
                        \ar[r, "\concL{\uncon{\me_2}}"] \ar[u, "\concL{\sigma_{\con{\me_2}}}"] \pbsquare{ur}
                        & {\concL{\sigma_3}}
                            \ar[u, "\concL{\sigma_3}"']
            \\
            {\sigma_{\concL{\con{\me_2 \comp \me_1}}}
             \pltimes\relax
             (\concL{\sigma_{\con{\me_2}}} \ominus \concL{\con{\me_2} \comp \me_1})}
                \ar[r] \ar[u] \pbsquare{ur}
                & {\concL{\sigma_{\con{\me_2}}} \ominus \concL{\con{\me_2} \comp \me_1}}
                    \ar[r, "\concL{\uncon{\me_2}}"' yshift=-4pt] \ar[u, "\compl{\concL{\con{\me_2} \comp \me_1}}"] \pbsquare{ur}
                    & {\concL{\sigma_3} \ominus \concL{\me_2 \comp \me_1}}
                        \ar[u, "\compl{\concL{\me_2 \comp \me_1}}"']
        \end{tikzcd}
    \end{center}
    Now by \cref{lem:dcon-pres-dclosed}, we have that $\con{\me_2}$ preserves disconnectivity of $\me_1$ and so
    \begin{center}
    \begin{tikzcd}
        {\anyL{\sigma_1}}
            \ar[r, "\anyL{\me_1}"] \ar[d, "\anyL{\con{\me_2 \comp \me_1}}"'] \pbsquare{dr}
                & {\anyL{\sigma_2}}
                    \ar[d, "\anyL{\con{\me_2}}"]
        \\
        {\anyL{\sigma_{\con{\me_2 \comp \me_1}}}}
            \ar[r, "\anyL{\lt{\uncon{\me_2 \comp \me_1}}}"']
                & {\anyL{\sigma_{\con{\me_2}}}}
    \end{tikzcd} 
    \end{center}
    which implies that $\res_{\compl{\concL{\con{\me_2} \comp \me_1}}}\concL{\lt{\uncon{\me_2 \comp \me_1}}}$ must factor as
    \begin{center}
    \begin{tikzcd}
            {\concL{\sigma_{\con{\me_2 \comp \me_1}}}
             \pltimes\relax
             (\concL{\sigma_{\con{\me_2}}} \ominus \concL{\con{\me_2} \comp \me_1})}
                \ar[rr, "\res_{\compl{\concL{\con{\me_2} \comp \me_1}}}\concL{\lt{\uncon{\me_2 \comp \me_1}}}"] \ar[dr, dashed]
                    && {\concL{\sigma_{\con{\me_2}}} \ominus \concL{\con{\me_2} \comp \me_1}}
            \\
                & {\concL{\sigma_{\con{\me_2}}} \ominus \concL{\con{\me_2}}}
                    \ar[ur, hook]
    \end{tikzcd}
    \end{center}
    As a result, all locations $\anyL{\sigma_{\con{\me_2 \comp \me_1}}} \pltimes{\anyL{\sigma_3}} (\concL{\sigma_3} \ominus \concL{\me_2 \comp \me_1})$
    in $\mvarM_{\me_2 \comp \me_1'}$ point to cells in $\mvarM_{\con{\me_2}}$.
    Since $\tinstU\me_1 = \tinstU\me_1'$, the locations in these cells factor through $\anyL{\sigma_{\con{\me_2 \comp \me_1}}}$ by \cref{def:reach-loc}
    for $\anyL{\con{\me_2 \comp \me_1}}$ at $\mvarM_{\me_2 \comp \me_1}$.

    We now use $\vh$ to show
    \begin{center}
        \begin{tikzcd}
            {\anyL{\sigma_1}} 
                \ar[r, Rightarrow] \ar[d, "\anyL{\con{\me_2 \comp \me_1'}}"'] \pbsquare{dr}
                & {\anyL{\sigma_1}}
                    \ar[d, "\anyL{\con{\me_2 \comp \me_1}}"'] \ar[rr, "\anyL{\me_1}"] \pbsquare{drr}
                        && {\anyL{\sigma_2}} 
                            \ar[d, "\anyL{\me_2}"]
            \\
            {\anyL{\sigma_{\con{\me_2 \comp \me_1'}}}}
                \ar[r, dashed, "\vh"] \ar[rrr, bend right=30, "\anyL{\uncon{\me_2 \comp \me_1'}}"']
                & {\anyL{\sigma_{\con{\me_2 \comp \me_1}}}}
                    \ar[rr, "\anyL{\uncon{\me_2 \comp \me_1}}"]
                        && {\anyL{\sigma_3}}
        \end{tikzcd}
    \end{center}
    The right square is a pullback by \cref{lem:dcon-pres-prop-equiv}, and the left square is a pullback by \cref{lem:pullback-iso-mono}. 
    As a result by pullback pasting the outer square is a pullback giving us that $\me_2$ preserves the disconnectivity of $\me_1'$.
\end{proof}

\Cref{lem:dcon-pres-data-agn} leads us to a definition of disconnectivity preservation w.r.t. morphisms in $\tempC$.

\begin{definition}[Disconnectedness Preserving Morphisms w.r.t. $\tempC$]
    We say that $\me_2 : \sigma_2 \to \sigma_3$ preserves the disconnectivity of $\sigma_1 \xrto{\rho_1} \sigma_2$ living in $\tempC$, 
    if for all $\mvarM_{\rho_1} : \dheapF(\concL{\sigma_2} \ominus \concL{\rho_1}, \anyL{\sigma_2})$,
        $\me_2$ preserves the disconnectivity of $\pair{\rho_1}{\mvarM_{\rho_1}}$.
\end{definition}

\begin{lemma}[Disconnectedness Preservation and L.I.C.]
    For a span $\sigma_1 \xfrom{\tinstU\me_1} \sigma \xrto{\rho_2} \sigma_2$ in $\tempC$,
    the induced morphism in $\tinstC$, $\ext^{\rho_2}{\me_1}$, given by \cref{lem:lic-exist} and \cref{rem:lic-T-data},
    satisfies: $\ext^{\rho_2}{\me_1}$ preserves the disconnectivity of $\rho_2$.
\end{lemma}
\begin{proof}
    For some $\mvarM_{\rho_2} : \dheapF(\concL{\sigma_2} \ominus \concL{\rho_2}, \anyL{\sigma_2})$, and $\me_2 = \pair{\rho_2}{\mvarM_{\rho_2}}$,
    we have by \cref{lem:ligu-rmono-exist} that
    \begin{tikzcd}
        & {\sigma}
            \ar[dl, "\me_1"'] \ar[dr, -rmonotip, "\me_2"]
    \\
    {\sigma_1}
        \ar[dr, "\ext^{\tinstU\me_1}\me_2"']
            && {\sigma_2}
                \ar[dl, "\ext^{\tinstU\me_2}\me_1"]
    \\
        & {\sigma_1 \liplus\sigma \sigma_2}
    \end{tikzcd}
    with data given by \cref{rem:lic-T-data} commutes in $\tinstC$. 
    Since this is a L.I.G.U. we have by \cref{lem:ligu-dcon} that $\ext^{\tinstU\me_2}\me_1$ preserves the disconnectivity of $\me_2$.
\end{proof}


\subsection{Families of Template Instantiations}\label{sec:famC-tinstC}

We model finite collections of templates and their instantiations using the $\famC\tinstC$ construction, the free finite product completion of $\tinstC$.

\begin{itemize}
    \item Objects: A finite set $\vI$, and an $\vI$-indexed set of templates $\tlistr{\sigma_\vi \in \objC\tinstC}{\vi \in \vI}$.
    \item Morphisms: A $\tlistr{\sigma_\vi}{\vi \in \vI} \xto\phi \tlistr{\tau_\vj}{\vj \in \vJ}$ consists of
                     a mapping $\indF\phi : \vJ \to \vI$ and a $\vJ$-indexed set of morphisms $\tlistr{\phi_\vj : \sigma_{\indF\phi(\vj)} \to \tau_\vj}{\vj \in \vJ}$
\end{itemize}
We have a forgetful functor $\indF\dash : \famC\tinstC \to \opC\finC$ that maps into the indexing sets.

Notice also that $\famC\tinstC$ is isomorphic to the opposite full subcategory of finite sums of covariant representable presheaves, 
\[ \tinstC \xto{\opC\yoneda} \opC{\funC\tinstC\setC}_\sqcup \iso \famC\tinstC \]
Which extends to an embedding of $\tinstC$ into it $\famC\tinstC$ by selecting a singleton family.

\begin{notation}
We use metavariable $\Sigma$ to range over objects in $\famC\tinstC$ and use
set membership notation $\sigma_\vi \in \Sigma$ to range over the templates in $\Sigma$.
We also use list notation $[ \sigma, \sigma, \tau ]$ to write objects in $\famC\tinstC$.
\end{notation}

\begin{lemma}[Coequalizers in $\famC\tinstC$]
    The category $\famC\tinstC$ has coequalizers.
\end{lemma}
\begin{proof}
    Since $\famC\tinstC \iso \opC{\funC\tinstC\setC}_\sqcup$ and inclusions reflect limits, 
    it is enough to check that
    the equalizer of \begin{tikzcd}{\Sigma_2} \ar[r, yshift=3pt, "\phi"] \ar[r, yshift=-3pt, "\psi"'] & {\Sigma_1}\end{tikzcd} in $\funC\tinstC\setC$
    can be represented as a finite sum of representables.

    Equalizers are given pointwise by
    \[ \eqz\phi\psi(\sigma) \iso \{ \inj_{\vj \in \indF{\Sigma_2}} (\me : \sigma_\vj \to \sigma) : \indF\phi(\vj) = \indF\psi(\vj) \Land \me \comp \phi_\vj = \me \comp \psi_\vj \} \]
    Notice that since every $\me$ is mono, we have $\phi_\vj = \psi_\vj$ for every such $\me$ above.

    So define the coequalizer $\coeqz\phi\psi$ in $\famC\tinstC$   
    as exactly the $\sigma_\vj$ that satisfy $\phi_\vj = \psi_\vj$
    \[ \coeqz\phi\psi \eqdef \{ \sigma_\vj \in \Sigma_2 : \indF\phi(\vj) = \indF\psi(\vj) \Land \phi_\vj = \psi_\vj \} \] 
\end{proof}

\begin{lemma}[Pushouts in $\famC\tinstC$]
    The category $\famC\tinstC$ has pushouts given by the general unifiers of \cref{sec:unification-theaps}.
\end{lemma}
\begin{proof}
    We define the pushout
\begin{tikzcd}
    {\Sigma}
        \ar[d, "\phi"'] \ar[r, "\psi"]
        & {\Sigma_2}
            \ar[d, "\ext^\psi\phi"]
    \\
    {\Sigma_1}
        \ar[r, "\ext^\phi\psi"'] 
        & {\Sigma_1 \liplus{\Sigma} \Sigma_2} \posquare{ul}
\end{tikzcd}
    by noticing that pullbacks in $\funC\tinstC\setC$ are given pointwise and
    \[
        \Sigma_1\sigma' \pltimes{\Sigma\sigma'} \Sigma_2\sigma'
                \iso \{ \pair{\inj_{\sigma_\vi}(\sigma_\vi \xto{\me_\vi} \sigma')}{\inj_{\sigma_\vj}(\sigma_\vj \xto{\me_\vj} \sigma'} :
                        \indF\phi(\vi) = \indF\psi(\vj) \Land
                        \begin{tikzcd}[sep=small]
                            & {\sigma_{\indF\phi(\vi)}} 
                                \ar[dl, "\phi_\vi"'] \ar[dr, "\psi_\vj"] 
                        \\
                        {\sigma_\vi}
                            \ar[dr, "\me_\vi"'] 
                            && {\sigma_\vj}
                                    \ar[dl, "\me_\vj"]
                        \\
                            & {\sigma'}
                        \end{tikzcd} \}
    \]
    By \cref{lem:general-unifier-prop}, every such $\me_\vi, \me_\vj$ pair factors through a general unifier, 
    and so $\Sigma_1 \liplus\Sigma \Sigma_2 \in \famC\tinstC$ must be isomorphic to the set of general unifiers of matching pairs $\vi, \vj$. 

    Specifically, we select all pairs that have matching indices 
    $\{ \pair\vi\vj \in \indF{\Sigma_1} \times \indF{\Sigma_2} : \indF\phi(\vi) = \indF\psi(\vj) \}$
    and for every such pair we have by \cref{lem:unification-dec} a finite set $\GU{\phi_\vi}{\psi_\vj}$ 
    which we turn into a finite family of templates $\Sigma_{\GU{\phi_\vi}{\psi_\vj}}$ by projecting out the first component. 
    \[ \Sigma_1 \liplus\Sigma \Sigma_2 \eqdef 
        \Prod{\pair\vi\vj \in \indF{\Sigma_1} \times \indF{\Sigma_2} \,:\, \indF\phi(\vi) = \indF\psi(\vj)} \Sigma_{\GU{\phi_\vi}{\psi_\vj}} \]

    For every vertex of a general unifier, $\sigma_\vk \in \Sigma_{\GU{\phi_\vi}{\psi_\vj}}$,
    we have morphisms $\me^\vk_\vi: \sigma_\vi \to \sigma_\vk$ and $\me^\vk_\vj: \sigma_\vj \to \sigma^\vk$,
    and so define the injections in $\Sigma_1 \liplus\Sigma \Sigma_2$ as
    \begin{align*}
            \ext^\psi\phi \eqdef \tlistr{\me^\vk_\vi}{\triple\vi\vj\vk \in \indF{\Sigma_1 \liplus\Sigma \Sigma_2}} : \Sigma_1 \to \Sigma_1 \pltimes\Sigma \Sigma_2 \\
            \ext^\phi\psi \eqdef \tlistr{\me^\vk_\vj}{\triple\vi\vj\vk \in \indF{\Sigma_1 \liplus\Sigma \Sigma_2}} : \Sigma_2 \to \Sigma_1 \pltimes\Sigma \Sigma_2
    \end{align*}
    The square $\ext^\phi\psi \comp \phi = \ext^\psi\phi \comp \psi$ commutes from $\me^\vk_\vi \comp \phi_\vi = \me^\vk_\vj \comp \psi_\vj$ for every $\triple\vi\vj\vk \in \indF{\Sigma_1 \liplus\Sigma \Sigma_2}$.
\end{proof}

\begin{lemma}
    The embedding $\opC\yoneda : \tinstC \to \famC\tinstC$ preserves equalizers and pullbacks from $\tinstC$.  
\end{lemma}
\begin{proof}
    See \cref{lem:ryoneda-coeqz,lem:ryoneda-pushouts}.
\end{proof}

\begin{definition}[Connected Morphisms]
    A morphism $\phi : \Sigma \to \Sigma'$ is said to be \emph{connected} if every $\phi_\vj$ is connected.
\end{definition}

\begin{definition}[Non-Overlapping Morphisms]
    A morphism $\phi : \Sigma \to \Sigma'$ is said to be \emph{non-overlapping} if 
    every distinct pair of morphisms with the same root, 
    $\sigma_\vi \xfrom{\phi_\vi} \sigma_\vk \xto{\phi_\vj} \sigma_\vj$, 
    do not unify.

    That is, for every $\vi, \vj \in \indF{\Sigma'}$ s.t. $\vi \neq \vj$ we have
    \[ 
        \indF\phi(\vi) = \indF\phi(\vj) \implies 
            \AllTwo{\sigma_\vi \xto{\me_\vi} \hat\sigma}{\sigma_\vj \xto{\me_\vj} \hat\sigma} 
        \me_\vi \comp \phi_\vi \neq \me_\vj \comp \phi_\vj
    \]
\end{definition}

\begin{lemma}\label{lem:fam-tinstC-epi}
     A morphism $\phi : \Sigma \to \Sigma'$ is epi iff it is connected and non-overlapping.
\end{lemma}
\begin{proof}
    An epi $\phi : \Sigma \to \Sigma' \in \famC\tinstC$ is a mono in $\funC\tinstC\setC_\sqcup$.
    Component-wise injectivity is exactly being connected and non-overlapping, and so it follows by \cref{lem:fsreprC-mono}.
\end{proof}

\begin{definition}[Total Morphisms]
    We say that a morphism $\phi : \Sigma \to \Sigma'$ is \emph{total} if for every morphism $\me : \Sigma \to 1_\vw$ 
    there exists a $\tilde\me : \Sigma' \to 1_\vw$ such that
    \begin{center}
    \begin{tikzcd}
        {\Sigma} \ar[rr, "\me"] \ar[dr, "\phi"']
                && {1_\vw}
        \\
            & {\Sigma'} \ar[ur, "\exists \tilde\me"']
    \end{tikzcd}
    \end{center}
    and further, this mapping is compatible w.r.t. connectivity
    \begin{center}
    \begin{tikzcd}
        {\Sigma} \ar[r, twoheadrightarrow, "\me"] \ar[dr, "\phi"']
                & {1_\vw} \ar[r, -rmonotip, "\me'"]
                    & {1_{\vw'}}
        \\
            & {\Sigma'} 
                \ar[u, "\exists \tilde\me"'] \ar[ur, "\exists \widetilde{\me' \comp \me}"']
    \end{tikzcd}
    \end{center}
    The induced morphism $\widetilde{\me'\comp \me}$ for a connected morphism $\me$ and unconnected $\me'$ factors through the induced morphism $\tilde\me$.
\end{definition}
\begin{lemma}
    Total morphisms are closed under composition.
\end{lemma}
\begin{proof}Easy exercise.\end{proof}

\begin{definition}[Complete Covers]
    A \emph{complete} cover is a connected, non-overlapping and total cover.
\end{definition}

\begin{lemma}\label{lem:compl-utotal-famC-tinstC}
    When the set of sorts are inhabited (\cref{def:inhabited-sorts}),
    a morphism $\phi : \Sigma \to \Sigma'$ is complete iff totality is satisfied uniquely
    \begin{center}
    \begin{tikzcd}
        {\Sigma} \ar[rr, "\forall\me"] \ar[dr, "\phi"']
                && {1_\vw}
        \\
            & {\Sigma'} \ar[ur, dashed, "\exists! \tilde\me"']
    \end{tikzcd}
    \end{center}
\end{lemma}
\begin{proof}
    For $\to$, we only have to show that totality is satisfied uniquely which is immediate from $\phi$ being epi.
    \nl
    For $\from$ direction, uniqueness of totality implies that the mapping is compatible w.r.t. connectivity and so it is a total morphism.  
    To show that $\phi$ is also epi, consider $\psi_1, \psi_2 : \Sigma' \to \Sigma''$ s.t. $\psi_1 \comp \phi = \psi_2 \comp \phi$.
    Notice that for every $\vj \in \indF{\Sigma''}$ we have
    \begin{center}
        \begin{tikzcd}
            {\Sigma} \ar[r, "\psi_1 \comp \phi"] \ar[dr, "\phi"']
                & {\Sigma''} \ar[r, "\proj_\vj"]
                    & {\sigma_\vj} \ar[r, "\me"]
                        & {1_{\hat\vw}}
            \\
                & {\Sigma'} \ar[rru, dashed]
        \end{tikzcd}
    \end{center}
    where $\me$ is given by \cref{def:inhabited-sorts-tinstC}.
    By uniqueness we get $\me \comp \proj_\vj \comp \psi_1 = \me \comp \proj_\vj \comp \psi_2$ which implies by $\me$ being mono that
    $\proj_\vj \comp \psi_1 = \proj_\vj \comp \psi_2$ for every $\vj \in \indF{\Sigma'}$.
    Thus $\psi_1 = \psi_2$.
\end{proof}

\begin{lemma}[Connectivity is stable under pushouts]\label{lem:con-pushout}
    In a pushout diagram, 
    \begin{tikzcd}
        {\Sigma} \ar[r, "\phi_2"] \ar[d, "\phi_1"']
            & {\Sigma_2} \ar[d, "\ext^{\phi_2}\phi_1"]
        \\
        {\Sigma_1} \ar[r, "\ext^{\phi_1}\phi_2"']
            & {\Sigma_1 \liplus\Sigma \Sigma_2} \posquare{ul}
    \end{tikzcd}
    if $\phi_1$ is connected, so is $\ext^{\phi_2}\phi_1$.
\end{lemma}
\begin{proof}
    Directly by \cref{lem:general-unifier-con-prop}.
\end{proof}

\begin{lemma}[Connectivity and Non-Overlapping stable under pushousts]\label{lem:epi-pushout}
    In a pushout diagram, 
    \begin{tikzcd}
        {\Sigma} \ar[r, "\phi_2"] \ar[d, "\phi_1"']
            & {\Sigma_2} \ar[d, "\ext^{\phi_2}\phi_1"]
        \\
        {\Sigma_1} \ar[r, "\ext^{\phi_1}\phi_2"']
            & {\Sigma_1 \liplus\Sigma \Sigma_2} \posquare{ul}
    \end{tikzcd}
    if $\phi_1$ is connected and non-overlapping, so is $\ext^{\phi_2}\phi_1$.
\end{lemma}
\begin{proof}
    Connected, non-overlapping morphisms are epis, and epis are stable under pushouts.
\end{proof}
% it requires the stencils to be connected for non-overlapping to be preserved:
% Consider a non-connected singleton cover, <a, 1_{a,b}> where the cell b is unreachable
%  and rebase it over a -> a,b
%  the resulting cover has two stencils, <a,b, 1_{a,b}, i^{a, b', b}> where in the second stencil, the cell b is unreachable.
%  we can come up with instantiations to 1_{a,b,b'} such that they yield the same stencil

\begin{lemma}[Totality is stable under pushouts]\label{lem:total-pushout}
    In a pushout diagram, 
    \begin{tikzcd}
        {\Sigma} \ar[r, "\phi_2"] \ar[d, "\phi_1"']
            & {\Sigma_2} \ar[d, "\ext^{\phi_2}\phi_1"]
        \\
        {\Sigma_1} \ar[r, "\ext^{\phi_1}\phi_2"']
            & {\Sigma_1 \liplus\Sigma \Sigma_2} \posquare{ul}
    \end{tikzcd}
    if $\phi_1$ is total, so is $\ext^{\phi_2}\phi_1$.
\end{lemma}
\begin{proof}
    To show totality for $\ext^{\phi_2}\phi_1$, it is enough to give an induced morphism $\tilde\me : \Sigma_1 \liplus\Sigma \Sigma_2 \to 1_{\hat\vw}$
    for every connected $\me : \Sigma_2 \to 1_{\hat\vw}$.
    By construction, the mapping will be compatible with connectivity.

    Totality of $\phi_1$ at $\Sigma \xto{\phi_2} \Sigma_2 \xto\me 1_{\hat\vw}$ gives us a $\widetilde{\me \comp \phi_2}$ s.t.
    $\me \comp \phi_2 = \phi_1 \comp \widetilde{\me \comp \phi_2}$ and so by the UP of the pushout we get $\tilde\me$.
\end{proof}

\begin{corollary}[Completeness is stable under pushouts]\label{lem:compl-pushout}
    In a pushout diagram, 
    \begin{tikzcd}
        {\Sigma} \ar[r, "\phi_2"] \ar[d, "\phi_1"']
            & {\Sigma_2} \ar[d, "\ext^{\phi_2}\phi_1"]
        \\
        {\Sigma_1} \ar[r, "\ext^{\phi_1}\phi_2"']
            & {\Sigma_1 \liplus\Sigma \Sigma_2} \posquare{ul}
    \end{tikzcd}
    if $\phi_1$ is complete, so is $\ext^{\phi_2}\phi_1$.
\end{corollary}



\section{The Finitary Monad}


\subsection{FGLS Monad}
\label{sec:fgls-monad}

We give an alternative description of the FGLS monad $\monFLS : \funC\worldC\setC \to \funC\worldC\setC$ using natural transformations. 
We utilise this presentation to define the finitary submonad from.
\[
    \monFLS\vX\vw \iso \homF{\funC\instC\setC}{\homF\worldC\vw{\instU(\dash)} \times \heapF}{\vP(\vX\instU \times \heapF)}
\]
where $\instU : \instC \to \worldC$ is the forgetful functor that forgets the heap data.
We further define the functorial action, unit and join of the FGLS monad in terms of this naturality style.

\subsubsection*{Functorial action}

\begin{definition}[Functorial action of $\monFLS$ on renamings]
    Let $\sigma :  \vw \injto \vw'$ be a renaming in $\worldC$ in
\begin{align*}
    \monFLS\vX\sigma & : \monFLS\vX\vw \to \monFLS\vX\vw' \\
    \phi & \mapsto \homF\worldC{\vw'}{\instU(\dash)} \times \heapF \xto{(\dash \comp \sigma) \times 1} 
    \homF\worldC\vw{\instU(\dash)} \times \heapF \xto\phi \vP(\vX\instU \times \heapF)
\end{align*}
\end{definition}

\subsubsection*{Unit of Monad}


\begin{definition}[Unit of $\monFLS$]
\begin{align*}
    \eta_\vX & : \vX \To \monFLS\vX \\
    \vx      & \mapsto_\vw \homF\worldC\vw{\instU(\dash)} \times \heapF \xto{\vX(\dash)(\vx) \times 1} \vX\instU \times \heapF 
                                                          \xto{\eta^\vP_{\vX\instU \times \heapF}} \vP(\vX\instU \times \heapF)
\end{align*}
where $\vX(\dash)(\vx)$ is simply the functorial action of $\vX$ on $\vx$.
\end{definition}

\subsubsection*{Join of Monad}

\begin{definition}[Multiplication of $\monFLS$]
\begin{align*}
    \mu_\vX & : \Exp\monFLS2\vX \To \monFLS\vX \\
    \phi & \mapsto_\vw \homF\worldC\vw{\tinstU(\dash)} \times \heapF \xto\phi \vP(\monFLS\vX\tinstU \times \heapF) \xto{\vP(\eval)} \Exp\vP2(\vX\tinstU \times \heapF)
    \xto{\mu^\vP_{\vX\tinstU \times \heapF}} \vP(\vX\tinstU \times \heapF)
\end{align*}
where $\eval \pair\phi\vh \overeq\vw \phi_\vw\pair{1_\vw}\vh$.
\end{definition}


\subsection{Matching}

\subsubsection{Stencils}

A \emph{stencil} is a template heap with a selection of public locations. 
They represent collections of heaps with a common segment determined by the public locations.

\begin{definition}[Stencils]
    Stencils are given by
    \begin{align*}
    \stencilF\relax & : \opC\worldC \to \funC\tinstC\setC \\
    \stencilF\vw & = \homF\tinstC{?_{\vw}}\dash \iso \homF\worldC\vw{\anyL\dash} \times \theapF
    \end{align*}
    with the obvious functorial actions. We call the contravariant operation $\stencilF\sigma$ as \emph{hiding}.

    \emph{Concrete stencils} are stencils over heaps, that is 
    \[ \stencilF\vw\inclITI\dash = \homF\tinstC{?_{\vw}}{\inclITI\dash} \iso \homF\worldC\vw{\vu(\dash)} \times \heapF : \instC \to \setC \]
\end{definition}

\begin{example}\label{ex:one}

Consider the following stencil given diagrammatically and assume all locations are of a single sort $\vC$
(underlined locations are the public locations, dashed boxes represent location variables)
\begin{center}
\renewcommand\ndsize{15}
\begin{tikzpicture}[node distance=5mm]
    \node[ccell=$\underline\va$] (na) {};
    \node[ccell=$\vb$] (nb) [below=0pt of na] {};
    \node[vcell=$\vc$] (nc) [below=0pt of nb] {};
    \draw[ptr] (na.center) to ++\east to ++\south to (nb.east);
    \draw[ptr] (nb.center) to ++\wwest to ++\south to (nc.west);
\end{tikzpicture}
\end{center}
It represents heaps that have this segment as starting point from $\va$, such as
\begin{center}
\renewcommand\ndsize{15}
\begin{tikzpicture}[node distance=5mm]
    \node[ccell=$\underline\va$] (na) {};
    \node[ccell=$\vb$] (nb) [below=0pt of na] {};
    \node[ccell=$\vc$] (nc) [below=0pt of nb] {};
    \draw[ptr] (na.center) to ++\east to ++\south to (nb.east);
    \draw[ptr] (nb.center) to ++\wwest to ++\south to (nc.west);
    \draw[ptr] (nc.center) to ++\east to ++\south to ++\west to (nc.south);
\end{tikzpicture}
\hspace{2em}
\begin{tikzpicture}[node distance=5mm]
    \node[ccell=$\underline\va$] (na) {};
    \node[ccell=$\vb$] (nb) [below=0pt of na] {};
    \node[ccell=$\vc$] (nc) [below=0pt of nb] {};
    \draw[ptr] (na.center) to ++\east to ++\south to (nb.east);
    \draw[ptr] (nb.center) to ++\wwest to ++\south to (nc.west);
    \draw[ptr] (nc.center) to ++\eeast to ++\north to ++\north to ++\north to ++\wwest to (na.north);
\end{tikzpicture}
\end{center}
\end{example}

\subsubsection{Covers}

\begin{notation}
    Let metavariable $\Sigma$ to range over finite families of templates, that is, objects in $\Sigma \in \objC{\famC\tinstC}$.
\end{notation}

\begin{definition}[Cover]
    Let $\Sigma$ be a finite family of templates, and $\vw \in \objC\worldC$.
    A \emph{cover} $\cov$ over $\pair\Sigma\vw$ is an indexed family of stencils $\Prod{\sigma_\vi \in \Sigma} \stencilF\vw{\sigma_\vi}$,
    that is, it is a morphism $\cov :\, ?_\vw \to \Sigma \in \famC\tinstC$.
\end{definition}

We inherit the terminology from \cref{sec:famC-tinstC}, and talk about connected, non-overlapping, total and complete covers.
We revisit these concepts here in their specialised form.

\begin{definition}[Connected Cover]
    A cover $\cov$ over $\pair\vw\Sigma$ is said to be connected if every stencil $\proj_\vi \comp \cov$ for $\sigma_\vi \in \Sigma$ is connected.
\end{definition}

\begin{definition}[Non-Overlapping Cover]
    A cover $\cov$ over $\pair\Sigma\vw$ is said to be \emph{non-overlapping} if 
    for every distinct $\sigma_\vi, \sigma_\vj \in \Sigma$, we have
    \[ \AllTwo{\sigma_\vi \xto{\me_\vi} \hat\sigma}{\sigma_\vj \xto{\me_\vj} \hat\sigma} 
        \me_\vi \comp \proj_\vi \comp \cov \neq \me_\vj \comp \proj_\vj \comp \cov \]
\end{definition}

\begin{definition}[Total Covers]
    A cover $\vS$ over $\pair\Sigma\vw$ is \emph{total} if for every connected, concrete stencil $\me :\, ?_\vw \to 1_{\vw'}$,
    there exists a $\sigma_\vi \in \Sigma$, $\tilde\me : \sigma_\vi \to 1_{\vw'}$ such that $\tilde\me \comp \proj_\vi \comp \cov = \me$.
\end{definition}

\begin{definition}[Complete Covers]
    A \emph{complete} cover is a connected, non-overlapping and total cover.
\end{definition}

\begin{definition}[Cover Intersection]
Given a cover $\cov$ over $\pair\vw\Sigma$ and a cover $\cov'$ over $\pair\vw{\Sigma'}$ 
we define their intersection $\cov \uwedge\vw \cov'$ as the biggest cover that is included in both of them.

It is given by pushouts in $\famC\tinstC$ 
\begin{center}
\begin{tikzcd}
    {?_\vw}
        \ar[dd, "\cov"'] \ar[rr, "\cov'"] \ar[ddrr, "\cov \uwedge\vw \cov'"]
        && {\Sigma'}
            \ar[dd, "\ext^{\cov'}\cov"] 
    \\\\
    {\Sigma}
        \ar[rr, "\ext^\cov\cov'"']
        && {\Sigma \pltimes{?_\vw} \Sigma'}
\end{tikzcd}
\end{center}
\end{definition}

\paragraph{Hiding Locations in a Cover}

We extend the contravariant hiding action on stencils to covers.
\begin{definition}[Hiding in Covers]
For a renaming $\sigma : \vw \injto \vw'$, and a cover $\cov$ over $\pair{\vw'}\Sigma$, we have a cover over $\pair\vw\Sigma$
by applying the hiding action on every stencil
\[ \hideSt_\sigma(\cov) \eqdef ?_\vw \xto{\sigma} ?_{\vw'} \xto \cov \Sigma \]
\end{definition}
Hiding does not preserve connectivity, non-overlapping, or totality in general.
% Counter-example: for two stencils with public locations a, b: { a |-> True, b |-> False } and { a |-> False, b |-> True },
%                  if we hide along the empty world, non-overlapping proprety breaks. It is not hard to see that completeness breaks as well

\paragraph{Rebasing a Cover}\label{sec:cov-rebase}

We define a covariant action on renamings.

\begin{definition}
For a $\sigma : \vw \injto \vw'$ and cover $\cov$ over $\pair\vw\Sigma$,
we can rebase $\cov$ to a cover $\cov^\sigma$ over $\pair{\vw'}{\Sigma^\sigma}$ for some induced family of templates ${\Sigma^\sigma}$.

We define the construction using pushouts in $\famC\tinstC$
\begin{center}
\begin{tikzcd}
    {?_\vw}
        \ar[d, "\cov"'] \ar[r, "?_\sigma"]
        & {?_{\vw'}}
            \ar[d, "\cov^\sigma"]
    \\
    {\Sigma} \ar[r, "\rebase^\sigma"' yshift=-2pt]
        & {\Sigma^\sigma} \posquare{ul}
\end{tikzcd}
\end{center}
\end{definition}

Cover rebasing preserves connectivity, non-overlapping property (when cover is connected), totality, and thus completeness (see \cref{lem:con-pushout,lem:epi-pushout,lem:total-pushout,lem:compl-pushout}).

\begin{example}
    Consider a singleton cover over $\pair{\va : \vC}{[1_{\va,\vb,\vc}]}$ given by
    \begin{center}
    \renewcommand\ndsize{12}
    \begin{tikzpicture}[node distance=5mm]
        \node[ccell=$\vb$] (nb) {};
        \node[ccell=$\underline\va$] (na) [below=0pt of nb] {};
        \node[ccell=$\vc$] (nc) [below=0pt of na] {};
        \draw[ptr] (na.center) to ++\east to ++\north to (nb.east);
        \draw[ptr] (na.center) to ++\east to ++\south to (nc.east);
    \end{tikzpicture}
    \end{center}
        
    and we rebase it over renaming $\va : \vC \inclto \va : \vC, \vb : \vC, \vc : \vC$, where $\phi(\vC) = \tyUnit + \tyRef\vC \tyProd \tyRef\vC$. 
    The resulting cover is over 
    \[  \pair{\{\va : \vC, \vb : \vC, \vc : \vC\}}
             {[1_{\va,\vb,\vc}, 1_{\va,\vb,\vc}, \incl^{\va,\vb',\vc \inclto \va,\vb',\vc,\vb}, \incl^{\va,\vb,\vc' \inclto \va,\vb,\vc',\vc}, \incl^{\va,\vb',\vc' \inclto \va,\vb',\vc',\vb, \vc}]} \]
    with stencils
    \begin{center}
    \renewcommand\ndsize{12}
    \begin{tikzpicture}[node distance=5mm]
            \node[ccell=$\underline\vb$] (nb1) {};
            \node[ccell=$\underline\va$] (na1) [below=0pt of nb1] {};
            \node[ccell=$\underline\vc$] (nc1) [below=0pt of na1] {};
            \draw[ptr] (na1.center) to ++\east to ++\north to (nb1.east);
            \draw[ptr] (na1.center) to ++\east to ++\south to (nc1.east);
    
            \node[ccell=$\underline\vc$] (nc2) [right=30pt of nb1] {};
            \node[ccell=$\underline\va$] (na2) [below=0pt of nc2] {};
            \node[ccell=$\underline\vb$] (nb2) [below=0pt of na2] {};
            \draw[ptr] (na2.center) to ++\east to ++\north to (nc2.east);
            \draw[ptr] (na2.center) to ++\east to ++\south to (nb2.east);

            \node[ccell=$\vb'$] (nb3) [right=30pt of nc2] {};
            \node[ccell=$\underline\va$] (na3) [below=0pt of nb3] {};
            \node[ccell=$\underline\vc$] (nc3) [below=0pt of na3] {};
            \node[vcell=$\underline\vb$] (nb3') [below=0pt of nc3] {};
            \draw[ptr] (na3.center) to ++\east to ++\north to (nb3.east);
            \draw[ptr] (na3.center) to ++\east to ++\south to (nc3.east);

            \node[ccell=$\underline\vb$] (nb4) [right=30pt of nb3] {};
            \node[ccell=$\underline\va$] (na4) [below=0pt of nb4] {};
            \node[ccell=$\vc'$] (nc4) [below=0pt of na4] {};
            \node[vcell=$\underline\vc$] (nc4') [below=0pt of nc4] {};
            \draw[ptr] (na4.center) to ++\east to ++\north to (nb4.east);
            \draw[ptr] (na4.center) to ++\east to ++\south to (nc4.east);

            \node[ccell=$\vb'$] (nb5) [right=30pt of nb4] {};
            \node[ccell=$\underline\va$] (na5) [below=0pt of nb5] {};
            \node[ccell=$\vc'$] (nc5) [below=0pt of na5] {};
            \node[vcell=$\underline\vb$] (nb5') [below=0pt of nc5] {};
            \node[vcell=$\underline\vc$] (nc5') [below=0pt of nb5'] {};
            \draw[ptr] (na5.center) to ++\east to ++\north to (nb5.east);
            \draw[ptr] (na5.center) to ++\east to ++\south to (nc5.east);
    \end{tikzpicture}
    \end{center}
\end{example}



\subsection{Evaluating}

\subsubsection{Allocation Morphisms}

We call morphisms in $\tempC$ and $\tinstC$ that are both variable preserving and reflecting as allocation morphisms.
We intent to capture the idea that allocations only introduce concrete locations and do not interfere with location variables.

\begin{definition}[Allocation Morphisms]
    A morphism $\rho : \sigma_1 \to \sigma_2$ is an allocation morphism if it satisfies 
        $\concL{\sigma_2} \pltimes{\anyL{\sigma_2}} \varL{\sigma_1} = \emptyset$ and
        $\varL{\sigma_2} \pltimes{\anyL{\sigma_2}} \varL{\sigma_1} = \emptyset$.
\end{definition}
\begin{example}
    The morphism $?_\va \xto{\pair{?_\vb}{\incl^{\va \inclto \va,\vb}}} \incl^{\vb \inclto \va,\vb}$ is an allocation morphism since it does not instantiate nor introduce any variables.
\end{example}

\begin{notation}
    An allocation morphism is represented using the arrow $\ato$.
\end{notation}

\begin{lemma}[Composition of Allocation Morphisms]\label{lem:comp-alloc}
    Allocation morphisms are closed under composition.
\end{lemma}
\begin{proof}
    Follows from \cref{lem:comp-vpres} and \cref{lem:comp-vrefl}.
\end{proof}

\begin{lemma}[Codomain of Allocation Morphisms]\label{lem:codomain-alloc}
    If $\rho : \sigma_1 \to \sigma_2$ is an allocation morphism, then 
    \[ \sigma_2 \iso \begin{tikzcd}
                        {\concL{\sigma_2} \mplus \varL{\sigma_1}} \\
                        {\concL{\sigma_2}} \ar[u, "\incl_1"']
                    \end{tikzcd}\]
\end{lemma}
\begin{proof} We have the canonical isomorphism
    \[ \sigma_2 \iso \begin{tikzcd}
                    {\concL{\sigma_2} \mplus \varL{\sigma_2}} \\
                    {\concL{\sigma_2}} \ar[u, "\incl_1"']
                 \end{tikzcd}\]
    and notice that $\varL{\sigma_1} \iso \varL{\sigma_2}$, 
    from $\varL\rho \comp \inv^{\varL\rho}\res|{\sigma_2} \comp \varL\rho\res|{\sigma_2} = \varL\rho \comp 1$ and
    $\varL{\sigma_2} \comp \varL\rho\res|{\sigma_2} \comp \inv^{\varL\rho}\res|{\sigma_2} = \varL{\sigma_2} \comp 1$.
\end{proof}
\begin{notation}
    If $\rho$ is an allocation morphism, let $\aiso\rho$ be the isomorphism given by \cref{lem:codomain-alloc}.
\end{notation}

\begin{lemma}[Local Independent Coproducts and Allocation Morphisms]\label{lem:lic-alloc}
    In a L.I.C. square where $\rho_2$ is an allocation morphism
    \begin{center}
    \begin{tikzcd}[sep=small]     
                                                     & {\sigma} \ar[ld, "{\rho_1}"'] \ar[-alloctip, rd, "{\rho_2}" black]  & \\
        {\sigma_1} \ar[rd, "{\ext^{\rho_1}\rho_2}"'] &                                                                     & {\sigma_2} \ar[ld, "{\ext^{\rho_2}\rho_1}"] \\
                                                     &              {\sigma_1 \liplus\sigma \sigma_2}                          &
    \end{tikzcd}
    \end{center}
    we have that $\ext^{\rho_1}\rho_2$ is an allocation morphism as well.
\end{lemma}
\begin{proof}
    Follows from \cref{lem:lic-vpres} and \cref{lem:lic-vrefl}.
\end{proof}

\begin{lemma}[Allocation Morphism Cancellation]\label{lem:alloc-cancl}
    In \begin{tikzcd}
        {\sigma_1} 
            \ar[rr, "\rho"] \ar[dr, "\rho_1"']
                && {\sigma_3}
        \\
            & {\sigma_2}
                \ar[ur, "\rho_2"']
    \end{tikzcd}
    if $\rho$ is an allocation morphism then $\rho_1$ is an allocation morphism iff $\rho_2$ is.
\end{lemma}
\begin{proof}
    For ($\to$), it follows from \cref{lem:vrefl-cancl,lem:vpres-vrefl-cancl}, and for ($\from$) it follows from \cref{lem:vpres-cancl,lem:vrefl-vpres-cancl}.
\end{proof}

\begin{lemma}[Allocation Morphisms and Epis]\label{lem:alloc-epi}
    Every $\sigma_1 \xato{\me_1} \sigma_2 \xsurjto{\me_2} \sigma_3$ has the following diagram in $\tempC$
    \begin{center}
    \begin{tikzcd}
        {\sigma_{\uncon{\me_2 \comp \me_1}}}
            \ar[r, dashed, "\exists"] \ar[d, "\compl{\uncon{\me_2 \comp \me_1}}"']
            & {\sigma_2 \ominus \me_1}
                \ar[d, "\compl{\me_1}"]
        \\
        {\sigma_3}
            & {\sigma_2}
                \ar[l, "\me_2"]
    \end{tikzcd}
    \end{center}
\end{lemma}
\begin{proof}
    Notice that if it holds at $\anyL\dash$ level it also holds at $\concL\dash$ from 
    \begin{center}
    \begin{tikzcd}
            & {\concL{\sigma_2}}
                \ar[dr, "\concL{\sigma_2}"]
        \\
        {\concL{\sigma_2} \ominus \concL{\me_1}}
            \ar[ur, "\compl{\concL{\me_1}}"]
            & {\anyL{\sigma_2} \ominus \anyL{\me_1}}
                \ar[l, "\inv^{\sigma_2}\res|{\me_1}"'] \ar[r, "\compl{\anyL{\me_1}}"] 
                & {\anyL{\sigma_2}}
                    \ar[d, "\anyL{\me_2}"'] 
        \\
        {\concL{\sigma_3} \ominus \concL{\uncon{\me_2 \comp \me_1}}}
            \ar[r, "\concL{\sigma_3}\res|{\uncon{\me_2 \comp \me_1}}"' yshift=-5pt]
            & {\anyL{\sigma_3} \ominus \anyL{\uncon{\me_2 \comp \me_1}}}
                \ar[u, dashed, "\exists \inv^{\anyL{\me_2}}\res|{\me_1}"]
                \ar[r, "\compl{\anyL{\uncon{\me_2 \comp \me_1}}}"' yshift=-5pt] 
                & {\anyL{\sigma_3}}
    \end{tikzcd}
    \end{center}

    We construct $\anyL{\sigma_3} \ominus \anyL{\uncon{\me_2 \comp \me_1}} \dashto \anyL{\sigma_2} \ominus \concL{\me_1}$, 
    using the canonical iso
    \[ (\anyL{\sigma_{\con{\me_2 \comp \me_1}}} \ominus \anyL{\con{\me_2 \comp \me_1}}) \oplus 
       (\anyL{\sigma_3} \ominus \anyL{\uncon{\me_2 \comp \me_1}}) \iso (\anyL{\sigma_2} \ominus \anyL{\me_1}) \oplus (\anyL{\sigma_3} \ominus \anyL{\me_2}) \]
    Observe that if there exists $\va, \va' : \vC$ s.t. $\anyL{\compl{\uncon{\me_2 \comp \me_1}}}(\va) = \compl{\anyL{\me_2}}(\va')$
    then there exists $\vb : \vC$, s.t. $\anyL{\uncon{\me_2}}\compl{\anyL{\con{\me_2}}}(\vb) = \va'$ since $\me_2$ is connected. 
    As a result, we have a $\vb : \vC$ that is not in the image of $\anyL{\sigma_2}$ and is reachable in $\me_2$.
    By \cref{def:reach-loc-hwalks}, we have a heap walk that leads to $\vb$ with root location variable in $\varL{\sigma_2}$ and since $\me_1$ is an allocation morphism
    by \cref{lem:vrefl-prop}, this variable comes from $\varL{\sigma_1}$. 
    As a result, $\vb$ must be reachable in $\me_2 \comp \me_1$ which is a contradiction as it maps to the unreachable segment by $\anyL{\uncon{\me_2}}\compl{\anyL{\con{\me_2}}}$.
\end{proof}

\begin{lemma}[Disconnectedness Preservation and L.I.G.U.]\label{lem:dcon-ligu}
    A diagram of the form
    \begin{tikzcd}
        & {\sigma_1} 
            \ar[dl, twoheadrightarrow, "\me_1'"'] \ar[dr, -alloctip, "\me_1"]
    \\
    {\sigma_2'}
        \ar[dr, -rmonotip, "\me_2'"']
            && {\sigma_2}
                    \ar[dl, twoheadrightarrow, "\me_2"]
    \\
        & {\sigma_3}
    \end{tikzcd}
    s.t. $\me_2$ preserves the disconnectivity of $\me_1$ implies that it is a L.I.G.U.
\end{lemma}
\begin{proof}
    We first show that $\me_2'$ is variable reflecting.
    From \cref{lem:alloc-epi,lem:vrefl-prop} we have
    \begin{center}
        \begin{tikzcd}
            {\anyL{\sigma_3} \ominus \anyL{\me_2'}} 
                \ar[r, dashed] 
                    & {\anyL{\sigma_2} \ominus \anyL{\me_1}}
                        \ar[d, "\inv^{\concL{\sigma_2}}\res|{\me_1}"]
            \\
            {\concL{\sigma_3} \ominus \concL{\me_2'}} 
                \ar[r, dashed] \ar[u, "\concL{\sigma_3}\res|{\me_2'}"]
                    & {\concL{\sigma_2} \ominus \concL{\me_1}} 
    \end{tikzcd}
    \end{center}
    which gives us a morphism $\anyL{\sigma_3} \ominus \anyL{\me_2'} \to \concL{\sigma_2} \ominus \concL{\me_1}$ and notice that
    post-composing this with $\concL{\me_2} \comp \compl{\concL{\me_1}}$ must reside in $\concL{\sigma_3} \ominus \concL{\me_2'}$ 
    for otherwise the square of $\me_2'$ would not commute.

    By \cref{lem:lic-vpres}, we need to show that both $\concL\dash$ and $\anyL\dash$ images are L.I.C. 
    Notice that by \cref{lem:dcon-pres-prop-equiv}, its $\anyL\dash$-image is a pullback, 
    and by \cref{lem:pullback-vpres} we have that the $\tinstU$-image is a pullback in $\tempC$,
    and so the $\concL\dash$-image is a pullback as well. 

    So it enough to show that the $\concL\dash$ and $\anyL\dash$ images of cospan $\sigma_2' \xto{\me_2'} \sigma_3 \xfrom{\me_2} \sigma_2$ are jointly surjective.
    We show that $\concL\dash$-image of the cospan is jointly surjective which implies that the $\anyL\dash$-image is jointly surjective
    by the variable reflecting property of $\me_2'$ (\cref{lem:vrefl-prop}).

    Joint surjectivity at $\concL\dash$ is immediate from \cref{lem:alloc-epi}, every $\va : \vC \in \concL{\sigma_3} \ominus \concL{\uncon{\me_2 \comp \me_1}}$ 
    factors through $\concL{\sigma_2} \ominus \concL{\me_1}$ and so it is in the image of $\concL{\me_2}$.
\end{proof}

\begin{definition}[Category of Template Allocations]
    We consider $\allocC\tempC \subseteq \tempC$, the wide subcategory of $\tempC$ where morphisms are allocation morphisms.
    Identities are allocation morphisms and by \cref{lem:comp-alloc}, they are also closed under composition.
    Similarly, we consider $\allocC\tinstC \subseteq \tinstC$ and notice that we have a forgetful functor $\allocC\tinstC \xto{\allocC\tinstU} \allocC\tempC$.
\end{definition}

\subsubsection{Resultlets}

\emph{Resultlets} represent collections of results in the spirit of the FGLS monad.

\begin{definition}
    For a $\vA : \tinstC \to \setC$ and template $\sigma$, we define the resultlet functor $\resF\vA_\sigma$ as
    \[ \resF\vA_\sigma \eqdef \sliceC{\sigma}{\allocC\tinstU} \xto{\vU} \allocC\tinstC \inclto \tinstC \xto\vA \setC \]

    A resultlet $\vr$ over $\pair\sigma\vA$ is an element of $\typeRes\vA\sigma$
    \[ \vr \in \typeRes\vA\sigma \eqdef \Coprod{\sigma \xato{\rho'} \sigma' \in \sliceC\sigma{\allocC\tinstU}}  \resF\vA_\sigma\rho' = 
        \Coprod{\sigma \xato{\rho'} \sigma' \in \sliceC\sigma{\allocC\tinstU}} \vA(\sigma') \]
\end{definition}

\begin{remark}
For FGLS we are interested in situations where $\vA = \vX\anyL\dash \times \theapF$ for a $\vX : \worldC \to \setC$.
A resultlet over $\pair\sigma{\vX\anyL\dash \times \theapF}$ then consists of
\begin{itemize}
    \item An allocation morphism $\sigma \xato\rho \sigma'$, which by \cref{lem:codomain-alloc} 
            is determined only by $\concL\sigma \xto{\concL\rho} \concL{\sigma'}$.
        \item An $\vX$-valued return value $\vx : \vX\anyL{\sigma'}$ which by \cref{lem:codomain-alloc}
                is isomorphic to a value in $\vX(\concL{\sigma'} \oplus \varL\sigma)$.
    \item A template heap $\mvarM : \dheapF(\sigma')$, which is isomorphic to 
            a heaplet $\dheapF(\concL{\sigma'}, \concL{\sigma'} \oplus \varL\sigma)$.
\end{itemize}
\end{remark}

We define actions of $\typeRes\relax$ on morphisms in $\funC\tinstC\setC$ and $\tinstC$ but we stress that some of these actions are not functorial.
In \cref{sec:tinstC-hiding-monad} we quotient $\typeRes\relax$ to form a functor out of it.

\begin{definition}[Action on morphisms in $\funC\tinstC\setC$]
    For a natural transformation $\alpha : \vA \to \vB$, and every $\sigma \in \objC\tempC$
    \begin{align*}
        \typeRes(\alpha)_\sigma & : \typeRes\vA\sigma \to \typeRes\vB\sigma \\
        \inj_{\sigma \xato{\rho'} \sigma'}(\va) & \mapsto \inj_{\rho'}(\alpha_{\sigma'}(\va)) 
    \end{align*}
    It changes a resultlet over $\pair\sigma\vA$ to one over $\pair\sigma\vB$.
\end{definition}

\begin{definition}[Covariant action on morphisms in $\tinstC$]
For a $\me : \sigma_1 \to \sigma_2 \in \tinstC$ and a resultlet $\inj_{\sigma_1 \xato{\rho'} \sigma'} \va$
we consider the following L.I.C. in $\tempC$
\begin{center}
\begin{tikzcd}
        & {\sigma_1} 
            \ar[ld, "\tinstU\me"'] \ar[-alloctip, rd, "\rho'"] 
    \\
    {\sigma_2} 
        \ar[-alloctip, rd, "\ext^{\tinstU\me}\rho'"'] 
            && {\sigma'} 
                \ar[ld, start anchor=south west, "\ext^{\rho'}\me"] 
    \\
        & {\sigma_2 \liplus{\sigma_1} \sigma'}
\end{tikzcd}
\end{center}
Notice that by \cref{lem:lic-alloc}, $\ext^{\tinstU\me}\rho'$ is an allocation morphism.
We further equip $\ext^{\rho'}\me$ with data using \cref{rem:lic-T-data} and so define the covariant action $\typeRes\vA\me$ as
\begin{align*}
    \typeRes\vA\me & : \typeRes\vA{\sigma_1} \to \typeRes\vA{\sigma_2} \\
    \inj_{\rho'}\va & \mapsto \inj_{\ext^{\tinstU\me}{\rho'}} \vA(\ext^{\rho'}\me)(\va)
\end{align*}
\end{definition}

\begin{definition}[Contravariant action on allocation morphisms in $\tempC$]
    For every allocation morphism $\sigma_2 \xato{\rho'} \sigma_1$ in $\tempC$ define
\begin{align*}
    \hideRes\vA{\rho'} & : \typeRes\vA\sigma_1 \to \typeRes\vA\sigma_2 \\ 
    \inj_\rho \va & \mapsto \inj_{\rho \comp \rho'} \va
\end{align*}
where $\rho \comp \rho'$ is an allocation morphism itself by \cref{lem:comp-alloc}.

We call this action \emph{hiding} as it forgets some of the public locations.
\end{definition}

\begin{example}
    Consider the allocation morphism
    \[ \incl^{\va \inclto \va,\vc} \xato\rho \incl^{\va,\vb \inclto \va,\vb,\vc} \] 
    that introduces a new concrete location $\vb$,
    and a resultlet $\inj_{\rho'}\va$ where $\va : \vA(\incl^{\va,\vb \inclto \va,\vb,\vc})$.

    We can apply the covariant action $\me = \incl^{\va \inclto \va,\vc} \xto{\pair{\incl^{\va \inclto \va,\vc}}{\incl^{\va,\vc \inclto \va,\vc,\vd}}} 1_{\va,\vc}$ that instantiates $\vc$ with some data $\mvarM_\me : \dheapF(\vc, \{\va,\vc,\vd\})$

    \begin{center}
    \begin{tikzcd}
            & {\incl^{\va \inclto \va,\vc}} 
                \ar[dl, "\tinstU\me"'] \ar[dr, "\rho"] 
        \\
        {\incl^{\va,\vc \inclto \va,\vc,\vd}} 
            \ar[dr, "\ext^{\tinstU\me}\rho"']
                    && {\incl^{\va,\vb \inclto \va,\vb,\vc}} 
                        \ar[dl, "\ext^\rho\me"] 
        \\
        & {\incl^{\va,\vb,\vc \inclto \va,\vb,\vc,\vd}}
    \end{tikzcd}
    \end{center}
    to get the resultlet 
        \[ \inj_{\ext^{\tinstU\me}\rho} \vA(\ext^\rho\me)(\va) \]
            with data $\mvarM_{\ext^{\rho}{\me}}$ being $\dheapF(1, \anyL{\ext^{\tinstU\me}\rho})(\mvarM_\me)$.

    Consider the allocation morphism $\rho' = ?_\vc \xto{\pair{?_\va}{\incl^{\vc \inclto \va,\vc}}} \incl^{\va \inclto \va,\vc}$.
    We can apply the contravariant action $\rho'$ to get the resultlet $\inj_{\pair{?_{\{\va,\vb\}}}{\incl^{\vc \inclto \va,\vb,\vc}}} \va$.
\end{example}

\paragraph{Monad Structure on Resultlets}

We equip resultlets with unit and multiplication structure.

\begin{definition}[Unit of a Resultlet]
    We define unit morphisms
    \begin{align*}
        \eta^\vA_\sigma & : \vA\sigma \to \typeRes\vA\sigma \\
        \va & \mapsto \inj_{1_\sigma} \va
    \end{align*}
\end{definition}

\begin{definition}[Multiplication of a Resultlet]
    We define multiplication morphisms
    \begin{align*}
        \mu^\vA_\sigma & : \typeRes(\typeRes\vA)\sigma \to \typeRes\vA \\
        \inj_\rho \vr & \mapsto \hideRes\vA\rho\vr
    \end{align*}
\end{definition}

\subsubsection{Hiding Monad on Template Heaps}
\label{sec:tinstC-hiding-monad}

\begin{definition}[Hiding Monad on Template Heaps]
We generalise the hiding monad $\hideMon$ from \cite{kammar2017monad} to template heaps by appropriately quotienting the type of resultlets 
\[\begin{array}{r@{\spac}l@{\spac}l}
    \thideMon              & : \funC\tinstC\setC \to \funC\tinstC\setC \\
    (\thideMon\vA)(\sigma) & = \Colimit{(\resF\vA_\sigma)} \ni \eqclass{\sigma \xato{\rho'} \sigma'} \va
\end{array}\]
which satisfies the natural (monad) isomorphism for every $\vA : \tinstC \to \setC$
\[ (\thideMon\vA) \comp \inclITI\dash \hmonIso \hideMon(\vA \comp \inclITI\dash) \;:\; \instC \to \setC \]

We inject resultlets into the monad using the colimit injections $\inj^{\resF\vA_\sigma}_{\rho'} : \resF\vA_\sigma\rho' \to \thideMon\vA\sigma$
\begin{align*}
    \quot^\vA_\sigma & : \typeRes\vA\sigma \to \thideMon\vA\sigma \\
    \inj_{\rho'}\va & \mapsto \inj^{\resF\vA_\sigma}_{\rho'}(\va)
\end{align*}
\end{definition}

\begin{lemma}[Quotienting in the Hiding Monad on Template Heaps]\label{lem:thmon-quotient}
    For $\inj_{\sigma \xato{\rho_1} \sigma_1} \va_1$, and $\inj_{\sigma \xato{\rho_2} \sigma_2} \va_2$, 
    we have that $\eqclass{\rho_1}\va_1 = \eqclass{\rho_2}\va_2 \in \thideMon\vA\sigma$ iff there exist 
    \begin{tikzcd}
        & {\sigma} 
            \ar[dl, -malloctip, "\rho_1"'] \ar[dr, -alloctip, "\rho_2"]
    \\ 
    {\sigma_1} 
        \ar[dr, -alloctip, "\exists \tinstU\me_1'"'] 
            && {\sigma_2} 
                \ar[dl, -malloctip, "\exists \tinstU\me_2'"]
    \\ 
        & {\hat\sigma}
    \end{tikzcd}
    s.t. $\vA(\me_1')(\va_1) = \vA(\me_2')(\va_2)$.
\end{lemma}
\begin{proof}
    This is a direct generalisation of the proof in \cref{lem:hmon-quotient}. 
    We first show that the RHS relation is an equivalence. 
    Reflexivity and symmetry is immediate.
    For transitivity, we have
    \begin{center}
        \begin{tikzcd}
                && {\sigma_{12} \liplus{\sigma_2} \sigma_{23}}
        \\
            & {\sigma_{12}} 
                \ar[ur, -alloctip, "\ext^{\tinstU\me_{22}}\me_{23}"]
                    && {\sigma_{23}}
                        \ar[lu, -malloctip, "\ext^{\tinstU\me_{23}}\me_{22}"']
        \\
        {\sigma_1}
            \ar[ur, -malloctip, "\tinstU\me_1"]
                && {\sigma_2}
                        \ar[ul, -malloctip, "\tinstU\me_{22}"] \ar[ur, -alloctip, "\tinstU\me_{23}"']
                        && {\sigma_3}
                            \ar[ul, -malloctip, "\tinstU\me_3"']
        \\\\
                && {\sigma}
                    \ar[uull, -malloctip, "\rho_1"] \ar[uu, -alloctip, "\rho_2"] \ar[uurr, -alloctip, "\rho_3"']
        \end{tikzcd}
    \end{center}
    Notice that the top square commutes by \cref{lem:ligu-rmono-exist}
    which implies that $\inj_{\rho_1}(\va_1) \sim \inj_{\rho_3}(\va_3)$.

    To prove that the generated relation by the colimit is exactly $\sim$, one shows direction $\to$ by inducting on the list of related elements,
    and for $\from$, notice that if $\inj_{\rho_1}(\va_1) \sim \inj_{\rho_2}(\va_2)$, then
    \[ \inj_{\rho_1}\va_1 \leq \inj_{\instU\me_1 \comp \rho_1}\vA(\me_1)(\va_1) \ge \inj_{\rho_2}\va_2 \]

\end{proof}

\begin{lemma}
    $\thideMon$ together with the covariant actions on morphisms, unit and multiplication from $\typeRes\relax$ form a monad.
\end{lemma}
\begin{proof}
    \nl
    We first show it is functorial on morphisms $\me : \sigma_1 \to \sigma_2 \in \tinstC$.
    Consider two related $\inj_{\sigma_1 \xto\rho \sigma'} \va'$, $\inj_{\sigma_1 \xto{\rho'} \sigma''} \va''$. 
    By \cref{lem:thmon-quotient} we have
    \begin{tikzcd}
        & {\sigma_1} 
            \ar[dl, -malloctip, "\rho'"'] \ar[dr, -alloctip, "\rho''"] \ar[dd, "\hat\rho"]
    \\ 
    {\sigma'} 
        \ar[dr, -alloctip, "\exists \tinstU\me'"'] 
            && {\sigma''} 
                \ar[dl, -malloctip, "\exists \tinstU\me''"]
    \\ 
        & {\hat\sigma}
    \end{tikzcd}
    with $\vA(\me')(\va') = \vA(\me'')(\va'')$.

    We form the following diagram
    \begin{center}
        \begin{tikzcd}[sep=small]
            & {\hat\sigma \liplus{\sigma'} (\sigma_2 \liplus{\sigma_1} \sigma')} \ar[rr, -alloctip, "\iso'"]
                && {\hat\sigma \liplus{\sigma_1} \sigma_2}
                    && {(\sigma_2 \liplus{\sigma_1} \sigma'') \liplus{\sigma''} \hat\sigma} \ar[ll, -malloctip, "\iso''"']
        \\
        {\hat\sigma} \ar[ur, "\ext^{\tinstU\me'}\ext^{\rho'}\me"'] \ar[urrr, bend left=50, "\ext^{\hat\sigma}\me"]
                && {\sigma' \liplus{\sigma_1} \sigma_2} \ar[ul, -malloctip, "\ext^{\ext^{\rho'}\tinstU\me}\me'"]
                    && {\sigma_2 \liplus{\sigma_1} \sigma''} \ar[ur, -alloctip, "\ext^{\ext^{\rho''}\tinstU\me}\me''"']
                        && {\hat\sigma} \ar[ul, "\ext^{\tinstU\me''}\ext^{\rho''}\me"] \ar[ulll, bend right=50, "\ext^{\hat\sigma}\me"]
        \\
            & {\sigma'} \ar[ul, -malloctip, "\tinstU\me'"'] \ar[ur, "\ext^{\rho'}\me"]
                && {\sigma_2} \ar[ul, -malloctip, "\ext^{\tinstU\me}\rho'"] \ar[ur, -alloctip, "\ext^{\tinstU\me}\rho''"'] \ar[uu, -alloctip, "\ext^{\tinstU\me}\hat\rho"]
                            && {\sigma''} \ar[ur, -alloctip, "\tinstU\me''"] \ar[ul, "\ext^{\rho''}\me"']
        \\\\
                        &&& {\sigma_1} \ar[uull, -malloctip, "\rho'"'] 
                                       \ar[uu, "\tinstU\me"] 
                                       \ar[uurr, -alloctip, "\rho''"] 
                                       \ar[uuulll, -malloctip, bend left=30, "\hat\rho"] 
                                       \ar[uuurrr, -alloctip, bend right=30, "\hat\rho"']
        \end{tikzcd}
    \end{center}
    where the isomorphisms are given by the pasting property of L.I.C. in $\tempC$.
    The induced data given by \cref{rem:lic-T-data} of the morphisms also commute as they are uniquely determined by $\mvarM_\me, \mvarM_{\me'}, \mvarM_{\me''}$.

    We relate $\inj_{\ext^{\tinstU\me}\rho'}\vA(\ext^{\rho'}\me)(\va')$ with $\inj_{\ext^{\tinstU\me}\rho''}\vA(\ext^{\rho''}\me)(\va'')$ by
    \begin{center}
    \begin{tikzcd}
        & {\sigma_2} \ar[dl, -malloctip, "\ext^{\tinstU\me}\rho'"'] \ar[dr, -alloctip, "\ext^{\tinstU\me}\rho''"]
    \\
    {\sigma_2 \liplus{\sigma_1} \sigma'} \ar[dr, -alloctip, "\iso' \comp \ext^{\ext^{\rho'}\tinstU\me}\me'"']
            && {\sigma_2 \liplus{\sigma_1} \sigma''} \ar[dl, -malloctip, "\iso'' \comp \ext^{\ext^{\rho''}\tinstU\me}\me''"]
    \\ 
        & {\hat\sigma \liplus{\sigma_1} \sigma_2}
    \end{tikzcd}
    \end{center}
    which give us
    \begin{align*}
        \vA(\iso' \comp \ext^{\ext^{\rho'}\tinstU\me}\me' \comp \ext^{\rho'}\me)(\va') 
                &= \vA(\iso' \comp \ext^{\tinstU\me'}\ext^{\rho'}\tinstU\me \comp \tinstU\me')(\va') \\
                &= \vA(\ext^{\hat\sigma}\me \comp \tinstU\me'')(\va'') \\
                &= \vA(\iso'' \comp \ext^{\tinstU\me''}\ext^{\rho''}\tinstU\me \comp \tinstU\me'')(\va'') \\
                &= \vA(\iso'' \comp \ext^{\ext^{\rho''}\tinstU\me}\me'' \comp \ext^{\rho''}\me)(\va'')
    \end{align*}

    To show functoriality, for preservation of identities consider
    \begin{tikzcd}[sep=small]
        & {\sigma} \ar[dl, Rightarrow, "1"'] \ar[dr, "\rho"]
    \\
    {\sigma} \ar[dr, "\ext^{1}\rho"']
            && {\sigma'} \ar[dl, "\ext^\rho1"]
    \\
        & {\sigma \liplus\sigma \sigma'}
    \end{tikzcd}
    and so take
    \begin{tikzcd}[sep=small]
        & {\sigma} \ar[dl, "\ext^{1}\rho"'] \ar[dr, "\rho"]
    \\
    {\sigma \liplus\sigma \sigma'} \ar[dr, Rightarrow, "1"']
            && {\sigma'} \ar[dl, "\ext^\rho1"]
    \\
        & {\sigma \liplus\sigma \sigma'}
    \end{tikzcd}
    and for composition one considers pasting law for L.I.C. 

    That multiplication is quotient respecting, consider
    \begin{tikzcd}
        & {\sigma} \ar[dl, "\rho_1'"'] \ar[dr, "\rho_2'"]
    \\
    {\sigma_1'} \ar[dr, "\tinstU\me_1"']
            && {\sigma_2'} \ar[dl, "\tinstU\me_2"]
    \\
        & {\hat\sigma}
    \end{tikzcd}
    for some $\vq_1 = \quot_{\sigma_1' \xto{\rho_1''} \sigma_1''} \va_1''$, and  
    $\vq_2 = \quot_{\sigma_2' \xto{\rho_2''} \sigma_2''} \va_2''$.
    with
    \[ \thideMon(\vA)(\me_1)(\vq_1) = \thideMon(\vA)(\me_2)(\vq_2) \]

    Unpacking this gives us a diagram
    \begin{center}
        \begin{tikzcd}
                && {\sigma} \ar[dl, "\rho_1'"'] \ar[dr, "\rho_2'"]
        \\
            & {\sigma_1'} \ar[dr, "\tinstU\me_1"] \ar[dl, "\rho_1''"']
                        && {\sigma_2'} \ar[dl, "\tinstU\me_2"'] \ar[dr, "\rho_2''"]
        \\
        {\sigma_1''} \ar[dr, "\ext^{\rho_1''}\me_1"']
                && {\hat\sigma} \ar[dl, "\ext^{\tinstU\me_1}\rho_1''"'] \ar[dr, "\ext^{\tinstU\me_2}\rho_2''"]
                        && {\sigma_2''} \ar[dl, "\ext^{\rho_2''}\me_2"]
        \\ 
            & {\sigma_1'' \liplus{\sigma_1'} \hat\sigma} \ar[dr, "\tinstU\check\me_1"']
                    && {\sigma_2'' \liplus{\sigma_2'} \hat\sigma} \ar[dl, "\tinstU\check\me_2"]
        \\
                && {\check\sigma}
        \end{tikzcd}
    \end{center}
with $\vA(\tinstU\check\me_1 \comp \ext^{\rho_1'}\me_1)(\va_1'') = \vA(\tinstU\check\me_2 \comp \ext^{\rho_2''}\me_2)(\va_2'')$. 

    Which is exactly what we need to show that 
    \[ \hideRes\vA{\rho_1'}\vq_1 = \hideRes\vA{\rho_2'}\vq_2 \]

    That the action on morphisms in $\funC\tinstC\setC$ is quotient respecting and natural is a simple exercise
    and the same goes for naturality of unit and multiplication and that they satisfy the monad laws.
\end{proof}

\begin{lemma}\label{lem:thidemon-hide-ext-com}
    Given a L.I.G.U. diagram
    \begin{tikzcd}[sep=small]
        & {\sigma} 
            \ar[dl, "\me_1"'] \ar[dr, -alloctip, "\me_2"] 
    \\ 
    {\sigma_1} 
        \ar[dr, -alloctip, "\ext^{\me_1}\me_2"']
            && {\sigma_2}
                \ar[dl, "\ext^{\me_2}\me_1"] 
    \\ 
        & {\sigma_1 \liplus\sigma \sigma_2}
    \end{tikzcd} and a $\vr \in \thideMon\vA(\sigma_2)$, we have
    \[ \hideRes\vA{\ext^{\me_1}\me_2}(\thideMon(\ext^{\me_2}\me_1)(\vr)) = \thideMon(\me_1)(\hideRes\vA{\me_2}(\vr)) \] 
\end{lemma}
\begin{proof}
    Let $\vr = \eqclass{\sigma_2 \xato\rho \hat\sigma} \va$, for some $\va : \vA(\hat\sigma)$.
    The LHS of the equation gives
    \[ 
        \eqclass{\sigma_1 \xato{\tinstU(\ext^\me_1\me_2)} \sigma_1 \liplus\sigma \sigma_2 
                          \xato{\ext^{\tinstU(\ext^{\me_2}\me_1)}\rho)} (\sigma_1 \liplus\sigma \sigma_2) \liplus{\sigma_2} \hat\sigma} 
                            \vA(\ext^\rho(\ext^{\me_2}\me_1)(\va))
    \]
    using the diagram
    \begin{center}
        \begin{tikzcd}[sep=small]
            & {\sigma_2} \ar[dr, -alloctip, "\rho"] \ar[dl, "\ext^{\me_2}\me_1"']
        \\
        {\sigma_1 \liplus\sigma \sigma_2} \ar[dr, -alloctip, "\ext^{\tinstU(\ext^{\me_2}\me_1)}\rho"']
                && {\hat\sigma} \ar[dl, "\ext^\rho(\ext^\me_2\me_1)"]
        \\
            & {(\sigma_1 \liplus\sigma \sigma_2) \liplus{\sigma_2} \hat\sigma}
        \end{tikzcd}
    \end{center}

    The RHS of the equation gives
    \[ 
        \eqclass{\sigma_1 \xato{\ext^{\tinstU\me_1}(\rho \comp \tinstU\me_2)} \sigma_1 \liplus\sigma \hat\sigma} \vA(\ext^{(\rho \comp \tinstU\me_2)}\me_1)(\va)
    \]
    using the diagram
    \begin{center}
        \begin{tikzcd}[sep=small]
                && {\sigma} \ar[dr, -alloctip, "\tinstU\me_2"] \ar[ddll, "\me_1"']
        \\
                    &&& {\sigma_2} \ar[dr, -alloctip, "\rho"]
        \\
        {\sigma_1} \ar[ddrr, -alloctip, "\ext^{\tinstU\me_1}(\rho \comp \tinstU\me_2)"']
                        &&&& {\hat\sigma} \ar[ddll, "\ext^{(\rho \comp \tinstU\me_2)}\me_1"]
        \\
        \\
                && {\sigma_1 \liplus\sigma \hat\sigma}
        \end{tikzcd}
    \end{center}

    We relate the two resultlets by
    \begin{center}
        \begin{tikzcd}
            & {\sigma} \ar[dl, "\me_1"'] \ar[dr, -alloctip, "\tinstU\me_2"]
        \\
        {\sigma_1} \ar[dr, -alloctip, "\tinstU(\ext^{\me_1}\me_2)"'] \ar[drr, -alloctip, "\ext^{\tinstU\me_1}(\rho \comp \tinstU\me_2)" near start]
                    && {\sigma_2} \ar[dr, -alloctip, end anchor=north east, "\rho"] \ar[dl, crossing over, "\ext^{\me_2}\me_1"']
        \\
                & {\sigma_1 \liplus\sigma \sigma_2} \ar[dr, -alloctip, "\ext^{\tinstU(\ext^{\me_2}\me_1)}\rho"']
                    & {\sigma_1 \liplus\sigma \hat\sigma} \ar[d, dashed, "\iso"', "\vh"]
                        & {\hat\sigma} \ar[l, "\ext^{(\rho \comp \tinstU\me_2)}\me_1"'] \ar[dl, "\ext^\rho(\ext^{\me_2}\me_1)"]
        \\
                && {(\sigma_1 \liplus\sigma \sigma_2) \liplus{\sigma_2} \hat\sigma}
        \end{tikzcd}
    \end{center}
    where by L.I.C. pasting (\cref{lem:lic-pasting}) the two outer L.I.C. squares are isomorphic,
    witnessed by the iso $\vh$.
    Further, $\vh$ is an allocation morphism by \cref{lem:alloc-cancl}.
    And so, by \cref{lem:thmon-quotient}, the resultlets are related via $\vh$ and the identity $1_{(\sigma_1 \liplus\sigma \sigma_2) \liplus{\sigma_2} \hat\sigma}$.
    (notice that the data of $\ext^\rho(\ext^{\me_2}\me_1)$ and $\ext^{(\rho \comp \tinstU\me_2)}\me_1$ are completely determined by $\mvarM_{\me_1}$ and so they match).
\end{proof}

\subsubsection{Evaluation and Realisation}

An \emph{evaluation} is a collection of resultlets. 

\begin{definition}[Evaluation]
    Let $\Sigma$ be a finite family of templates and $\vA : \tinstC \to \setC$.
    An $\vA$-valued \emph{evaluation} $\ev$ over $\Sigma$ is an indexed family of resultlets $\ev \in \Prod{\sigma_\vi \in \Sigma} \typeRes\vA\sigma_\vi$.
\end{definition}

Let $\vA : \tinstC \to \setC$, total cover $\cov$ over $\pair\vw\Sigma$, $\vA$-valued evaluation $\ev$ over $\Sigma$, 
and let $\phi$ be
\begin{center}
    \begin{tikzcd}
        {\homF\worldC\vw{\instU(\dash)} \times \heapF} \ar[r, "\phi"]
            & {\hideMon{(\vA \comp \inclITI\dash)}} \ar[d, "\hmonIso"] 
        \\
        {\stencilF\vw\inclITI\dash} \ar[r, dashed] \ar[u, "\pair{\anyL\dash}{\mvarM_\dash}", "\iso"']
            & {\thideMon\vA\inclITI\dash}
    \end{tikzcd}
\end{center}
which induces a morphism from $\stencilF\vw\inclITI\dash \dashto \thideMon\vA\inclITI\dash$.

\begin{definition}[Realisation]
We say that $\cov$ and $\ev$ realise $\phi$, $\realise\cov\ev\phi$ if
for every concrete stencil $\me :\, ?_\vw \to 1_{\hat\vw}$, we have
\[ \hmonIso(\phi\,\pair{\anyL\me}{\mvarM_\me}) = \thideMon(\vA)(\tilde\me)(\proj_{\indF{\tilde\me}(\star)} \ev) \]  
where the morphism $\tilde\me : \sigma_{\indF{\tilde\me}(\star)} \to 1_{\hat\vw}$ is given by totality.
\end{definition}


\subsection{Finitary FGLS Monad}
\label{sec:finitary-monad}

A $\phi \in \monFLS\vX\vw$ is finitary when it is generated by its actions on a finite number of \emph{stencils}. 

\begin{definition}[Finitary FGLS Monad]
\[\begin{array}{r@{\;}l@{\,}l@{\,}}
    \monFLSFin         :& \multicolumn{2}{l}{\funC\worldC\setC \to \funC\worldC\setC} \\
    \monFLSFin\vX\vw   \eqdef& \{ \phi \in \monFLS\vX\vw :& \textit{exists a finite family of templates $\Sigma$,} \\
                       &&                                   \textit{a complete cover $\cov$ over $\pair\vw\Sigma$,} \\
                       &&                                   \textit{an $(\vX\anyL\dash \times \theapF)$-valued evaluation $\ev$ over $\Sigma$,} \\
                       &&                                   \textit{such that $\cov,\ev \Vdash \phi$} \}
\end{array}\]
\end{definition}
We give a functor and monad structure to $\monFLSFin$ by showing that the structure of $\monFLS$ lifts through the finitary setting. 

\subsubsection[Action on Morphisms in W]{Action on Morphisms in $\worldC$}

\begin{theorem}
The factorial action of the $\monFLSFin$ monad on morphisms lifts through $\monFLS$, 
that is for every $\sigma : \vw \to \vw'$ we have
\begin{center}
\begin{tikzcd}
    {\monFLSFin\vX\vw}
        \ar[rr, dashed, "\exists\monFLSFin\vX\sigma"] \ar[d, hook]
            && {\monFLSFin\vX\vw'} \ar[d, hook]
    \\
    {\monFLS\vX\vw}
        \ar[rr, "\monFLS\vX\sigma"] 
            && {\monFLS\vX\vw'}
\end{tikzcd}
\end{center}
\end{theorem}
\begin{proof}
Consider a finitary $\phi \in \monFLSFin\vX\vw$, that is, we have a complete cover $\cov$ over $\pair\vw\Sigma$ and an $(\vX\anyL\dash \times \theapF)$-valued evaluation $\ev$ over $\Sigma$ s.t. $\cov,\ev \Vdash \phi$.
The factorial action on $\phi$ is given by precomposition expressed as
\begin{center}
\begin{tikzcd}
    {\homF\worldC{\vw'}{\instU(\dash)} \times \heapF} 
        \ar[rr, "(\dash \comp \sigma) \times \heapF"] \ar[rrr, bend left=15, "\monFLS\vX\sigma(\phi)"]
            && {\homF\worldC\vw{\instU(\dash)} \times \heapF} \ar[r, "\phi"]
                & {\hideMon(\vX\instU \times \heapF)}  \ar[d, "\hmonIso"]
    \\
    {\stencilF{\vw'}\inclITI\dash} \ar[u, "\pair{\anyL\dash}{\mvarM_\dash}"', "\iso"] \ar[rr, "\stencilF\sigma \comp \inclITI\dash"]
                && {\stencilF\vw\inclITI\dash} \ar[u, "\pair{\anyL\dash}{\mvarM_\dash}"', "\iso"] \ar[r, dashed]
                & {\thideMon(\vX\anyL\dash \times \theapF)\inclITI\dash}
\end{tikzcd}
\end{center}

We have to find $\cov^\sigma$, $\ev^\sigma$ s.t. $\cov^\sigma,\ev^\sigma \Vdash \phi \comp (\dash \comp \sigma) \times \heapF$,
So first rebase cover $\cov$ using $\sigma$ (\cref{sec:cov-rebase}),
\begin{center}
\begin{tikzcd}
    {?_\vw}
        \ar[d, "\cov"'] \ar[r, "?_\sigma"]
        & {?_{\vw'}}
            \ar[d, "\cov^\sigma"]
    \\
    {\Sigma} \ar[r, "\rebase^\sigma"' yshift=-2pt]
        & {\Sigma^\sigma} \posquare{ul}
\end{tikzcd}
\end{center}
which gives us a complete cover $\cov^\sigma$.

Define evaluation $\ev^\sigma$ at $\sigma_\vk \in \Sigma^\sigma$,
via $\proj_\vk \comp \rebase^\sigma : \sigma_{\indF{\rebase^\sigma}(\vk)} \to \sigma_\vk$, by 
\[ \proj_\vk \ev^\sigma \mapsto \typeRes{(\vX\anyL\dash \times \theapF)}(\proj_\vk \comp \rebase^\sigma)(\proj_{\indF{\rebase^\sigma}(\vk)} \ev) \]

That $\cov^\sigma,\ev^\sigma \Vdash \phi \comp (\dash \comp \sigma) \times \heapF$, consider $?_{\vw'} \xto\me 1_{\hat\vw}$
\begin{center}
\begin{tikzcd}
    {?_\vw}
        \ar[d, "\cov"'] \ar[r, "?_\sigma"]
        & {?_{\vw'}}
            \ar[d, "\cov^\sigma"] \ar[dr, bend left=30, "\me"]
    \\
    {\Sigma} \ar[r, "\rebase^\sigma"' yshift=-2pt] \ar[rr, dashed, bend right=50, "\widetilde{\me \comp ?_\sigma}"']
        & {\Sigma^\sigma} \ar[r, dashed, "\tilde\me"] \posquare{ul}
            & {1_{\hat\vw}}
\end{tikzcd}
\end{center}
which gives
\begin{align*} 
    \hmonIso(\phi\pair{\anyL{\me \comp ?_\sigma}}{\mvarM_\me})
                    & = \thideMon(\vX\anyL\dash \times \theapF)(\widetilde{\me \comp ?_\sigma})(\proj_{\indF{\widetilde{\me \comp ?_\sigma}}(\star)} \ev) \\
                    & = \thideMon(\vX\anyL\dash \times \theapF)(\tilde\me \comp \rebase^\sigma)(\proj_{\indF{\rebase^\sigma}\indF{\tilde\me}(\star)} \ev) \\
                    & = \thideMon(\vX\anyL\dash \times \theapF)(\tilde\me)(\proj_{\indF{\tilde\me}\star} \ev^\sigma)
\end{align*} 
\end{proof}

\subsubsection[Action on Morphisms in {[W,Set]}]{Action on Morphisms in $\funC\worldC\setC$}
\begin{theorem}
The factorial action of the $\monFLSFin$ monad on morphisms in $\funC\worldC\setC$ lifts through $\monFLS$, 
that is for every $\alpha : \vX \to \vY$ we have
\begin{center}
\begin{tikzcd}
    {\monFLSFin\vX\vw}
        \ar[rr, dashed, "\exists\monFLSFin\alpha_\vw"] \ar[d, hook]
            && {\monFLSFin\vY\vw} \ar[d, hook]
    \\
    {\monFLS\vX\vw}
        \ar[rr, "\monFLS\alpha_\vw"] 
            && {\monFLS\vY\vw}
\end{tikzcd}
\end{center}
\end{theorem}
\begin{proof}
Consider a finitary $\phi \in \monFLSFin\vX\vw$, that is, we have a complete cover $\cov$ over $\pair\vw\Sigma$ and an $(\vX\anyL\dash \times \theapF)$-valued evaluation $\ev$ over $\Sigma$ s.t. $\cov,\ev \Vdash \phi$.

The factorial action on $\phi$ is given by post-composition
\begin{center}
\begin{tikzcd}
    {\homF\worldC\vw{\instU(\dash)} \times \heapF} \ar[r, "\phi"] \ar[rrr, bend left=15, "\monFLS\alpha_\vw(\phi)"]
        & {\hideMon(\vX\instU \times \heapF)} \ar[rr, "\hideMon(\alpha\tinstU \times \heapF)"] \ar[d, "\hmonIso"]
            && {\hideMon(\vY\instU \times \heapF)} \ar[d, "\hmonIso"]
    \\
    {\stencilF\vw\inclITI\dash} \ar[r, dashed] \ar[u, "\pair{\anyL\dash}{\mvarM_\dash}"]
        & {\thideMon(\vX\anyL\dash \times \theapF)\inclITI\dash} \ar[rr, "\thideMon(\alpha\anyL\dash \times \theapF)\comp \inclITI\dash"' yshift=-2pt]
            && {\thideMon(\vY\anyL\dash \times \theapF)\inclITI\dash}
\end{tikzcd}
\end{center}

Define a new evaluation as
\[ \proj_{\sigma_\vi} \ev' \eqdef \typeRes(\alpha\anyL\dash \times \theapF)_{\sigma_\vi}(\proj_{\sigma_\vi} \ev) : \Prod{\sigma_\vi \in \Sigma} \typeRes(\vY\anyL\dash \times \theapF)(\sigma_\vi) \]
Then for $\me :\, ?_\vw \to \inclITI{\hat\vw}$ and induced morphism $\Sigma \xto{\tilde\me} \inclITI{\hat\vw}$.
\begin{align*}
    \hmonIso(\hideMon(\alpha\tinstU \times \heapF)_{\hat\vw}(\phi\pair{\anyL\me}{\mvarM_\me})) & = 
    \thideMon(\alpha\anyL\dash \times \theapF)_{\inclITI{\hat\vw}} (\hmonIso(\phi\pair{\anyL\me}{\mvarM_\me})) \\
        &= \thideMon(\alpha\anyL\dash \times \theapF)_{\inclITI{\hat\vw}}(\thideMon(\vX\anyL\dash \times \theapF)(\tilde\me)(\proj_{\indF{\tilde\me}(\star)} \ev)) \\
        &= \thideMon(\vY\anyL\dash \times \heapF)(\tilde\me)(\thideMon(\alpha\anyL\dash \times \theapF)_{\sigma_\vi}(\proj_{\indF{\tilde\me}(\star)} \ev)) \\
        &= \thideMon(\vY\anyL\dash \times \heapF)(\tilde\me)(\proj_{\indF{\tilde\me}(\star)} \ev') \\
\end{align*}
    and so $\cov, \ev' \Vdash \hideMon(\alpha\tinstU \times \heapF) \comp \phi$.
\end{proof}

\subsubsection{Unit of Monad}

\begin{theorem}
    The unit of the $\monFLSFin$ monad lifts through $\monFLS$, that is we have
\begin{center}
\begin{tikzcd}
    {\vX\vw} \ar[r, dashed, "\eta^\fin_\vw"] \ar[dr, "\eta_\vw"']
        & {\monFLSFin\vX\vw} \ar[d, hook]
    \\
        & {\monFLS\vX\vw}
\end{tikzcd}
\end{center}
\end{theorem}

\begin{proof}
$\eta_\vw(\vx) \in \monFLS\vX\vw$ is given by
\begin{center}
\begin{tikzcd}
    {\homF\worldC\vw{\instU(\dash)} \times \heapF} 
        \ar[rr, "\vX(\dash)(\vx) \times 1"] \ar[rrr, bend left=15, "\eta_\vw(\vx)"]
        && {\vX\instU \times \heapF}
            \ar[r, "\eta^\hideMon"]
            & {\hideMon(\vX\instU \times \heapF)} \ar[d, "\hmonIso"]
    \\
    {\stencilF\vw\inclITI\dash} \ar[u, "\iso"] \ar[rr, "(\vX\anyL\dash(\vx) \times \theapF)\inclITI\dash"]
        && {(\vX\anyL\dash \times \theapF)\inclITI\dash} \ar[r, "\mvarM^{\thideMon} \comp \inclITI\dash"] \ar[u, "\iso"]
            & {\thideMon(\vX\anyL\dash \times \theapF)\inclITI\dash}
\end{tikzcd}
\end{center}

To show that $\eta_\vw(\vx)$ is finitary consider the singleton template $\star \mapsto ?_\vw$,
and the singleton cover over $\pair\vw{\star \mapsto ?_\vw}$, given by the identity stencil $\cov_\star \eqdef\; ?_\vw \xto1 ?_\vw$.
Notice that $\cov$ is an isomorphism, and so it is complete.
Define the $(\vX\anyL\dash \times \theapF)$-valued evaluation at $\star$ to be the resultlet
\[ \proj_\star \ev \eqdef\; \eta_{?_\vw}\pair\vx\emptyset : \typeRes{(\vX\anyL\dash \times \thideMon)}(?_\vw) \]

That $\cov,\ev \Vdash \eta_\vw(\vx)$, consider $?_\vw \xto\me 1_{\hat\vw}$. We have
\begin{align*}
    \hmonIso(\eta_\vw(\vx)\pair{\anyL\me}{\mvarM_\me})
                       & = \eta^\hideMon_{\hat\vw}\pair{\vX\anyL\me(\vx)}{\mvarM_\me} \\
                       & = \eta^\thideMon_{\inclITI{\hat\vw}}\pair{\vX\anyL\me(\vx)}{\mvarM_\me} \\
                       & = \thideMon(\vX\anyL\dash \times \theapF)(\me)(\eta_{?_\vw}\pair\vx\emptyset) \\
                       & = \thideMon(\vX\anyL\dash \times \theapF)(\me)(\proj_\star \ev)
\end{align*}
\end{proof}

\subsubsection{Join of Monad}

\begin{theorem}
    The join of the $\monFLSFin$ monad on morphisms lifts through $\monFLS$.
\begin{center}
\begin{tikzcd}
    {\monFLSFin\monFLSFin\vX\vw}
        \ar[rr, dashed, "\exists \mu^\fin_\vw"] \ar[d, hook]
            && {\monFLSFin\vX\vw} \ar[dd, hook]
    \\
    {\monFLS\monFLSFin\vX\vw}
        \ar[d, "\monFLS\incl_\vw"]
    \\
    {\monFLS\monFLS\vX\vw}
        \ar[rr, "\mu_\vw"] 
            && {\monFLS\vX\vw}
\end{tikzcd}
\end{center}
\end{theorem}
\begin{proof}
    For a finitary $\phi \in \monFLSFin\monFLSFin\vX\vw$, $\mu_\vw(\phi)$ is given by
    \begin{center}
        \begin{tikzcd}
            {\homF\worldC\vw{\instU(\dash)} \times \heapF}
                \ar[r, "\phi"] \ar[rrr, bend left=15, "\mu_\vw(\phi)"]
                & {\hideMon{(\monFLSFin\vX\tinstU \times \heapF)}}
                    \ar[r, "\hideMon(\eval \comp (\vi \times 1))"] \ar[d, "\hmonIso"']
                        & {\Exp\hideMon2{(\vX\tinstU \times \heapF)}}
                            \ar[r, "\mu^\hideMon"]
                            & {\hideMon{(\vX\tinstU \times \heapF)}} \ar[d, "\hmonIso"]
            \\
            {\stencilF\vw\inclITI\dash} \ar[u, "\pair{\anyL\dash}{\mvarM_\dash}", "\iso"'] \ar[r, dashed]
                & {\thideMon(\monFLSFin\vX\anyL\dash \times \theapF)\inclITI\dash} \ar[u, xshift=5pt, "\inv^{\hmonIso}"'] \ar[rr, dashed]
                            && {\thideMon(\vX\anyL\dash \times \theapF)\inclITI\dash}
        \end{tikzcd}
    \end{center}
    which requires
    \begin{enumerate}
        \item A finite family of templates $\Sigma^{\mu\phi}$
        \item A complete cover $\cov^{\mu\phi}$ over $\pair\vw{\Sigma^{\mu\phi}}$
        \item An $(\vX\anyL\dash \times \theapF)$-valued evaluation $\ev^{\mu\phi}$ over $\Sigma^{\mu\phi}$
        \item That $\cov^{\mu\phi},\ev^{\mu\phi} \Vdash \mu_\vw(\phi)$
    \end{enumerate}

    From finitariness of $\phi$, we have a complete cover $\cov$ over $\pair\vw\Sigma$ and $(\monFLSFin\vX\anyL\dash \times \theapF)$-valued evaluation $\ev$ 
    s.t. $\cov, \ev \Vdash \phi$.
    Unpacking this, the resultlet at every $\sigma_\vi \in \Sigma$, $\proj_\vi \ev$ gives
    \begin{itemize}
        \item An allocation morphism $\rho_\vi : \sigma_\vi \ato \sigma_\vi'$.
        \item A finitary computation $\phi_\vi : \monFLSFin\vX\anyL{\sigma_\vi'}$
            \begin{itemize}
                \item A finite family of templates $\Sigma_\vi$
                \item A complete cover $\cov_\vi :\, ?_{\anyL{\sigma_\vi'}} \to \Sigma_\vi$
                \item An evaluation $\ev_\vi : \Prod{\sigma_\vj \in \Sigma_\vi} \typeRes{(\vX\anyL\dash \times \theapF)}(\sigma_\vj)$
            \end{itemize}
        \item A template heap $\mvarM_\vi : \theapF{\sigma_\vi'}$ which is a morphism $?_{\anyL{\sigma_\vi'}} \xto{\mvarM_\vi} \sigma_\vi' \in \tinstC$.
    \end{itemize}
    We lift morphism $\rho_\vi$ to $\tinstC$ by selecting the allocation data using $\dheapF(\concL{\rho_\vi}, 1)(\mvarM_\vi)$,
    and we denote it by $\me_\vi : \sigma_\vi \ato \sigma_\vi' \in \tinstC$.

    \paragraph{Family of templates $\Sigma^{\mu\phi}$}

    We first extract all the ways $\mvarM_\vi$ can unify with cover $\cov_\vi$, that is, we consider the pushout in $\famC\tinstC$
    \begin{center}
        \begin{tikzcd}
            {\Prod{\vi \in \indF\Sigma} ?_{\anyL{\sigma_\vi'}}}
                \ar[rr, "\Prod{\vi \in \indF\Sigma} \cov_\vi"] \ar[dd, "\Prod{\vi \in \indF\Sigma} \mvarM_\vi"'] 
                && {\Prod{\vi \in \indF\Sigma} \Sigma_\vi}
                    \ar[dd, "\Prod{\vi \in \indF\Sigma} \ext^{\cov_\vi}\mvarM_\vi"] 
            \\\\
            {\Prod{\vi \in \indF{\Sigma}} \sigma_\vi'}
                \ar[rr, "\Prod{\vi \in \indF\Sigma}\ext^{\mvarM_\vi}\cov_\vi"']
                && {\Prod{\vi \in \indF\Sigma} \sigma_\vi' \liplus{?_{\anyL{\sigma_\vi'}}} \Sigma_\vi)} \posquare{uull}
        \end{tikzcd}
    \end{center}
    Notice that $\famC\tinstC$ is a coextensive category and thus the pushout above is given as the product of the pushouts at every component.

    Consider a general unifier for a $\mvarM_\vi$, that is consider a $\sigma_\vk \in \sigma_\vi' \liplus{?_{\anyL{\sigma_\vi'}}} \Sigma_\vi$,
    and let $\vj \eqdef \indF{\ext^{\cov_\vi}\mvarM_\vi}(\vk)$ be the index of the stencil in $\cov_\vi$ that matches.
    We have a diagram (connectivity follows from \cref{lem:general-unifier-con-prop})
    \begin{center}
    \begin{tikzcd}
        & {?_{\anyL{\sigma_\vi'}}} \ar[dl, twoheadrightarrow, "\mvarM_\vi"'] \ar[dr, twoheadrightarrow, "\proj_\vj \comp \cov_\vi"]
    \\
    {\sigma_\vi'} \ar[dr, twoheadrightarrow, "\me^\vk_1 \eqdef \proj_\vk \comp \ext^{\mvarM_\vi}\cov_\vi"']
            && {\sigma_\vj} \ar[dl, twoheadrightarrow, "\proj_\vk \comp \ext^{\cov_\vi}\mvarM_\vi \defeq \me^\vk_2"]
    \\
        & {\sigma_\vk}
    \end{tikzcd}
    \end{center}

    Intuitively, $\me^\vk_1$ instantiates the variables $\varL{\sigma_\vi}$ with enough data to match the stencil $\proj_\vj \comp \cov_\vi$.
    But we need to be careful, variables $\varL{\sigma_\vi}$ represent the extensions on the starting heap before the allocation occurs.
    As a result, we need to select only the $\me^\vk_1$ where the instantiated data do not point into the newly allocated data $\mvarM_{\me_\vi}$.
    Semantically, this is expressed as $\me^\vk_1$ should preserve the disconnectivity of $\me_\vi$ (\cref{def:dcon-pres-morph}),
    which is by \cref{lem:dcon-ligu,lem:ligu-dcon} it is equivalent to requiring that
\begin{center}
    \begin{tikzcd}
        & {\sigma_\vi} 
            \ar[dl, twoheadrightarrow, "\con{\me^\vk_1 \comp \me_\vi}"'] \ar[dr, -alloctip, "\me_\vi"]
    \\
    {\sigma_{\con{\me^\vk_1 \comp \me_\vi}}}
        \ar[dr, -rmonotip, "\uncon{\me^\vk_1 \comp \me_\vi}"']
            && {\sigma_\vi'}
                \ar[dl, twoheadrightarrow, "\me^\vk_1"]
    \\
        & {\sigma^\vk_\vi}
    \end{tikzcd}
\end{center}
    is a L.I.G.U.
    When this is the case, the reachable component $\sigma_{\con{\me^\vk_1 \comp \me_\vi}}$ represents only instantiations of the variables $\varL{\sigma_\vi}$
    with data that do not point to the allocated locations.

    Let $\dpres_{\me_\vi}(\sigma_\vi' \pltimes{?_{\anyL{\sigma_\vi'}}} \Sigma_\vi) \subset \sigma_\vi' \pltimes{?_{\anyL{\sigma_\vi'}}} \Sigma_\vi)$
    represent the selection of general unifiers from $\sigma_\vk \in \sigma_\vi' \pltimes{?_{\anyL{\sigma_\vi'}}} \Sigma_\vi$ such that $\me^\vk_1$ preserves the disconnectivity of $\me_\vi$ as described.

    \begin{align*}
        \Sigma 
            \xto{\overbrace{\Prod{\sigma_\vi \in \Sigma} \ipair\vk{\con{\me^\vk_1 \comp \me_\vi}}}^\embed}
        \overbrace{
            \Prod{\sigma_\vi \in \Sigma} \Prod{\sigma_\vk \in \dpres_{\me_\vi}(\sigma_\vi' \pltimes{?_{\anyL{\sigma_\vi'}}} \Sigma_\vi)}\sigma_{\con{\me^\vk_1 \comp \me_\vi}}
        }^{\Sigma^{\mu\phi}} 
    \end{align*}

\paragraph{Cover $\cov^{\mu\phi}$}
    It is enough to show that morphism $\embed : \Sigma \to \Sigma^{\mu\phi}$ is complete, 
    and take the composition $\cov^{\mu\phi} \eqdef\, ?_\vw \xto\cov \Sigma \xto{\embed} \Sigma^{\mu\phi}$ as our cover (complete morphisms are closed under composition).

    That $\embed$ is connected is immediate, so we need to show that is non-overlapping and total.
    For non-overlappingness, it is enough to show that all diagrams of this form satisfy $\vk = \vl$.
    \begin{center}
    \begin{tikzcd}
        & {\sigma_\vi} \ar[dl, "{\con{\me^\vk_1 \comp \me_\vi}}"'] \ar[dr, "\sigma_{\con{\me^\vl_1 \comp \me_1}}"] \ar[dd, "\hat\me"]
    \\
    {\sigma_{\con{\me^\vk_1 \comp \me_1}}} \ar[dr, "\hat\me_\vk"']
            && {\sigma_{\con{\me^\vl_1 \comp \me_\vi}}} \ar[dl, "\hat\me_\vl"]
    \\
        & {\hat\sigma}
    \end{tikzcd}
    \end{center}

    Notice that all the inner squares are L.I.G.U, and by \cref{lem:ligu-pasting,lem:ligu-uniq,lem:ligu-rmono-exist} we have the given isos
    \begin{center}
    \begin{tikzcd}
            && {\sigma_\vi} \ar[dl, twoheadrightarrow, "\con{\me^\vk_1 \comp \me_\vi}"'] \ar[ddll, bend right=40, "\hat\me"'] \ar[ddrr, bend left=40, "\hat\me"] \ar[dr, twoheadrightarrow, "\con{\me^\vl_1 \comp \me_\vi}"] \ar[dd, -alloctip, "\me_\vi"]
    \\
        & {\sigma_{\con{\me^\vk_1 \comp \me_\vi}}} \ar[dl, "\hat\me_\vk"'] \ar[dd, -alloctip, "\uncon{\me^\vk_1 \comp \me_\vi}"']
                && {\sigma_{\con{\me^\vl_1 \comp \me_\vi}}} \ar[dr, "\hat\me_\vl"] \ar[dd, -alloctip, "\uncon{\me^\vl_1 \comp \me_\vi}"]
    \\
    {\hat\sigma} \ar[dd, -alloctip] \ar[dddrr, -alloctip, "\ext^{\hat\me}\me_\vi" near start]
            && {\sigma_\vi'} \ar[ddd, "\ext^{\me_\vi}\hat\me"] \ar[dl, twoheadrightarrow, "\me^\vk_1"'] \ar[dr, twoheadrightarrow, "\me^\vl_1"]
                    && {\hat\sigma} \ar[dd, -alloctip] \ar[dddll, -malloctip, "\ext^{\hat\me}\me_\vi"' near start]
    \\
        & {\sigma^\vk} \ar[dl, crossing over]
                && {\sigma^\vl} \ar[dr, crossing over]
    \\ 
    {\hat\sigma \liplus{\sigma_{\con{\me^\vk_1 \comp \me_\vi}}} \sigma^\vk} \ar[drr, crossing over, "\iso"]
                    &&&& {\hat\sigma \liplus{\sigma_{\con{\me^\vl_1 \comp \me_\vi}}} \sigma^\vl} \ar[dll, crossing over, "\iso"']
    \\
            && {\hat\sigma \liplus{\sigma_\vi} \sigma_\vi'}
    \end{tikzcd}
    \end{center}

    Let $\vj \eqdef \indF{\ext^{\cov_\vi}\mvarM_\vi}(\vk)$ and $\vj' \eqdef \indF{\ext^{\cov_\vi}\mvarM_\vi}(\vl)$ be the indices of their matching stencils.
    We have a diagram
    \begin{center}
    \begin{tikzcd}
        & {?_{\anyL{\sigma_\vi'}}} \ar[dl, "\mvarM_\vi"'] \ar[dr, "\proj_\vj \cov_\vi"'] \ar[drr, "\proj_{\vj'} \cov_\vi"]
    \\
    {\sigma_\vi'} \ar[dr, "\me^\vk_1"'] \ar[drr, "\me^\vl_1" near start]
            && {\sigma_\vj} \ar[dl, crossing over, "\me^\vk_2"' near start]
                & {\sigma_{\vj'}} \ar[dl, "\me^\vl_2"]
    \\
        & {\sigma_\vk} \ar[dr]
            & {\sigma_\vl} \ar[d]
    \\
            && {\hat\sigma \liplus{\sigma_\vi} \sigma_\vi'}
    \end{tikzcd}
    \end{center}
and by $\cov_\vi$ being complete, we have that $\vj = \vj'$, and by \cref{lem:general-unifier-prop}, the general unifiers $\vk, \vl$ must be iso which means they are equal in our setting since $\sigma_\vi' \liplus{?_{\anyL{\sigma_\vi'}}} \Sigma_\vi$ picks representatives for each isomorphism class of general unifiers.

    For totality of $\embed$, we need to show that $\tilde\me$ exists for any $\me$ 
    \begin{tikzcd}
        {\Sigma} \ar[rr, "\me"] \ar[dr, "\embed"']
                && {1_{\hat\vw}}
        \\
            & {\Sigma^{\mu\phi}} \ar[ur, "\exists \tilde\me"']
    \end{tikzcd}

    For $\vi \eqdef \indF\me(\star)$, consider the L.I.G.U. given by \cref{lem:ligu-rmono-exist}
    \begin{tikzcd}
        & {\sigma_\vi} \ar[dl, "\me"'] \ar[dr, -alloctip, "\me_\vi"]
    \\
    {1_{\hat\vw}} \ar[dr, -alloctip, "\ext^\me\me_\vi"']
            && {\sigma_\vi'} \ar[dl, "\ext^{\me_\vi}\me"]
    \\
        & {1_{\hat\vw} \liplus{\sigma_\vi} \sigma_\vi'}
    \end{tikzcd}
    and notice that $1_{\hat\vw} \liplus{\sigma_\vi} \sigma_\vi'$ is a concrete template since it has a variable preserving morphism into it with domain a concrete template.
   
    By totality of $\cov_\vi$, and further by \cref{lem:general-unifier-prop} we have
    \begin{tikzcd} 
        & {?_{\anyL{\sigma_\vi'}}} \ar[dl, "\mvarM_\vi"'] \ar[dr, "\proj_\vj \cov_\vi"] 
    \\
    {\sigma_\vi'} \ar[ddr, "\ext^{\me_\vi}\me"'] \ar[dr, "\me^\vk_1"] 
            && {\sigma_\vj} \ar[dl, "\me^\vk_2"'] \ar[ddl]
    \\
        & {\sigma^\vk} \ar[d, dashed, "\vh"]
    \\
        & {1_{\hat\vw} \liplus{\sigma_\vi} \sigma_\vi}
    \end{tikzcd}
    where $\sigma^\vk$ is its general unifier.
     
    To show that general unifier $\sigma_\vk$ with $\me^\vk_1, \me^\vk_2$ is in $\Sigma^{\mu\phi}$ we only need to show that  
    $\me^\vk_1$ preserves the disconnectivity of $\me_\vi$ which follows from \cref{lem:ligu-dcon,lem:dcon-pres-dclosed}
    \begin{center}
    \begin{tikzcd}
        {\sigma_\vi} \ar[r, -rmonotip, "\me_\vi"] 
            & {\sigma_\vi'} \ar[rr, bend right=20, "\ext^{\me_\vi}\me"'] \ar[r, "\me^\vk_1"] 
                & {\sigma^\vk} \ar[r, dashed, "\vh"] 
                    & {1_{\hat\vw} \liplus{\sigma_\vi} \sigma_\vi'}
    \end{tikzcd}
    \end{center}
    As a result, let the morphism $\tilde\me$ be given by general unifier $\vk$ and the induced morphism by the factorisation system
    \begin{tikzcd}
        {\sigma_\vi} \ar[r, "\me"] \ar[dr, twoheadrightarrow, "\con{\me^\vk_1 \comp \me_\vi}"']
            & {1_{\hat\vw}}
                \ar[r, -alloctip, "\ext^\me\me_\vi"]
                & {1_{\hat\vw} \liplus{\sigma_\vi} \sigma_\vi'}
        \\
            & {\sigma_{\con{\me^\vk_1 \comp \me_\vi}}}
                \ar[u, dashed, "\tilde\me"'] \ar[r, -rmonotip, "\uncon{\me^\vk_1 \comp \me_\vi}"']
                & {\sigma^\vk}
                    \ar[u, "\vh"']
    \end{tikzcd}
      
    \paragraph{Evaluation $\ev^{\mu\phi}$}
    
    We need an evaluation 
    \[ \ev^{\mu\phi} : \Prod{\sigma_{\con{\me^\vk_1 \comp \me_\vi}} \in \Sigma^{\mu\phi}} \typeRes(\vX\anyL\dash \times \theapF)(\sigma_{\con{\me^\vk_1 \comp \me_\vi}}) \]
    We define the resultlet at a $\sigma_{\con{\me^\vk_1 \comp \me_\vi}}$. 
    Let $\vj = \indF{\ext^{\cov_\vi}\mvarM_\vi}(\vk)$ be the index of its matching stencil, and from the diagram
    \begin{tikzcd}
        & {?_{\anyL{\sigma_\vi'}}} \ar[dl, "\mvarM_\vi"'] \ar[dr, "\proj_\vj \cov_\vi"]
    \\
    {\sigma_\vi'} \ar[dr, "\me^\vk_1"']
            && {\sigma_\vj} \ar[dl, "\me^\vk_2"]
    \\
        & {\sigma^\vk}
    \end{tikzcd}
    define the resultlet as
    \[ \proj_{\sigma_{\con{\me^\vk_1 \comp \me_\vi}}} \ev^{\mu\phi} \eqdef 
            \hideRes\relax{\uncon{\me^\vk_1 \comp \me_\vi}} (\typeRes(\vX\anyL\dash \times \theapF)(\me^\vk_2)(\proj_\vj \ev_\vi)) \]

    \paragraph{Realisation $\cov^{\mu\phi},\ev^{\mu\phi} \Vdash \mu_\vw(\phi)$}
    For a $\me : \, ?_\vw \to 1_{\hat\vw}$, we have by totality of $\cov^{\mu\phi}$
    \begin{tikzcd}
        {?_\vw} \ar[r, "\me"] \ar[d, "\cov"']
            & {1_{\hat\vw}}
        \\
        {\Sigma} \ar[ur, "\tilde\me"] \ar[r, "\embed"']
            & {\Sigma^{\mu\phi}} \ar[u, "\me^{\mu\phi}"']
    \end{tikzcd}
    which gives
    \begin{itemize}
        \item An $\vi \eqdef \indF{\tilde\me}(\star)$, and so a $\sigma_\vi \in \Sigma$, and $\tilde\me : \sigma_\vi \to 1_{\hat\vw} \in \tinstC$
        \item A $\vk \eqdef \indF{\me^{\mu\phi}}(\star)$, and a $\me^{\mu\phi} : \sigma_{\con{\me^\vk_1 \comp \me_\vi}} \to \inclITI{\hat\vw}$ with $\tilde\me = \me^{\mu\phi} \comp \embed_\vk$ given by (let $\vj \eqdef \indF{\ext^{\cov_\vi}\mvarM_\vi}(\vk)$)

            \begin{center}
                \begin{tikzcd}
                    & {?_{\anyL{\sigma_\vi'}}} \ar[dl, twoheadrightarrow, "\mvarM_\vi"'] \ar[dr, twoheadrightarrow, "\proj_\vj \, \cov_\vi"]
                \\
                {\sigma_\vi'} \ar[dr, twoheadrightarrow, "\me^\vk_1 \eqdef \proj_\vk \, \ext^{\mvarM_\vi}\cov_\vi"']
                        && {\sigma_\vj} \ar[dl, twoheadrightarrow, "\proj_\vk \, \ext^{\cov_\vi}\mvarM_\vi \defeq \me^\vk_2"]
                \\
                    & {\sigma_\vk}
                \end{tikzcd}
                \begin{tikzcd}
                  {\sigma_\vi} \ar[d, twoheadrightarrow, "\embed_\vk = \con{\me^\vk_1 \comp \me_\vi}"'] \ar[r, -alloctip, "\me_\vi"] \pbsquare{dr}
                      & {\sigma_\vi'} \ar[d, twoheadrightarrow, "\me^\vk_1"]
                  \\
                  {\proj_\vk \Sigma^{\mu\phi} \eqdef \sigma_{\con{\me^\vk_1 \comp \me_\vi}}}
                    \ar[r, -alloctip, "\uncon{\me^\vk_1 \comp \me_\vi}"']
                      & {\sigma^\vk_\vi}
                \end{tikzcd}
            \end{center}
            where the RHS square is a L.I.G.U. by \cref{lem:dcon-ligu}.

        \item Notice that the resultlet $\proj_\vi \ev = \inj_{\sigma_\vi \xato{\rho_\vi} \sigma_\vi'} \pair{\phi_\vi}{\mvarM_\vi}$ extends to
            \[ 
                \typeRes(\monFLSFin\vX\anyL\dash \times \theapF)(\tilde\me)(\proj_\vi \ev) = 
                    \inj_{\inclITI{\hat\vw} \xato{\ext^{\tilde\me}\me_\vi} {\inclITI{\hat\vw} \liplus{\sigma_\vi} \sigma_\vi'}}
                        \pair{\monFLSFin\vX\anyL{\ext^{\me_\vi}{\tilde\me}}(\phi_\vi)}{\theapF(\ext^{\me_\vi}{\tilde\me})(\mvarM_\vi)}
            \]
            using the L.I.G.U. (\cref{lem:ligu-rmono-exist})

            \begin{tikzcd}[sep=small]
                & {\sigma_\vi} \ar[dr, -alloctip, "\me_\vi"]  \ar[dl, "\tilde\me"']
            \\
            {\inclITI{\hat\vw}} \ar[dr, -alloctip, "\ext^{\tilde\me}\me_\vi"']
                    && {\sigma_\vi'} \ar[dl, "\ext^{\me_\vi}{\tilde\me}"]
            \\
                & {\inclITI{\hat\vw} \liplus{\sigma_\vi} \sigma_\vi'}
            \end{tikzcd}
            which further factors as (the outer square is a L.I.S. by \cref{lem:pullback-factoring}, so $\vh$ is induced by \cref{lem:ligu-prop})
            \begin{center}
            \begin{tikzcd}
                {\sigma_\vi} \ar[r, twoheadrightarrow, "\con{\me^\vk_1 \comp \me_\vi}"'] \ar[rr, bend left=25, "\tilde\me"] \ar[d, -alloctip, "\me_\vi"']
                    & {\sigma_{\con{\me^\vk_1 \comp \me_\vi}}} \ar[d, -alloctip, "\uncon{\me^\vk_1 \comp \me_\vi}"] \ar[r, "\me^{\mu\phi}"]
                        & {\inclITI{\hat\vw}} \ar[d, -alloctip, "\ext^{\tilde\me}\me_\vi"]
                \\
                {\sigma_\vi'} \ar[r, "\me^\vk_1"] \ar[rr, bend right=15, "\ext^{\me_\vi}\tilde\me"']
                    & {\sigma_\vk} \ar[r, dashed, "\vh"]
                        & {\sigma_\vi' \liplus{\sigma_\vi} \inclITI{\hat\vw}}
            \end{tikzcd}
            \end{center}
            and notice that the RHS square is itself a L.I.G.U. by \cref{lem:ligu-pasting}.

            Completeness of $\cov_\vi$, gives us
            \[ \hmonIso(\phi_\vi\pair{\anyL{\ext^{\me_\vi}{\tilde\me}}}{\theapF(\ext^{\me_\vi}{\tilde\me})(\mvarM_\vi)}) =
                        \thideMon(\vX\anyL\dash \times \theapF)(\vh)(\typeRes(\vX\anyL\dash \times \theapF)(\me^\vk_2)(\proj_\vj \ev_\vi)) \]
    \end{itemize}

    Lastly we prove $\cov^{\mu\phi},\ev^{\mu\phi} \Vdash \mu_\vw(\phi)$ by
    \begin{align*}
        \hmonIso(\mu_\vw(\phi)\pair{\anyL\me}{\mvarM_\me}) 
            &= \hmonIso(\mu^\hideMon(\hideMon(\eval \comp (\incl \times 1))(\phi\pair{\anyL\me}{\mvarM_\me})))          & \tag{$\cov,\ev \Vdash \phi$} 
\\
            &= \hmonIso(\mu^\hideMon(\hideMon(\eval \comp (\incl \times 1))(
                        \inv^\hmonIso(\thideMon(\monFLSFin\vX\anyL\dash \times \theapF)(\tilde\me)(\proj_\vi \ev)))))
\\
            &= \hmonIso(\mu^\hideMon(\hideMon(\eval \comp (\incl \times 1))(
                    \inv^\hmonIso(\eqclass{\inclITI{\hat\vw} \xato{\ext^{\tilde\me}\me_\vi} {\inclITI{\hat\vw} \liplus{\sigma_\vi} \sigma_\vi'}}
                                          \pair{\monFLSFin\vX\anyL{\ext^{\me_\vi}{\tilde\me}}(\phi_\vi)}{\theapF(\ext^{\me_\vi}{\tilde\me})(\mvarM_\vi)})))) 
\\
            &= \hmonIso(\mu^\hideMon(\eqclass{\inclITI{\hat\vw} \xato{\ext^{\tilde\me}\me_\vi} {\inclITI{\hat\vw} \liplus{\sigma_\vi} \sigma_\vi'}}
                                             \phi_\vi\pair{\anyL{\ext^{\me_\vi}{\tilde\me}}}{\theapF(\ext^{\me_\vi}{\tilde\me})(\mvarM_\vi)})) 
\\
            &= \hmonIso(\hideRes\relax{\ext^{\tilde\me}\me_\vi}\phi_\vi\pair{\anyL{\ext^{\me_\vi}{\tilde\me}}}{\theapF(\ext^{\me_\vi}{\tilde\me})(\mvarM_\vi)})
\\
            &= \hideRes\relax{\ext^{\tilde\me}\me_\vi}(
                    \thideMon(\vX\anyL\dash \times \theapF)(\vh)(\typeRes(\vX\anyL\dash \times \theapF)(\me^\vk_2)(\proj_\vj \ev_\vi)))
                                                                                                                        & \tag{by \cref{lem:thidemon-hide-ext-com}}
\\ 
            &= \thideMon(\vX\anyL\dash \times \theapF)(\me^{\mu\phi})(
                    \hideRes\relax{\uncon{\me^\vk_1 \comp \me_\vi}} (\typeRes(\vX\anyL\dash \times \theapF)(\me^\vk_2)(\proj_\vj \ev_\vi)))
                                                                                                                        & \tag{def $\ev^{\mu\phi}$}
\\ 
            &= \thideMon(\vX\anyL\dash \times \theapF)(\me^{\mu\phi})(\proj_\vk \ev^{\mu\phi})
    \end{align*}
\end{proof}


\appendix

\section{Appendix}


\subsection{Pullback Lemmas}

\begin{lemma}[Pullbacks and mono post-composition]\label{lem:pullback-mono-post-comp}
    Let $\vk$ be mono. Then every square of the form 
    \begin{arrow-diagram}
        {} \ar[d, "\pi_2"'] \ar[r, "\pi_1"] & {} \ar[d, "\vf"] \\
        {} \ar[r, "\vg"'] & {}
    \end{arrow-diagram}
    is a pullback square iff
    \begin{arrow-diagram}
        {} \ar[dd, "\pi_2"'] \ar[rr, "\pi_1"] && {} \ar[d, "\vf"] \\
        {}                                    && {} \ar[d, "\vk"] \\
        {} \ar[r, "\vg"'] & {} \ar[r, "\vk"']  &  {}
    \end{arrow-diagram}
    is a pullback square.
\end{lemma}
\begin{proof} \nl
    $\to$) Take $\vk \comp \vf \comp \kappa_1 = \vk \comp \vg \comp \kappa_2$.
           From $\vk$ being mono we have $\vf \comp \kappa_1 = \vg \comp \kappa_2$. 
           The induced morphism from the pullback on the left is exactly the one we need.
    \nl
    $\from$) Take $\vf \comp \kappa_1 = \vg \comp \kappa_2$ which gives us $\vk \comp \vf \comp \kappa_1 = \vk \comp \vg \comp \kappa_2$.
              The induced morphism from the pullback on the right is exactly the one we need.
\end{proof}

\begin{lemma}[Pullbacks and isos]\label{lem:pullback-iso-mono}
    A square
    \begin{arrow-diagram}
        {} \ar[d, "\pi_2"'] \ar[r, "\pi_1"] & {} \ar[d, "\vf"] \\
        {} \ar[r, "\vg"'] & {}
    \end{arrow-diagram}
    where $\pi_2$ is iso and $\vf$ is mono, is a pullback square.
\end{lemma}
\begin{proof}
    For $\vg \comp \vk = \vf \comp \vl$, define the unique morphism $\vh = \inv^{\pi_2} \comp \vk$.
    Then $\vk = \pi_2 \comp \inv^{\pi_2} \comp \vk$ and $\vl = \pi_1 \comp \inv^{\pi_2} \comp \vk$ by post-composition with monomorphism $\vf$.
    $\vh$ is unique by $\pi_2$ being mono.
\end{proof}

\begin{lemma}[Pullbacks with Initial Objects]\label{lem:pullback-initial-obj}
    Given a category with a strict initial object and a diagram
    \begin{tikzcd}
        {0} \ar[r, "?_\vA"] \ar[d, "?_\vB"']
            & {A} \ar[r, "\vh"] \ar[d, "\vf"]
                & {D} \ar[d, "\vk"]
        \\
        {B} \ar[r, "\vg"]
            & {C} \ar[r, "\vl"]
                & {E}
    \end{tikzcd}
    where the outer square is a pullback, we have that
    \begin{tikzcd}
        {0} \ar[r, "?_\vA"] \ar[d, "?_\vB"']
            & {A} \ar[d, "\vf"]
        \\
        {B} \ar[r, "\vg"]
            & {C}
    \end{tikzcd} is also a pullback.
\end{lemma}
\begin{proof}
    Consider $\vf \comp \phi = \vg \comp \xi$. 
    Then $\vk \comp \vh \comp \phi = \vl \comp \vf \comp \phi = \vl \comp \vg \comp \xi$ and so we have a morphism into $0$ which by stictness of the initial object, $\phi$, and $\xi$ must be the initial morphisms.
\end{proof}

\begin{lemma}[Pullback preserved along factoring]\label{lem:pullback-factoring}
    Given a pullback
    \begin{tikzcd}
        {\vA \pltimes\vC \vB} \ar[r, "\proj_2"] \ar[d, "\proj_1"] \pbsquare{dr}
            & {\vB} \ar[d, "\vg"]
        \\
        {\vA} \ar[r, "\vf"]
            & {\vC}
    \end{tikzcd}
    and a factoring of $\proj_2$ as 
    \begin{tikzcd}
            & {\vB'} \ar[d, "\vl"]
        \\
        {\vA \pltimes\vC \vB} \ar[ur, "\vk"] \ar[r, "\proj_2"]
            & {\vB}
    \end{tikzcd} where $\vl$ is mono, then 
    \begin{tikzcd}
        {\vA \pltimes\vC \vB} \ar[r, "\vk"] \ar[d, "\proj_1"]
            & {\vB'} \ar[d, "\vg \comp \vl"]
        \\
        {\vA} \ar[r, "\vf"]
            & {\vC}
    \end{tikzcd} is a pullback.
\end{lemma}
\begin{proof}
    Consider $\vg \comp \vl \comp \vp_2 = \vf \comp \vp_1$. 
    We get a unique morphism $\vh$ s.t. $\vp_1 = \proj_1 \comp \vh$ and $\vl \comp \vk \comp \vh = \vl \comp \vp_2$
    which by $\vl$ being mono, we have $\vk \comp \vh = \vp_2$ and uniqueness follows from the pullback.
\end{proof}

\begin{lemma}[Monos and Coparallel map in extensive categories]\label{lem:monos-coparallel-extensive}
    Given an extensive category $\mcatC$, and morphisms $\vf : \vA \to \vB, \vg : \vA' \to \vB' \in \mcatC$, 
    we have that the coparallel map $\vf + \vg : \vA + \vA' \to \vB + \vB'$ is mono iff $\vf,\vG$ are.
\end{lemma}
\begin{proof}
    For $\to$) direction, extensiveness gives us the pullbacks
    \begin{center}
        \begin{tikzcd}
            {\vA} \ar[r, "\inj_\vA"] \ar[d, "\vf"] \pbsquare{dr}
                & {\vA + \vA'}
                    \ar[d, dashed, "\vf + \vg"]
                    & {\vA'} 
                        \ar[l, "\inj_{\vA'}"'] \pbsquare{dl} \ar[d, "\vg"]
            \\
            {\vB} \ar[r, "\inj_\vB"]
                & {\vB + \vB'}
                    & {\vB'} 
                        \ar[l, "\inj_{\vB'}"]
        \end{tikzcd}
    \end{center}
    which imply $\vf$ and $\vg$ are mono.

    For $\from$, extensiveness gives us that for every $\vA'$, $\vB'$, we have an equivalence of categories
    \[ \sliceC\mcatC{\vA'} \times \sliceC\mcatC{\vB'} \equivC \sliceC\mcatC{\vA' + \vB'} \]
    given by the coparallel map.
    Since equivalences preserve monos, the coparallel map will also be mono.
\end{proof}


\subsection{Yoneda Lemmas}

\begin{lemma}[Restricted Yoneda Embedding preserves Coequalizers]\label{lem:ryoneda-coeqz}
    The yoneda embedding $\yoneda : \mcatC \to \fsreprC\mcatC$ preserves coequalizers.
\end{lemma}
\begin{proof}
    Given a coequalizing diagram 
    \begin{tikzcd}{\vA} \ar[r, yshift=3pt, "\vf"] \ar[r, yshift=-3pt, "\vg"'] & {\vB} \ar[r, "\ve"] & {\vC}\end{tikzcd}
    we show that 
    \begin{tikzcd}{\yoneda\vA} \ar[r, yshift=3pt, "\yoneda\vf"] \ar[r, yshift=-3pt, "\yoneda\vg"'] & {\yoneda\vB} \ar[r, "\yoneda\ve"] & {\yoneda\vC}\end{tikzcd}
    is also coequalizing.

    Consider a $\yoneda\vB \xto\alpha \Gamma$ s.t. $\alpha \comp \yoneda\vf = \alpha \comp \yoneda\vg$.
    By yoneda, $\alpha : \Gamma\vB$ is a morphism $\vB \xto{\check\alpha} \vC'$ for some $\vC' \in \Gamma$ s.t. $\check\alpha \comp \vf = \check\alpha \comp \vg$
    and so we have a unique $\vC \xto\vh \vC'$, which by yoneda a morphism  $\yoneda\vC \to \Gamma$.
    Uniqueness follows directly from uniqueness in the base category.
\end{proof}

\begin{lemma}[Restricted Yoneda Embedding preserves Pushouts]\label{lem:ryoneda-pushouts}
    The yoneda embedding $\yoneda : \mcatC \to \fsreprC\mcatC$ preserves pushouts.
\end{lemma}
\begin{proof}
    Given a pushout diagram 
    \begin{tikzcd}
        {\vA} \ar[r, "\vg"] \ar[d, "\vf"'] 
            & {\vC} \ar[d, "\inj_2"] 
        \\
        {\vB} \ar[r, "\inj_1"']
            & {\vB \liplus\vA \vC} \posquare{ul}
    \end{tikzcd}
    we show that 
    \begin{tikzcd}
        {\yoneda\vA} \ar[r, "\yoneda\vg"] \ar[d, "\yoneda\vf"'] 
            & {\yoneda\vC} \ar[d, "\yoneda(\inj_2)"] 
        \\
        {\yoneda\vB} \ar[r, "\yoneda(\inj_1)"']
            & {\yoneda(\vB \liplus\vA \vC)}
    \end{tikzcd}
    is also a pushout.

    Consider a $\yoneda\vB \xto\alpha \Gamma, \yoneda\vC \xto\beta \Gamma$ s.t. $\alpha \comp \yoneda\vf = \beta \comp \yoneda\vg$.
    By yoneda, we have morphisms $\vB \xto{\check\alpha} \vD$ and $\vC \xto{\check\beta} \vD$ 
    s.t. $\check\alpha \comp \vf = \check\beta \comp \vg$,
    and so we have a unique $\vB \liplus\vA \vC \xto\vh \vD$, which by yoneda is a morphism $\yoneda(\vB \liplus\vA \vC) \to \Gamma$.
    Uniqueness follows directly from uniqueness in the base category.
\end{proof}

\begin{lemma}\label{lem:fsreprC-mono}
    A $\alpha : \Gamma_1 \to \Gamma_2 \in \fsreprC\mcatC$ is mono iff it is component-wise injective.
\end{lemma}
\begin{proof}
    The $\from$ direction is clear.
    For $\to$, consider $\inj_\vi(\vf_\vi : \vC \to \vC_\vi) \neq \inj_{\vj}(\vf_\vj : \vC \to \vC_\vj)$.
    We want to show that
    \[ \alpha(\inj_\vi(\vf_\vi)) = \Gamma_2(\vf_\vi)(\alpha(\inj_\vi1_{\vC_\vi})) \neq \Gamma_2(\vf_\vj)(\alpha(\inj_{\vj}1_{\vC_\vj})) =  \alpha(\inj_\vj(\vf_\vj)) \] 
    so assume $\Gamma_2(\vf_\vi)(\alpha(\inj_\vi1)) = \Gamma_2(\vf_\vj)(\alpha(\inj_\vj1))$
    and we show a contradiction.

    We have that $\indF\alpha(\vi) = \indF\alpha(\vj)$ and
    \begin{center}
        \begin{tikzcd}
            & {\vC}
                \arrow[dl, "\vf_\vi"'] \arrow[dr, "\vf_\vj"]
        \\
            {\vC_\vi}
                \arrow[dr, "\alpha_\vi"']
                    && {\vC_\vj}
                        \arrow[dl, "\alpha_\vj"]
        \\ 
            & {\vC_{\alpha(\vi)}}
        \end{tikzcd}
    \end{center}

    We have morphisms $\inj_\vi(\vf_\vi),\, \inj_\vj(\vf_\vj) : \yoneda\vC \to \Gamma_1$.
    And by copairing we have morphisms $\copair1{\inj_\vi(\vf_\vi)}, \copair1{\inj_\vj(\vf_\vj)} : \Gamma_1 + \yoneda\vC \to \Gamma_1$.
    
    The assumption $\inj_\vi(\vf_\vi) \neq \inj_{\vj}(\vf_\vj)$, gives us exactly that they two copairing morphisms are not equal.
    Notice also that $\alpha \comp \copair1{\inj_\vi(\vf_\vi)} = \alpha \comp \copair1{\inj_\vj(\vf_\vj)}$ by construction.

    But $\alpha$ being mono gives us that $\copair1{\inj_\vi(\vf_\vi)} = \copair1{\inj_\vj(\vf_\vj)}$ which is a contradiction.
\end{proof}

\begin{corollary}\label{lem:fsreprC-base-mono}
    If $\alpha : \Sigma_1 \to \Sigma_2 \in \fsreprC\mcatC$ is mono then every $\alpha(\inj_\vi(1_{\vC_\vi}))$ for $\vC_\vi \in \Sigma$ is mono.
\end{corollary}
\begin{proof}
    Immediate from \cref{lem:fsreprC-mono}.
\end{proof}



\printbibliography

\end{document}
\endinput
